% Limites et continuité — Mathématiques 1re Tech — Cours signé OnBuch+
% Programme officiel MINESEC. Compiler avec : tectonic N21-limites-continuite.tex
\newcommand{\DOCMATIERE}{Mathématiques}
\newcommand{\DOCNIVEAU}{1re Tech}
\newcommand{\DOCMODULE}{Mathématiques – classe de Première technique}
\newcommand{\DOCLECON}{N21}
\newcommand{\DOCTITRE}{Limites et continuité}
% OnBuch+ — préambule « cours signés » ENRICHI (compilé avec tectonic / XeLaTeX)
% Système visuel « Pop » d'OnBuch+ : fond crème (jamais blanc), bordures encre
% épaisses, ombres « dures » décalées, titres Archivo Black en capitales.
% Version MATHÉMATIQUES : graphes (pgfplots/tikz), boîtes pédagogiques
% (définition, propriété, méthode, attention, le savais-tu, exercice résolu,
% exercice, TP, bilan), page de garde et signature OnBuch+ ancrée en bas.
%
% Macros attendues dans main.tex AVANT \input{preamble} :
%   \DOCMATIERE  \DOCNIVEAU  \DOCMODULE  \DOCLECON (numéro)  \DOCTITRE
\documentclass[11pt]{article}

\usepackage{fontspec}
\usepackage[french]{babel}
\usepackage[a4paper,margin=2cm,top=3.4cm,bottom=2.8cm,headheight=1.4cm,footskip=1.3cm]{geometry}
\usepackage{amsmath,amssymb,amsfonts}
\usepackage{mathtools} % \xmapsto, \coloneqq... souvent utilisés par les modèles
\usepackage{cancel} % simplifications barrées
% \abs{z} (module, valeur absolue) et \norm{u} (norme), utilisés par les modèles
\providecommand{\abs}{}\let\abs\relax\DeclarePairedDelimiter{\abs}{\lvert}{\rvert}
\providecommand{\norm}{}\let\norm\relax\DeclarePairedDelimiter{\norm}{\lVert}{\rVert}
\usepackage[table]{xcolor} % [table] : fournit aussi \rowcolors (alternance de lignes)
\usepackage{pagecolor}
\usepackage{tikz}
\usetikzlibrary{arrows.meta,positioning,calc,shapes.geometric,shapes.misc,shapes.arrows,decorations.pathmorphing,decorations.pathreplacing,decorations.markings,patterns,shadows,mindmap,trees,fit,backgrounds,matrix,intersections,angles,quotes}
\usepackage{pgfplots}
\usepackage{pgf-pie} % diagrammes circulaires \pie (statistiques)
\pgfplotsset{compat=1.18}
\usepgfplotslibrary{fillbetween,groupplots} % aires entre courbes (calcul intégral)
\usepackage{enumitem}
\usepackage{fancyhdr}
% [most] charge toutes les bibliothèques tcolorbox (listings, minted, raster,
% documentation...) et gonfle inutilement la mémoire TeX sur un document long
% (~300 boîtes) : on ne charge que ce qui est réellement utilisé.
\usepackage[skins,breakable]{tcolorbox}
\usepackage{graphicx}
\usepackage{colortbl}
\usepackage{booktabs}
\usepackage{tabularx}
\usepackage{diagbox} % case d'en-tête diagonale des tableaux à double entrée
% Colonnes extensibles centrée / gauche / droite (C, L, R), souvent utilisées par les modèles
\newcolumntype{C}{>{\centering\arraybackslash}X}
\newcolumntype{L}{>{\raggedright\arraybackslash}X}
\newcolumntype{R}{>{\raggedleft\arraybackslash}X}
\usepackage{array}
\usepackage{multicol}
\usepackage{siunitx}
% parse-numbers=false : accepte aussi \SI{2,5 \times 10^{-3}}{...}
\sisetup{locale=FR,output-decimal-marker={,},group-minimum-digits=4,parse-numbers=false}
\DeclareSIUnit\ohm{\ensuremath{\symup{\Omega}}} % U+2126 absent de la police
\usepackage{qrcode}
% \degree : commande fréquemment utilisée par les modèles pour un angle ou une
% température en dehors de \SI{}{} (non fournie par les paquets chargés).
\providecommand{\degree}{^{\circ}}
% \qed (fin de démonstration), utilisé par les modèles sans amsthm
\providecommand{\qed}{\hfill$\square$}
% \barre : double barre des valeurs interdites dans un tableau de variations (tkz-tab)
\providecommand{\barre}{\ensuremath{\|}}
% environnement proof (amsthm non chargé) utilisé par les modèles
\ifdefined\proof\else\newenvironment{proof}[1][Démonstration]{\par\noindent\textit{#1.}\ }{\qed\par}\fi
\usepackage{lastpage}
\usepackage{hyperref}
\hypersetup{hidelinks}

% ============================================================
% Palette « Pop » OnBuch+ — mêmes hex que la classe P (plus_ui.dart)
% ============================================================
\definecolor{poporange}{HTML}{F59321}
\definecolor{popdark}{HTML}{E07A0C}
\definecolor{popcream}{HTML}{FFF6E8}
\definecolor{popink}{HTML}{1C1714}
\definecolor{poppurple}{HTML}{7A5AE0}
\definecolor{popgreen}{HTML}{2E9E6A}
\definecolor{poppink}{HTML}{E0457B}
\definecolor{popblue}{HTML}{2F7DE1}
\definecolor{popgold}{HTML}{C98A12}
\definecolor{popmuted}{HTML}{8A7F76}
\definecolor{popline}{HTML}{EADFD2}
\definecolor{popgray}{HTML}{8A7F76}
\colorlet{popgrey}{popgray}
% Alias tolérants (le modèle nomme parfois les couleurs différemment)
\colorlet{popwhite}{white}
\colorlet{popblack}{popink}
\colorlet{popred}{poppink}
\colorlet{popyellow}{popgold}
% Teintes claires pour fonds de tableaux / zones
\colorlet{popblueL}{popblue!10!white}
\colorlet{poporangeL}{poporange!14!white}
\colorlet{popgreenL}{popgreen!12!white}
\colorlet{poppurpleL}{poppurple!12!white}
\colorlet{poppinkL}{poppink!12!white}
\colorlet{popgoldL}{popgold!14!white}

\pagecolor{popcream}

% ============================================================
% Typographie
% ============================================================
\defaultfontfeatures{Ligatures=TeX}
\setmainfont{PlusJakartaSans-Medium.ttf}[Path=../../../fonts/,BoldFont=PlusJakartaSans-Bold.ttf]
% Maths en sans-serif (Fira Math) avec chiffres et lettres droites de Plus
% Jakarta, pour que les formules se fondent dans le texte.
\usepackage[math-style=ISO,bold-style=ISO]{unicode-math}
\setmathfont{FiraMath-Regular.otf}[Path=../../../fonts/]
\setmathfont{PlusJakartaSans-Medium.ttf}[Path=../../../fonts/,range={up/num,up/{Latin,latin},\mathup}]
\usepackage{icomma} % virgule décimale sans espace en mode math
% Caractères mathématiques Unicode absents des polices de texte (Plus Jakarta,
% Archivo Black) : rendus par leur équivalent mathématique, en texte comme en
% formule — sinon ils s'affichent en carré vide (ex. « Arithmétique dans ℤ »).
\usepackage{newunicodechar}
\newunicodechar{ℝ}{\ensuremath{\mathbb{R}}}
\newunicodechar{ℤ}{\ensuremath{\mathbb{Z}}}
\newunicodechar{ℂ}{\ensuremath{\mathbb{C}}}
\newunicodechar{ℕ}{\ensuremath{\mathbb{N}}}
\newunicodechar{ℚ}{\ensuremath{\mathbb{Q}}}
\newunicodechar{𝔻}{\ensuremath{\mathbb{D}}}
\newunicodechar{α}{\ensuremath{\alpha}}
\newunicodechar{β}{\ensuremath{\beta}}
\newunicodechar{γ}{\ensuremath{\gamma}}
\newunicodechar{δ}{\ensuremath{\delta}}
\newunicodechar{ε}{\ensuremath{\varepsilon}}
\newunicodechar{θ}{\ensuremath{\theta}}
\newunicodechar{λ}{\ensuremath{\lambda}}
\newunicodechar{μ}{\ensuremath{\mu}}
\newunicodechar{σ}{\ensuremath{\sigma}}
\newunicodechar{τ}{\ensuremath{\tau}}
\newunicodechar{φ}{\ensuremath{\varphi}}
\newunicodechar{ψ}{\ensuremath{\psi}}
\newunicodechar{ω}{\ensuremath{\omega}}
\newunicodechar{Σ}{\ensuremath{\Sigma}}
% majuscules produites par \MakeUppercase dans les titres (α-aminés -> Α)
\newunicodechar{Α}{\ensuremath{\alpha}}
\newunicodechar{Β}{\ensuremath{\beta}}
\newunicodechar{Γ}{\ensuremath{\Gamma}}
\newunicodechar{Δ}{\ensuremath{\Delta}}
\newunicodechar{Ε}{\ensuremath{\varepsilon}}
\newunicodechar{Θ}{\ensuremath{\Theta}}
\newunicodechar{Λ}{\ensuremath{\Lambda}}
\newunicodechar{Μ}{\ensuremath{\mu}}
\newunicodechar{Τ}{\ensuremath{\tau}}
\newunicodechar{Φ}{\ensuremath{\Phi}}
\newunicodechar{Ψ}{\ensuremath{\Psi}}
\newunicodechar{Ω}{\ensuremath{\Omega}}
\newunicodechar{↦}{\ensuremath{\mapsto}}
\newunicodechar{∈}{\ensuremath{\in}}
\newunicodechar{∉}{\ensuremath{\notin}}
\newunicodechar{∧}{\ensuremath{\wedge}}
\newunicodechar{⇔}{\ensuremath{\Leftrightarrow}}
\newunicodechar{⟺}{\ensuremath{\Longleftrightarrow}}
\newunicodechar{⇒}{\ensuremath{\Rightarrow}}
\newunicodechar{⟹}{\ensuremath{\Longrightarrow}}
\newunicodechar{∘}{\ensuremath{\circ}}
\newunicodechar{∪}{\ensuremath{\cup}}
\newunicodechar{∩}{\ensuremath{\cap}}
\newunicodechar{≡}{\ensuremath{\equiv}}
\newunicodechar{⊂}{\ensuremath{\subset}}
\newunicodechar{⊥}{\ensuremath{\perp}}
\newunicodechar{∃}{\ensuremath{\exists}}
\newunicodechar{∀}{\ensuremath{\forall}}
\newunicodechar{∛}{\ensuremath{\sqrt[3]{\,}}}
\newunicodechar{∜}{\ensuremath{\sqrt[4]{\,}}}
\newunicodechar{⃗}{\ensuremath{{}^{\to}}}
\newunicodechar{ⁿ}{\ifmmode^{n}\else\textsuperscript{\ensuremath{n}}\fi}
\newunicodechar{ˣ}{\ifmmode^{x}\else\textsuperscript{\ensuremath{x}}\fi}
\newunicodechar{ᵇ}{\ifmmode^{b}\else\textsuperscript{\ensuremath{b}}\fi}
\newunicodechar{ʳ}{\ifmmode^{r}\else\textsuperscript{\ensuremath{r}}\fi}
\newunicodechar{ᵘ}{\ifmmode^{u}\else\textsuperscript{\ensuremath{u}}\fi}
\newunicodechar{ʸ}{\ifmmode^{y}\else\textsuperscript{\ensuremath{y}}\fi}
\newunicodechar{ᵃ}{\ifmmode^{a}\else\textsuperscript{\ensuremath{a}}\fi}
\newunicodechar{ᵉ}{\ifmmode^{e}\else\textsuperscript{\ensuremath{e}}\fi}
\newunicodechar{ᵏ}{\ifmmode^{k}\else\textsuperscript{\ensuremath{k}}\fi}
\newunicodechar{ᵖ}{\ifmmode^{p}\else\textsuperscript{\ensuremath{p}}\fi}
\newunicodechar{ᴺ}{\ifmmode^{N}\else\textsuperscript{\ensuremath{N}}\fi}
\newunicodechar{ᶜ}{\ifmmode^{c}\else\textsuperscript{\ensuremath{c}}\fi}
\newunicodechar{ʰ}{\ifmmode^{h}\else\textsuperscript{\ensuremath{h}}\fi}
\newunicodechar{ᵠ}{\ifmmode^{q}\else\textsuperscript{\ensuremath{q}}\fi}
\newunicodechar{ₙ}{\ifmmode_{n}\else\textsubscript{\ensuremath{n}}\fi}
\newunicodechar{ₐ}{\ifmmode_{a}\else\textsubscript{\ensuremath{a}}\fi}
\newunicodechar{ᵢ}{\ifmmode_{i}\else\textsubscript{\ensuremath{i}}\fi}
\newunicodechar{ᵣ}{\ifmmode_{r}\else\textsubscript{\ensuremath{r}}\fi}
\newunicodechar{ₖ}{\ifmmode_{k}\else\textsubscript{\ensuremath{k}}\fi}
\newunicodechar{ᵦ}{\ifmmode_{\beta}\else\textsubscript{\ensuremath{\beta}}\fi}
\newunicodechar{ₓ}{\ifmmode_{x}\else\textsubscript{\ensuremath{x}}\fi}
\newfontfamily\popdisplay{ArchivoBlack-Regular.ttf}[Path=../../../fonts/]
\newfontfamily\poplogo{SpaceGrotesk-Bold.ttf}[Path=../../../fonts/]
\newfontfamily\popsemi{PlusJakartaSans-SemiBold.ttf}[Path=../../../fonts/]
\newfontfamily\popxbold{PlusJakartaSans-ExtraBold.ttf}[Path=../../../fonts/]
\newfontfamily\popxboldls{PlusJakartaSans-ExtraBold.ttf}[Path=../../../fonts/,LetterSpace=3.0]

\setlist[enumerate]{leftmargin=1.7em,itemsep=5pt,topsep=4pt}
\setlist[itemize]{leftmargin=1.7em,itemsep=4pt,topsep=4pt,label={\color{poporange}\textbullet}}
\setlength{\parindent}{0pt}
\setlength{\parskip}{5pt}
\renewcommand{\arraystretch}{1.25}

% pgfplots : style Pop par défaut
\pgfplotsset{
  every axis/.append style={axis line style={popink,line width=1pt},tick style={popink},
    grid style={popline,line width=0.5pt},label style={font=\small},tick label style={font=\footnotesize}},
  popcurve/.style={poporange,line width=1.8pt,no markers,smooth},
}
% Même style disponible dans un \draw TikZ simple (hors pgfplots)
\tikzset{popcurve/.style={poporange,line width=1.8pt}}

% ============================================================
% Pill / squiggle
% ============================================================
\newcommand{\poppill}[3]{%
  \tikz[baseline=-0.55ex]\node[draw=popink,line width=1.2pt,rounded corners=2.6pt,%
    fill=#2,inner xsep=6.5pt,inner ysep=3pt]{%
    \popxboldls\fontsize{7}{7}\selectfont\color{#3}\MakeUppercase{#1}};%
}
\newcommand{\popsquiggle}[1]{%
  \tikz[baseline=-0.4ex]{
    \draw[#1,line width=2.6pt,line cap=round]
      plot [smooth] coordinates {(0,0.09) (0.34,0.01) (0.68,0.21) (1.02,0.01) (1.36,0.10)};
  }%
}
% Mot-clé surligné (marqueur orange sous le texte)
\newcommand{\cle}[1]{{\popxbold\color{popdark}#1}}

% ============================================================
% En-tête / pied
% ============================================================
\pagestyle{fancy}
\fancyhf{}
\renewcommand{\headrulewidth}{0pt}
\renewcommand{\footrulewidth}{0pt}
\lhead{\raisebox{-0.15ex}{\poplogo\fontsize{13}{13}\selectfont\color{popink}OnBuch\color{poporange}+}}
\rhead{\poppill{\DOCMATIERE\ \textbullet\ \DOCNIVEAU\ \textbullet\ Leçon \DOCLECON}{white}{popink}}
\lfoot{\popxbold\fontsize{7.6}{7.6}\selectfont\color{popmuted}\MakeUppercase{Cours signé OnBuch+}\ \textbullet\ {\popsemi\DOCTITRE}}
\rfoot{\tikz[baseline=-0.6ex]\node[rounded corners=8.5pt,fill=popink,minimum height=17pt,minimum width=17pt,inner xsep=5pt,inner sep=0]{\popxbold\fontsize{8}{8}\selectfont\color{popcream}\thepage\,/\,\pageref*{LastPage}};}

% ============================================================
% Titres
% ============================================================
\newcommand{\chaptitre}[2]{%
  \begin{tcolorbox}[enhanced,colback=poporange,colframe=popink,boxrule=2.6pt,arc=11pt,
      drop shadow=popink,left=15pt,right=15pt,top=12pt,bottom=11pt,before skip=3pt,after skip=18pt]
    \poppill{#2}{popink}{popcream}\par\vspace{9pt}
    {\raggedright\hyphenpenalty=10000\exhyphenpenalty=10000\popdisplay\fontsize{22}{24}\selectfont\color{popcream}\MakeUppercase{#1}\par}
  \end{tcolorbox}
}
% Section numérotée : pastille orange avec le numéro + titre Archivo
\newcounter{popsec}
\newcommand{\coursec}[1]{%
  \stepcounter{popsec}\setcounter{popsub}{0}%
  \par\vspace{16pt}\noindent
  \begin{minipage}{\linewidth}
  \tikz[baseline=-0.7ex]\node[draw=popink,line width=1.6pt,fill=poporange,rounded corners=4pt,minimum size=22pt,inner sep=0]{\popdisplay\fontsize{12}{12}\selectfont\color{popcream}\thepopsec};%
  \hspace{8pt}{\popdisplay\fontsize{15}{17}\selectfont\color{popink}\MakeUppercase{#1}}\par
  \vspace{4pt}\noindent{\color{poporange}\rule{36pt}{3.6pt}}
  \end{minipage}\par\nopagebreak\vspace{8pt}
}
\newcounter{popsub}
\newcommand{\courssub}[1]{%
  \stepcounter{popsub}%
  \par\vspace{9pt}\noindent{\popxbold\fontsize{11.5}{13}\selectfont\color{popink}{\color{poporange}\thepopsec.\thepopsub}\hspace{6pt}#1}\par\nopagebreak\vspace{4pt}
}

% ============================================================
% Boîtes pédagogiques — même recette : fond blanc, bordure épaisse colorée,
% ombre dure encre, badge plein. Titre optionnel [#1].
% ============================================================
\tcbset{popbase/.style={breakable,enhanced,colback=white,boxrule=2.3pt,arc=9pt,
  shadow={3pt}{-3pt}{0pt}{popink},left=14pt,right=14pt,top=11pt,bottom=11pt,before skip=11pt,after skip=15pt,
  fontupper=\popsemi\fontsize{10.6}{14.5}\selectfont\color{popink},
  fontlower=\popsemi\fontsize{10.6}{14.5}\selectfont\color{popink}}}
% Coche dessinée (le glyphe ✓ n'existe pas dans les polices du document)
\newcommand{\popcheck}[1]{\tikz[baseline=-0.5ex]\draw[#1,line width=1.6pt,line cap=round,line join=round](-0.13,0.0)--(-0.04,-0.09)--(0.13,0.1);}
\newcommand{\popboxhead}[3]{\poppill{#1}{#2}{white}\ifx\relax#3\relax\else\hspace{6pt}{\popxbold\fontsize{10.6}{12}\selectfont\color{popink}#3}\fi\par\vspace{7pt}}

\newtcolorbox{aretenir}[1][]{popbase,colframe=popgreen,before upper={\popboxhead{À retenir}{popgreen}{#1}}}
\newtcolorbox{exemplebox}[1][]{popbase,colframe=poporange,before upper={\popboxhead{Exemple}{poporange}{#1}}}
\newtcolorbox{definition}[1][]{popbase,colframe=popblue,before upper={\popboxhead{Définition}{popblue}{#1}}}
\newtcolorbox{propriete}[1][]{popbase,colframe=poppurple,before upper={\popboxhead{Propriété}{poppurple}{#1}}}
\newtcolorbox{methode}[1][]{popbase,colframe=popgold,before upper={\popboxhead{Méthode}{popgold}{#1}}}
\newtcolorbox{attention}[1][]{popbase,colframe=poppink,before upper={\popboxhead{Attention}{poppink}{#1}}}
\newtcolorbox{experience}[1][]{popbase,colframe=popblue,colback=popblueL,before upper={\popboxhead{Expérience}{popblue}{#1}}}
% « Le savais-tu ? » : carte violette pleine (contexte, culture, Cameroun)
\newtcolorbox{savaistu}[1][]{popbase,colback=poppurple,colframe=popink,
  fontupper=\popsemi\fontsize{10.4}{14.2}\selectfont\color{white},
  before upper={\poppill{Le savais-tu\,?}{white}{poppurple}\ifx\relax#1\relax\else\hspace{6pt}{\popxbold\color{white}#1}\fi\par\vspace{7pt}}}
% Exercice résolu : énoncé en haut, \tcblower puis la solution sur fond crème
\newcounter{popexo}\newcounter{popexr}
\newtcolorbox{exoresolu}[1][]{popbase,colframe=poporange,colbacklower=popcream,
  segmentation style={popink,line width=1.2pt,dashed},
  before upper={\stepcounter{popexr}\popboxhead{Exercice résolu \thepopexr}{poporange}{#1}},
  before lower={\poppill{Solution}{popink}{popcream}\par\vspace{6pt}}}
% Exercice (non résolu en place) — la correction est dans la partie Corrigés
\newtcolorbox{exercice}[1][]{popbase,colframe=popink,
  before upper={\stepcounter{popexo}\popboxhead{Exercice \thepopexo}{popink}{#1}}}
% Corrigé
\newtcolorbox{corrige}[1][]{popbase,colframe=popgreen,colback=popgreenL,
  before upper={\popboxhead{Corrigé}{popgreen}{#1}}}
% Objectifs / prérequis / bilan
\newtcolorbox{objectifs}{popbase,colframe=popink,colback=poporangeL,
  before upper={\popboxhead{Objectifs de la leçon}{popink}{}}}
\newtcolorbox{prerequis}{popbase,colframe=popink,colback=popblueL,
  before upper={\popboxhead{Prérequis}{popblue}{}}}
\newtcolorbox{bilan}{popbase,colframe=popink,colback=white,boxrule=2.8pt,
  before upper={\popboxhead{Fiche bilan}{poporange}{L'essentiel à connaître pour l'examen}}}
% Figure encadrée avec légende
\newtcolorbox{popfigure}[1][]{enhanced,breakable=false,colback=white,colframe=popline,boxrule=1.4pt,arc=8pt,
  left=10pt,right=10pt,top=10pt,bottom=8pt,before skip=10pt,after skip=12pt,halign=center,
  lower separated=false,halign lower=center,
  fontlower=\popsemi\fontsize{9}{11.5}\selectfont\color{popmuted},
  #1}
\newcommand{\legende}[1]{\par\vspace{6pt}{\popsemi\fontsize{9}{11.5}\selectfont\color{popmuted}#1}}

% ============================================================
% Signature OnBuch+
% ============================================================
\newcommand{\poplogoinline}[1]{{\poplogo\fontsize{#1}{#1}\selectfont\color{popink}OnBuch\color{poporange}+}}
\newcommand{\signatureonbuch}{%
  \begin{tcolorbox}[enhanced,colback=popink,colframe=popink,boxrule=2.2pt,arc=10pt,
      drop shadow=poporange,left=14pt,right=14pt,top=10pt,bottom=10pt,before skip=0pt,after skip=0pt]
    \tikz[baseline=-0.7ex]\node[circle,fill=popgreen,draw=popcream,line width=1.3pt,minimum size=22pt,inner sep=0]{\popcheck{white}};%
    \hspace{9pt}\begin{minipage}[c]{0.62\linewidth}
      {\popxbold\fontsize{12}{13}\selectfont\color{popcream}Signé\ \textbullet\ {\poplogo OnBuch\color{poporange}+}}\par\vspace{2pt}
      {\raggedright\popsemi\fontsize{8}{10.5}\selectfont\color{popline}\MakeUppercase{Équipe pédagogique OnBuch+}\\\MakeUppercase{Conforme au programme officiel MINESEC}\par}
    \end{minipage}\hfill
    {\poplogo\fontsize{22}{22}\selectfont\color{popcream}OnBuch\color{poporange}+}
  \end{tcolorbox}%
}

% ============================================================
% Page de garde
% #1 titre · #2 matière · #3 niveau · #4 module · #5 n° leçon · #6 durée ·
% #7 description / objectifs (texte ou itemize)
% ============================================================
\newcommand{\pagegarde}[7]{%
  \thispagestyle{empty}
  \newgeometry{a4paper,margin=1.7cm,top=1.6cm,bottom=1.4cm}%
  \begin{tikzpicture}[remember picture,overlay]
    \fill[poporange!18] ([xshift=-3.2cm,yshift=-3.6cm]current page.north east) circle (4.6cm);
    \fill[poppurple!14] ([xshift=2.4cm,yshift=7.6cm]current page.south west) circle (3.8cm);
  \end{tikzpicture}%
  {\poplogo\fontsize{30}{30}\selectfont\color{popink}OnBuch\color{poporange}+}\hfill
  \raisebox{0.6ex}{\poppill{Cours signé \textbullet\ Premium}{popink}{popcream}}\par
  \vspace{1pt}\popsquiggle{poppurple}\par
  \vspace{8pt}
  \begin{tcolorbox}[enhanced,colback=poporange,colframe=popink,boxrule=2.8pt,arc=13pt,
      drop shadow=popink,left=18pt,right=18pt,top=12pt,bottom=13pt,after skip=10pt]
    \poppill{#2 \textbullet\ #3}{popink}{popcream}\hspace{5pt}\poppill{#4}{white}{popink}\par\vspace{8pt}
    {\popdisplay\fontsize{14}{14}\selectfont\color{popink}LEÇON #5}\par\vspace{4pt}
    {\raggedright\hyphenpenalty=10000\exhyphenpenalty=10000\popdisplay\fontsize{25}{27}\selectfont\color{popcream}\MakeUppercase{#1}\par}
    \vspace{8pt}
    {\popxbold\fontsize{9.5}{9.5}\selectfont\color{popink}\MakeUppercase{Durée conseillée : #6}}
  \end{tcolorbox}
  \begin{tcolorbox}[enhanced,colback=white,colframe=popink,boxrule=2.2pt,arc=10pt,
      drop shadow=popink,left=16pt,right=16pt,top=10pt,bottom=10pt,after skip=8pt]
    \poppill{Dans cette leçon}{poppurple}{white}\par\vspace{5pt}
    \popsemi\fontsize{9.3}{12.6}\selectfont\color{popink}#7
  \end{tcolorbox}
  \noindent\begin{minipage}[t]{0.487\linewidth}
    \begin{tcolorbox}[enhanced,colback=popgreen,colframe=popink,boxrule=2.2pt,arc=10pt,
        drop shadow=popink,left=11pt,right=11pt,top=10pt,bottom=10pt,height=3.1cm,valign=center]
      \tcbox[colback=white,colframe=popink,boxrule=1.3pt,arc=5pt,boxsep=0pt,nobeforeafter,
        left=3pt,right=3pt,top=3pt,bottom=3pt,valign=center]{\qrcode[height=1.9cm]{https://play.google.com/store/apps/details?id=com.luvvix.onbuch&hl=fr}}%
      \hspace{8pt}\begin{minipage}[c]{0.5\linewidth}
        {\popdisplay\fontsize{11}{12.5}\selectfont\color{white}\MakeUppercase{Télécharge OnBuch}}\par\vspace{4pt}
        {\raggedright\popsemi\fontsize{8.4}{11}\selectfont\color{white}Gratuit \textbullet\ Google Play\\Tuteur IA, annales, concours, résultats\par}
      \end{minipage}
    \end{tcolorbox}
  \end{minipage}\hfill
  \begin{minipage}[t]{0.487\linewidth}
    \begin{tcolorbox}[enhanced,colback=poppurple,colframe=popink,boxrule=2.2pt,arc=10pt,
        drop shadow=popink,left=11pt,right=11pt,top=10pt,bottom=10pt,height=3.1cm,valign=center]
      \tikz[baseline=-0.7ex]\node[circle,fill=white,draw=popink,line width=1.3pt,minimum size=1.6cm,inner sep=0]{\popdisplay\fontsize{24}{24}\selectfont\color{poppurple}+};%
      \hspace{8pt}\begin{minipage}[c]{0.6\linewidth}
        {\popdisplay\fontsize{11}{12.5}\selectfont\color{white}\MakeUppercase{Passe à OnBuch+}}\par\vspace{4pt}
        {\raggedright\popsemi\fontsize{8.4}{11}\selectfont\color{white}Tous les cours signés, Tuteur IA illimité, fiches PDF — onglet OnBuch+ de l'app\par}
      \end{minipage}
    \end{tcolorbox}
  \end{minipage}
  \vspace{8pt}
  \popcontenu
  \vfill
  \signatureonbuch
  \restoregeometry
  \newpage
}

% Rangée « Ce cours contient » (page de garde)
\newcommand{\poptile}[3]{% #1 couleur #2 gros texte #3 légende
  \begin{tcolorbox}[enhanced,colback=#1,colframe=popink,boxrule=2pt,arc=9pt,shadow={2.5pt}{-2.5pt}{0pt}{popink},
      left=6pt,right=6pt,top=9pt,bottom=9pt,halign=center,valign=center,height=1.75cm,width=\linewidth,nobeforeafter]
    {\popdisplay\fontsize{12.5}{13}\selectfont\color{white}\MakeUppercase{#2}}\par\vspace{3pt}
    {\centering\popsemi\fontsize{7.8}{9.5}\selectfont\color{white}#3\par}
  \end{tcolorbox}}
\newcommand{\popcontenu}{%
  \par\noindent{\popxboldls\fontsize{7.5}{7.5}\selectfont\color{popmuted}CE COURS SIGNÉ CONTIENT}\par\vspace{7pt}
  \noindent\begin{minipage}[t]{0.235\linewidth}\poptile{popblue}{Cours}{complet, illustré, pas à pas}\end{minipage}\hfill
  \begin{minipage}[t]{0.235\linewidth}\poptile{poporange}{Méthodes}{et exercices résolus}\end{minipage}\hfill
  \begin{minipage}[t]{0.235\linewidth}\poptile{poppink}{Exercices}{type examen + corrigés}\end{minipage}\hfill
  \begin{minipage}[t]{0.235\linewidth}\poptile{popgreen}{Fiche bilan}{l'essentiel à retenir}\end{minipage}\par}

% Sommaire
\newtcolorbox{sommaire}{enhanced,colback=white,colframe=popink,boxrule=2.2pt,arc=10pt,
  drop shadow=popink,left=18pt,right=18pt,top=13pt,bottom=13pt,before skip=6pt,after skip=20pt,
  before upper={\poppill{Sommaire}{poppurple}{white}\par\vspace{9pt}\popsemi\fontsize{10.6}{14.5}\selectfont\color{popink}}}

% Signature de fin de cours — poussée tout en bas de la dernière page
\newcommand{\findecours}{%
  \par\vfill\nopagebreak
  \begin{center}\popsquiggle{poporange}\end{center}\vspace{4pt}
  \signatureonbuch
}
\AtBeginDocument{\def\llbracket{\mathopen{[\mkern-3.5mu[}}\def\rrbracket{\mathclose{]\mkern-3.5mu]}}} % intervalles d'entiers (glyphe ⟦ absent de la police)
% Glyphes absents de Fira Math : redéfinis à partir de glyphes disponibles
\AtBeginDocument{%
  \def\setminus{\mathbin{\backslash}}\let\smallsetminus\setminus
  \def\vdots{\mathord{\vbox{\baselineskip3.2pt\lineskiplimit0pt\kern4pt\hbox{.}\hbox{.}\hbox{.}}}}%
  \def\ddots{\mathinner{\mkern1mu\raise6.5pt\hbox{.}\mkern2mu\raise3.6pt\hbox{.}\mkern2mu\raise0.7pt\hbox{.}\mkern1mu}}%
  \def\triangle{\mathord{\scalebox{0.9}{$\symup{\Delta}$}}}%
  \def\nabla{\mathord{\raisebox{\depth}{\scalebox{1}[-1]{$\symup{\Delta}$}}}}%
  \def\checkmark{\popcheck{popdark}}%
  \def\star{\mathbin{\ast}}\def\bigstar{\mathord{\ast}}%
  \def\diamond{\mathbin{\scalebox{0.6}{$\Diamond$}}}%
  \def\bigcup{\mathop{\mathchoice{\raisebox{-0.3em}{\scalebox{1.7}{$\cup$}}}{\scalebox{1.25}{$\cup$}}{\cup}{\cup}}\displaylimits}%
  \def\bigcap{\mathop{\mathchoice{\raisebox{-0.3em}{\scalebox{1.7}{$\cap$}}}{\scalebox{1.25}{$\cap$}}{\cap}{\cap}}\displaylimits}%
  \def\lesseqqgtr{\mathrel{\lessgtr}}\def\bigoplus{\mathop{\oplus}\displaylimits}%
}
\providecommand{\ssi}{si et seulement si} % abréviation fréquente
\AtBeginDocument{\providecommand{\wideparen}{\overparen}} % arc de cercle

%% FIGFIX-BEGIN v3 — correctifs globaux des figures (docs/onbuch-generation/tools/figfix)
\usepackage{adjustbox}\usepackage{etoolbox}
% toute figure TikZ/pgfplots plus large que la ligne est réduite à la largeur disponible
\BeforeBeginEnvironment{tikzpicture}{\begin{adjustbox}{max width=\linewidth}}
\AfterEndEnvironment{tikzpicture}{\end{adjustbox}}
% légendes pgfplots lisibles par-dessus les courbes
\pgfplotsset{every axis/.append style={legend style={fill=white,fill opacity=0.92,text opacity=1,draw=black!15,inner sep=3pt}}}
% étiquettes d'arêtes (syntaxe edge["..."]) avec fond blanc
\tikzset{every edge quotes/.append style={fill=white,inner sep=1.2pt}}
% bulles de cartes mentales : texte à la ligne dans la bulle, bulle un peu plus grande
\tikzset{every concept/.append style={text width=2.0cm,align=center,minimum size=2.3cm,inner sep=2pt}}
%% FIGFIX-END
\begin{document}

\pagegarde{Limites et continuité}{Mathématiques}{1re Tech}{Mathématiques – classe de Première technique}{N21}{3 séances (fiche de progression harmonisée)}{Cette leçon introduit de manière intuitive puis opérationnelle les notions de limite d'une fonction (en un point et à l'infini) et de continuité. Elle fournit les outils de calcul pratique des limites (opérations, formes indéterminées simples) indispensables pour l'étude des fonctions en Première et en Terminale technique.
\vspace{4pt}
\begin{multicols}{2}\small
\begin{itemize}
\item À la fin de cette leçon, je sais décrire intuitivement, à l'aide d'un graphique ou d'un tableau de valeurs, la limite d'une fonction en un point ou à l'infini
\item À la fin de cette leçon, je sais lire graphiquement une limite finie ou infinie et reconnaître une asymptote verticale ou horizontale
\item À la fin de cette leçon, je sais citer les limites de référence des fonctions usuelles (puissances, inverse, racine carrée)
\item À la fin de cette leçon, je sais calculer une limite en utilisant les opérations (somme, produit, quotient) sur les limites
\item À la fin de cette leçon, je sais lever une forme indéterminée simple (polynôme, fonction rationnelle) en factorisant par le terme de plus haut degré
\item À la fin de cette leçon, je sais déterminer la limite en un point d'une fonction rationnelle par factorisation et simplification
\end{itemize}\end{multicols}}

\begin{sommaire}
\begin{multicols}{2}
\begin{enumerate}
\item Approche intuitive de la notion de limite
\item Limites de référence et interprétation graphique
\item Opérations sur les limites
\item Calculs pratiques : lever les formes indéterminées
\item Continuité d'une fonction
\item Applications à des situations concrètes
\item Activité d'intégration
\item Méthodes et exercices résolus
\item Exercices
\item Corrigés des exercices
\item Fiche bilan
\end{enumerate}
\end{multicols}
\end{sommaire}

\begin{objectifs}
\begin{itemize}
\item À la fin de cette leçon, je sais décrire intuitivement, à l'aide d'un graphique ou d'un tableau de valeurs, la limite d'une fonction en un point ou à l'infini
\item À la fin de cette leçon, je sais lire graphiquement une limite finie ou infinie et reconnaître une asymptote verticale ou horizontale
\item À la fin de cette leçon, je sais citer les limites de référence des fonctions usuelles (puissances, inverse, racine carrée)
\item À la fin de cette leçon, je sais calculer une limite en utilisant les opérations (somme, produit, quotient) sur les limites
\item À la fin de cette leçon, je sais lever une forme indéterminée simple (polynôme, fonction rationnelle) en factorisant par le terme de plus haut degré
\item À la fin de cette leçon, je sais déterminer la limite en un point d'une fonction rationnelle par factorisation et simplification
\item À la fin de cette leçon, je sais reconnaître graphiquement et justifier la continuité d'une fonction en un point ou sur un intervalle
\item À la fin de cette leçon, je sais appliquer ces notions à des situations concrètes de la vie courante (coûts, vitesses, remplissage) et rédiger une démarche justifiée
\end{itemize}
\end{objectifs}

\begin{prerequis}
\begin{itemize}
\item Fonctions usuelles : fonctions affines, carré, inverse, racine carrée, polynômes (Seconde)
\item Lecture et interprétation de tableaux de valeurs et de représentations graphiques
\item Factorisation d'expressions algébriques et identités remarquables (4ème/3ème, révisées en Seconde)
\item Calcul algébrique sur les fractions et les puissances
\item Notion d'intervalle et d'ensemble de définition d'une fonction
\end{itemize}
\end{prerequis}

\newpage

\coursec{Approche intuitive de la notion de limite}

\courssub{Situation d'introduction : le bénéfice de Maman Ngo Bell}

Maman Ngo Bell vend des beignets au marché Mokolo. Chaque matin, elle achète de l'huile pour sa friture. Elle a remarqué un phénomène : lorsque le prix du litre d'huile augmente et se rapproche de \SI{1500}{FCFA}, son bénéfice journalier diminue et se rapproche de zéro. Si l'on note $x$ le prix du litre d'huile et $B(x)$ son bénéfice, on peut modéliser la situation par une formule du type
\[ B(x) = 3000 - 2x \quad \text{(en FCFA).} \]
Quand $x$ vaut \num{1400}, $B(x) = 200$ ; quand $x$ vaut \num{1490}, $B(x) = 20$ ; quand $x$ vaut \num{1499}, $B(x) = 2$. Le bénéfice \emph{se rapproche} de $0$ au fur et à mesure que $x$ \emph{se rapproche} de \num{1500}.

De même, un moto-taxi qui freine devant un feu rouge à Yaoundé voit sa vitesse se rapprocher de \SI{0}{km/h} : elle passe de \SI{30}{km/h} à \SI{10}{km/h}, puis \SI{3}{km/h}, puis \SI{0,5}{km/h}... La vitesse « tend vers » zéro.

\textbf{Problématique.} Comment les mathématiques décrivent-elles précisément ce comportement de « se rapprocher de » ? Comment prévoir la valeur vers laquelle une quantité se stabilise, ou au contraire reconnaître qu'elle grandit sans limite ? C'est toute l'idée de la notion de \cle{limite}, qui permettra aussi, dans la suite de la leçon, d'étudier si une fonction évolue « sans saut » (on dira qu'elle est \cle{continue}), comme la température au cours d'une journée à Douala, qui ne passe jamais brutalement de \SI{25}{\celsius} à \SI{35}{\celsius}.

\courssub{Se rapprocher d'une valeur finie : limite en un point $a$}

\textbf{Observation graphique.} Considérons une fonction $f$ et un nombre réel $a$. On s'intéresse à la question suivante : lorsque la variable $x$ se rapproche de $a$ (sans forcément être égale à $a$), que devient l'image $f(x)$ ?

Il y a deux façons de se rapprocher de $a$ :
\begin{itemize}
\item \cle{par valeurs inférieures} (on dit aussi « par la gauche ») : $x$ prend les valeurs $a - 1$ ; $a - 0{,}1$ ; $a - 0{,}01$ ; \ldots
\item \cle{par valeurs supérieures} (on dit aussi « par la droite ») : $x$ prend les valeurs $a + 1$ ; $a + 0{,}1$ ; $a + 0{,}01$ ; \ldots
\end{itemize}

Dans le cas du bénéfice de Maman Ngo Bell, $B(x)$ se rapproche de $0$ quand $x$ se rapproche de \num{1500}, aussi bien par la gauche (\num{1490}, \num{1499}) que par la droite (\num{1510} donnerait $B = -20$, une perte proche de zéro). On dit que $B(x)$ a pour \cle{limite} $0$ quand $x$ tend vers \num{1500}.

\begin{definition}[Limite finie en un point — approche intuitive]
Soit $f$ une fonction définie au voisinage d'un réel $a$ (sauf éventuellement en $a$ lui-même) et $\ell$ un nombre réel. Dire que $f(x)$ a pour \cle{limite} $\ell$ quand $x$ tend vers $a$ signifie que $f(x)$ peut être rendu \textbf{aussi proche de $\ell$ qu'on veut}, pourvu que $x$ soit suffisamment proche de $a$. On écrit :
\[ \lim_{x \to a} f(x) = \ell. \]
\end{definition}

Lisons cette définition avec soin. « Aussi proche qu'on veut » veut dire : si l'on exige que $f(x)$ soit à moins de $0{,}01$ de $\ell$, on y arrive en prenant $x$ assez proche de $a$ ; si l'on exige $0{,}0001$, on y arrive aussi, en se rapprochant encore davantage de $a$. La valeur $\ell$ est comme un « aimant » vers lequel les images $f(x)$ sont attirées.

\begin{exemplebox}[Le bénéfice de Maman Ngo Bell]
Avec $B(x) = 3000 - 2x$, vérifions que $\displaystyle\lim_{x \to 1500} B(x) = 0$.
\begin{center}
\begin{tabular}{|c|c|c|c|c|c|c|}
\hline
\rowcolor{popblueL} $x$ & \num{1400} & \num{1490} & \num{1499} & \num{1501} & \num{1510} & \num{1600} \\
\hline $B(x)$ & $200$ & $20$ & $2$ & $-2$ & $-20$ & $-200$ \\
\hline
\end{tabular}
\end{center}
Quand $x$ se rapproche de \num{1500} (par la gauche ou par la droite), $B(x)$ se rapproche de $0$. La limite est bien $0$.
\end{exemplebox}

\courssub{Conjecturer une limite à l'aide d'un tableau de valeurs}

Comment deviner une limite quand on ne la connaît pas à l'avance ? L'outil de base est le \cle{tableau de valeurs} : on calcule $f(x)$ pour des valeurs de $x$ de plus en plus proches de $a$, par la gauche et par la droite, et on observe vers quelle valeur les résultats se stabilisent.

\begin{methode}[Conjecturer une limite en un point]
Pour conjecturer $\displaystyle\lim_{x \to a} f(x)$ :
\begin{enumerate}
\item Compléter un tableau de valeurs en prenant $x$ de plus en plus proche de $a$ \textbf{par valeurs inférieures} : $a - 0{,}1$ ; $a - 0{,}01$ ; $a - 0{,}001$ ; \ldots
\item Faire de même \textbf{par valeurs supérieures} : $a + 0{,}1$ ; $a + 0{,}01$ ; $a + 0{,}001$ ; \ldots
\item Observer les résultats : s'ils se stabilisent tous vers le même nombre $\ell$, on conjecture que la limite est $\ell$.
\item Si les valeurs deviennent de plus en plus grandes (ou de plus en plus négatives) sans se stabiliser, on conjecture une limite infinie.
\end{enumerate}
\end{methode}

\begin{exoresolu}[Une fonction « trouée » qui a pourtant une limite]
\textbf{Énoncé.} Soit $f$ la fonction définie par $f(x) = \dfrac{x^2 - 1}{x - 1}$.
\begin{enumerate}
\item La fonction $f$ est-elle définie en $x = 1$ ?
\item Compléter le tableau de valeurs suivant et conjecturer $\displaystyle\lim_{x \to 1} f(x)$.
\end{enumerate}
\begin{center}
\begin{tabular}{|c|c|c|c|c|c|c|}
\hline
\rowcolor{popblueL} $x$ & $0{,}9$ & $0{,}99$ & $0{,}999$ & $1{,}001$ & $1{,}01$ & $1{,}1$ \\
\hline $f(x)$ & & & & & & \\
\hline
\end{tabular}
\end{center}
\tcblower
\textbf{Solution détaillée.}
\begin{enumerate}
\item En $x = 1$, le dénominateur vaut $1 - 1 = 0$. Or on ne divise jamais par zéro : $f$ \textbf{n'est pas définie en $1$}. La courbe de $f$ possède donc un « trou » au point d'abscisse $1$.
\item Calculons les images. On remarque que $x^2 - 1 = (x-1)(x+1)$, donc pour $x \neq 1$ :
\[ f(x) = \frac{(x-1)(x+1)}{x-1} = x + 1. \]
Ainsi :
\begin{align*}
f(0{,}9) &= 1{,}9 & f(0{,}99) &= 1{,}99 & f(0{,}999) &= 1{,}999 \\
f(1{,}001) &= 2{,}001 & f(1{,}01) &= 2{,}01 & f(1{,}1) &= 2{,}1
\end{align*}
\begin{center}
\begin{tabular}{|c|c|c|c|c|c|c|}
\hline
\rowcolor{popblueL} $x$ & $0{,}9$ & $0{,}99$ & $0{,}999$ & $1{,}001$ & $1{,}01$ & $1{,}1$ \\
\hline $f(x)$ & $1{,}9$ & $1{,}99$ & $1{,}999$ & $2{,}001$ & $2{,}01$ & $2{,}1$ \\
\hline
\end{tabular}
\end{center}
Par la gauche, $f(x)$ vaut $1{,}9$ ; $1{,}99$ ; $1{,}999$ : les valeurs se rapprochent de $2$. Par la droite, $f(x)$ vaut $2{,}1$ ; $2{,}01$ ; $2{,}001$ : elles se rapprochent aussi de $2$. On conjecture donc :
\[ \lim_{x \to 1} \frac{x^2 - 1}{x - 1} = 2. \]
\end{enumerate}
\end{exoresolu}

\begin{attention}[Limite et valeur en $a$ : deux choses différentes]
Une fonction peut \textbf{ne pas être définie en $a$} et pourtant \textbf{admettre une limite en $a$}. C'est exactement le cas de $f(x) = \dfrac{x^2-1}{x-1}$ en $a = 1$ : $f(1)$ n'existe pas, mais $\displaystyle\lim_{x \to 1} f(x) = 2$. Chercher une limite, ce n'est \emph{pas} remplacer $x$ par $a$ ; c'est observer le comportement de $f(x)$ \emph{autour} de $a$. Piège classique au Probatoire : conclure « pas de limite » parce que le dénominateur s'annule. Il faut d'abord transformer l'expression (factoriser, simplifier) avant de conclure.
\end{attention}

\courssub{Limite infinie en un point : le résultat dépasse toute valeur}

Reprenons la fonction \cle{inverse} $f(x) = \dfrac{1}{x}$ et observons ce qui se passe quand $x$ se rapproche de $0$.

\begin{center}
\begin{tabular}{|c|c|c|c|c|c|c|}
\hline
\rowcolor{popblueL} $x$ & $0{,}1$ & $0{,}01$ & $0{,}001$ & $-0{,}001$ & $-0{,}01$ & $-0{,}1$ \\
\hline $f(x) = \dfrac{1}{x}$ & $10$ & $100$ & $1000$ & $-1000$ & $-100$ & $-10$ \\
\hline
\end{tabular}
\end{center}

Quand $x$ se rapproche de $0$ par valeurs \emph{positives}, $f(x)$ prend des valeurs de plus en plus grandes : $10$, $100$, $1000$, et on peut dépasser n'importe quelle valeur fixée à l'avance (un million, un milliard...) en prenant $x$ assez proche de $0$. Les valeurs ne se stabilisent pas : elles « explosent » vers le haut. On dit que $f(x)$ tend vers $+\infty$. Quand $x$ se rapproche de $0$ par valeurs \emph{négatives}, $f(x)$ descend vers $-\infty$.

\begin{definition}[Limite infinie en un point — approche intuitive]
Dire que $f(x)$ tend vers $+\infty$ quand $x$ tend vers $a$ signifie que $f(x)$ \textbf{dépasse toute valeur donnée à l'avance}, aussi grande soit-elle, pourvu que $x$ soit suffisamment proche de $a$. On écrit :
\[ \lim_{x \to a} f(x) = +\infty. \]
On définit de même une limite égale à $-\infty$ quand $f(x)$ devient inférieur à toute valeur donnée.
\end{definition}

\begin{aretenir}[Limite finie ou limite infinie ?]
\begin{itemize}
\item \textbf{Limite finie} $\ell$ : les valeurs de $f(x)$ \emph{se stabilisent} autour d'un nombre précis $\ell$.
\item \textbf{Limite infinie} : les valeurs de $f(x)$ \emph{dépassent toute borne} ; elles grandissent sans fin ($+\infty$) ou décroissent sans fin ($-\infty$).
\end{itemize}
L'infini n'est donc pas un nombre que l'on « atteint » : c'est un \emph{comportement}.
\end{aretenir}

\begin{popfigure}
\begin{center}
\begin{tikzpicture}
\begin{axis}[
  width=0.85\linewidth, height=7cm,
  axis lines=middle, xlabel={$x$}, ylabel={$y$},
  xmin=-4, xmax=4, ymin=-5, ymax=5,
  xtick={-3,-2,-1,1,2,3}, ytick={-4,-2,2,4},
  samples=200, clip=true]
\addplot[popcurve,domain=0.22:4] {1/x};
\addplot[popcurve,domain=-4:-0.22] {1/x};
\draw[poporange, dashed] (axis cs:0,-5) -- (axis cs:0,5);
\draw[popgreen, dashed] (axis cs:-4,0) -- (axis cs:4,0);
\draw[-{Stealth}, poporange, thick] (axis cs:1.6,0.35) -- (axis cs:0.35,0.35) node[midway, above, font=\small] {$x \to 0^{+}$};
\draw[-{Stealth}, poporange, thick] (axis cs:-1.6,-0.35) -- (axis cs:-0.35,-0.35) node[midway, below, font=\small] {$x \to 0^{-}$};
\draw[-{Stealth}, poppurple, thick] (axis cs:2.2,1.6) -- (axis cs:3.6,0.5) node[midway, above right=-2pt, font=\small] {$x \to +\infty$};
\draw[-{Stealth}, poppurple, thick] (axis cs:-2.2,-1.6) -- (axis cs:-3.6,-0.5) node[midway, below left=-2pt, font=\small] {$x \to -\infty$};
\end{axis}
\end{tikzpicture}
\end{center}
\legende{Courbe de $f(x) = \dfrac{1}{x}$. Quand $x \to 0^{+}$, la courbe monte vers $+\infty$ ; quand $x \to 0^{-}$, elle descend vers $-\infty$ ; quand $x \to \pm\infty$, elle se rapproche de l'axe des abscisses : la limite est $0$.}
\end{popfigure}

\courssub{Comportement à l'infini : quand $x$ devient très grand}

Jusqu'ici, $x$ se rapprochait d'un nombre $a$. On peut aussi faire grandir $x$ \emph{sans limite} : $x$ prend les valeurs $10$ ; $100$ ; $1000$ ; $10\,000$ ; \ldots On dit alors que $x$ \cle{tend vers} $+\infty$, et on écrit $x \to +\infty$. De même, $x$ peut prendre des valeurs négatives de plus en plus grandes en valeur absolue : $-10$ ; $-100$ ; $-1000$ ; \ldots ; on écrit $x \to -\infty$.

La question est alors : que devient $f(x)$ ?

\begin{exemplebox}[Deux comportements opposés à l'infini]
\textbf{Cas 1 : $f(x) = \dfrac{1}{x}$.}
\begin{center}
\begin{tabular}{|c|c|c|c|c|}
\hline
\rowcolor{popblueL} $x$ & $10$ & $100$ & $1000$ & $100\,000$ \\
\hline $f(x)$ & $0{,}1$ & $0{,}01$ & $0{,}001$ & $0{,}00001$ \\
\hline
\end{tabular}
\end{center}
Les valeurs se rapprochent de $0$ : $\displaystyle\lim_{x \to +\infty} \frac{1}{x} = 0$. De même $\displaystyle\lim_{x \to -\infty} \frac{1}{x} = 0$.

\textbf{Cas 2 : $f(x) = x^2$.}
\begin{center}
\begin{tabular}{|c|c|c|c|c|}
\hline
\rowcolor{popblueL} $x$ & $10$ & $100$ & $1000$ & $-1000$ \\
\hline $f(x)$ & $100$ & $10\,000$ & $1\,000\,000$ & $1\,000\,000$ \\
\hline
\end{tabular}
\end{center}
Les valeurs dépassent toute borne : $\displaystyle\lim_{x \to +\infty} x^2 = +\infty$ et $\displaystyle\lim_{x \to -\infty} x^2 = +\infty$ (le carré d'un nombre négatif est positif).
\end{exemplebox}

\begin{aretenir}[Vocabulaire et notations]
\begin{itemize}
\item $x \to a$ : $x$ se rapproche du nombre $a$ (par la gauche et/ou par la droite).
\item $x \to +\infty$ : $x$ devient aussi grand qu'on veut ; $x \to -\infty$ : $x$ devient négatif, aussi grand qu'on veut en valeur absolue.
\item $\displaystyle\lim_{x \to a} f(x) = \ell$ : limite \textbf{finie}, $f(x)$ se stabilise vers le nombre $\ell$.
\item $\displaystyle\lim_{x \to a} f(x) = +\infty$ (ou $-\infty$) : limite \textbf{infinie}, $f(x)$ dépasse toute valeur.
\end{itemize}
\end{aretenir}

\courssub{Lecture graphique de limites}

Sur une courbe donnée, on lit les limites en observant le comportement des extrémités et des « cassures » :
\begin{itemize}
\item Si la courbe \emph{se plaque} contre une droite horizontale d'équation $y = \ell$ quand $x \to +\infty$ ou $x \to -\infty$, alors la limite est $\ell$ (c'est le cas de $y = 0$ pour la fonction inverse).
\item Si la courbe \emph{monte ou descend sans fin} quand $x$ se rapproche d'une valeur $a$ (comme $\frac{1}{x}$ près de $0$), la limite est infinie.
\item Si la courbe \emph{se rapproche d'un point} $(a\,;\,\ell)$, même si ce point n'appartient pas à la courbe (le « trou » de $\frac{x^2-1}{x-1}$ en $1$), la limite en $a$ est $\ell$.
\end{itemize}

\begin{exercice}[Conjectures par tableaux de valeurs]
\begin{enumerate}
\item À l'aide d'un tableau de valeurs ($x = 1{,}9$ ; $1{,}99$ ; $1{,}999$ ; $2{,}001$ ; $2{,}01$ ; $2{,}1$), conjecturer $\displaystyle\lim_{x \to 2} (3x - 1)$.
\item Soit $g(x) = \dfrac{1}{x - 3}$. Compléter un tableau de valeurs pour $x$ proche de $3$ par valeurs supérieures, puis conjecturer $\displaystyle\lim_{\substack{x \to 3 \\ x > 3}} g(x)$.
\item Conjecturer $\displaystyle\lim_{x \to +\infty} \dfrac{2}{x}$ et $\displaystyle\lim_{x \to +\infty} (x^2 + 1)$.
\end{enumerate}
\end{exercice}

\begin{corrige}[Exercice : Conjectures par tableaux de valeurs]
\begin{enumerate}
\item $f(1{,}9) = 4{,}7$ ; $f(1{,}99) = 4{,}97$ ; $f(1{,}999) = 4{,}997$ ; $f(2{,}001) = 5{,}003$ ; $f(2{,}01) = 5{,}03$ ; $f(2{,}1) = 5{,}3$. Les valeurs se stabilisent vers $5$ : $\displaystyle\lim_{x \to 2}(3x-1) = 5$.
\item $g(3{,}1) = 10$ ; $g(3{,}01) = 100$ ; $g(3{,}001) = 1000$. Les valeurs dépassent toute borne : la limite est $+\infty$.
\item $\dfrac{2}{x}$ vaut $0{,}02$ ; $0{,}002$ ; $0{,}0002$ pour $x = 100$ ; $1000$ ; $10\,000$ : la limite est $0$. Pour $x^2 + 1$, les valeurs grandissent sans fin : la limite est $+\infty$.
\end{enumerate}
\end{corrige}

\begin{savaistu}[D'où viennent « lim » et le symbole $\infty$ ?]
Le symbole $\infty$ a été introduit par le mathématicien anglais John Wallis en 1655, et la notation « lim » a été popularisée au XVIII\textsuperscript{e} siècle, notamment grâce au Suisse Leonhard Euler, l'un des plus grands mathématiciens de l'histoire. Ces outils ont permis de fonder l'analyse moderne, celle qui décrit les mouvements, les vitesses et les variations. Retiens bien : l'infini \textbf{n'est pas un nombre}, c'est un \emph{comportement} : celui d'une quantité qui dépasse toute valeur qu'on pourrait fixer. C'est pourquoi on ne peut pas faire des calculs ordinaires avec $\infty$ comme avec $3$ ou $-5$ — nous verrons dans la suite de la leçon les règles précises pour calculer avec les limites.
\end{savaistu}

\coursec{Opérations sur les limites}

\courssub{Règles de calcul avec des limites finies}

Lorsque les limites en jeu sont des nombres réels (finis), les opérations habituelles se transmettent à la limite, sous réserve que les expressions aient un sens.

\begin{propriete}[Théorèmes sur les opérations — Limites finies]
Soient $f$ et $g$ deux fonctions et $a \in \mathbb{R} \cup \{+\infty, -\infty\}$. On suppose que
\[ \lim_{x \to a} f(x) = \ell \in \mathbb{R} \quad \text{et} \quad \lim_{x \to a} g(x) = \ell' \in \mathbb{R}. \]
Alors :
\begin{itemize}
\item \textbf{Somme :} $\displaystyle\lim_{x \to a} \bigl(f(x) + g(x)\bigr) = \ell + \ell'$.
\item \textbf{Produit :} $\displaystyle\lim_{x \to a} \bigl(f(x) \cdot g(x)\bigr) = \ell \cdot \ell'$.
\item \textbf{Inverse :} Si $\ell' \neq 0$, $\displaystyle\lim_{x \to a} \frac{1}{g(x)} = \frac{1}{\ell'}$.
\item \textbf{Quotient :} Si $\ell' \neq 0$, $\displaystyle\lim_{x \to a} \frac{f(x)}{g(x)} = \frac{\ell}{\ell'}$.
\item \textbf{Composition :} Si de plus $g$ est définie au voisinage de $\ell$ et $\displaystyle\lim_{y \to \ell} g(y) = L$, alors $\displaystyle\lim_{x \to a} g(f(x)) = L$.
\end{itemize}
\textit{(Démonstrations non exigibles au Probatoire, mais les conditions d'application le sont.)}
\end{propriete}

\begin{exemplebox}[Application directe des règles]
Soit $f(x) = 3x^2 - 5x + 2$. Calculer $\displaystyle\lim_{x \to 2} f(x)$.
\[ \lim_{x \to 2} (3x^2 - 5x + 2) = 3 \cdot (\lim_{x \to 2} x)^2 - 5 \cdot (\lim_{x \to 2} x) + 2 = 3 \cdot 4 - 5 \cdot 2 + 2 = 12 - 10 + 2 = 4. \]
Pour une fonction polynôme $P$, on a toujours $\displaystyle\lim_{x \to a} P(x) = P(a)$.
\end{exemplebox}

\courssub{Règles de calcul avec des limites infinies}

Lorsque l'une au moins des limites est infinie, des règles s'appliquent encore, mais d'autres situations conduisent à des \cle{formes indéterminées} qu'il faudra lever.

\begin{propriete}[Règles opératoires avec $+\infty$ et $-\infty$]
Soient $\ell \in \mathbb{R}$ (non nul pour les produits/quotients) et les limites infinies. On adopte les conventions suivantes (valables pour $x \to a$ avec $a \in \mathbb{R} \cup \{\pm\infty\}$) :
\begin{center}
\begin{tabular}{|c|c|c|}
\hline
\rowcolor{popblueL} Opération & Résultat & Condition \\
\hline
$\ell + (+\infty)$ & $+\infty$ & $\ell \in \mathbb{R}$ \\
$\ell + (-\infty)$ & $-\infty$ & $\ell \in \mathbb{R}$ \\
$(+\infty) + (+\infty)$ & $+\infty$ & \\
$(-\infty) + (-\infty)$ & $-\infty$ & \\
\hline
$\ell \times (+\infty)$ & $+\infty$ si $\ell > 0$, $-\infty$ si $\ell < 0$ & $\ell \neq 0$ \\
$\ell \times (-\infty)$ & $-\infty$ si $\ell > 0$, $+\infty$ si $\ell < 0$ & $\ell \neq 0$ \\
$(+\infty) \times (+\infty)$ & $+\infty$ & \\
$(+\infty) \times (-\infty)$ & $-\infty$ & \\
$(-\infty) \times (-\infty)$ & $+\infty$ & \\
\hline
$\dfrac{\ell}{+\infty} = \dfrac{\ell}{-\infty}$ & $0$ & $\ell \in \mathbb{R}$ \\
$\dfrac{+\infty}{\ell}$ & $+\infty$ si $\ell > 0$, $-\infty$ si $\ell < 0$ & $\ell \neq 0$ \\
$\dfrac{-\infty}{\ell}$ & $-\infty$ si $\ell > 0$, $+\infty$ si $\ell < 0$ & $\ell \neq 0$ \\
\hline
\end{tabular}
\end{center}
\end{propriete}

\begin{attention}[Formes indéterminées — Interdiction de calculer]
Les combinaisons suivantes \textbf{n'ont pas de sens défini} (on ne peut pas conclure) : ce sont les \cle{formes indéterminées}.
\begin{itemize}
\item $+\infty - \infty$ (ou $-\infty + \infty$)
\item $0 \times \infty$ (ou $0 \times (+\infty)$, $0 \times (-\infty)$)
\item $\dfrac{0}{0}$
\item $\dfrac{\infty}{\infty}$ (ou $\dfrac{+\infty}{+\infty}$, $\dfrac{-\infty}{-\infty}$, etc.)
\item $1^{\infty}$ (rencontré plus tard avec les exponentielles)
\end{itemize}
\textbf{Au Probatoire :} Si vous obtenez une forme indéterminée, vous \textbf{devez} transformer l'expression (factorisation, mise au même dénominateur, conjugaison, changement de variable) avant de recalculer la limite. Écrire « $= \frac{0}{0} $ » comme résultat final est une erreur grave.
\end{attention}

\begin{exoresolu}[Application des règles et détection de formes indéterminées]
On suppose $\displaystyle\lim_{x \to +\infty} f(x) = +\infty$ et $\displaystyle\lim_{x \to +\infty} g(x) = 3$.
Déterminer, si possible, les limites suivantes :
\begin{enumerate}
\item $\displaystyle\lim_{x \to +\infty} \bigl(f(x) + g(x)\bigr)$
\item $\displaystyle\lim_{x \to +\infty} \bigl(f(x) - g(x)\bigr)$
\item $\displaystyle\lim_{x \to +\infty} \dfrac{g(x)}{f(x)}$
\item $\displaystyle\lim_{x \to +\infty} \dfrac{f(x)}{g(x)}$
\item $\displaystyle\lim_{x \to +\infty} \bigl(f(x) \cdot (g(x) - 3)\bigr)$
\end{enumerate}
\tcblower
\textbf{Solution.}
\begin{enumerate}
\item $+\infty + 3 = +\infty$.
\item $+\infty - 3 = +\infty$.
\item $\dfrac{3}{+\infty} = 0$.
\item $\dfrac{+\infty}{3} = +\infty$.
\item $f(x) \to +\infty$ et $(g(x)-3) \to 0$. C'est une forme indéterminée $+\infty \times 0$ (ou $0 \times \infty$). \textbf{On ne peut pas conclure} sans connaître la forme exacte de $f$ et $g$.
\end{enumerate}
\end{exoresolu}

\coursec{Calculs pratiques : lever les formes indéterminées}

\courssub{Forme $\dfrac{0}{0}$ : factorisation et simplification}

C'est la forme la plus fréquente pour les fonctions rationnelles ou les expressions algébriques en un point $a$ où numérateur et dénominateur s'annulent.

\begin{methode}[Lever une forme $\frac{0}{0}$ en $a$]
\begin{enumerate}
\item Factoriser le numérateur et le dénominateur (identités remarquables : $x^2-a^2=(x-a)(x+a)$, $x^3-a^3=(x-a)(x^2+ax+a^2)$, etc.).
\item Simplifier le facteur commun $(x-a)$ (ou l'expression qui tend vers $0$).
\item Calculer la limite de l'expression simplifiée par substitution directe.
\end{enumerate}
\end{methode}

\begin{exoresolu}[Forme $\frac{0}{0}$ par factorisation]
Calculer $\displaystyle\lim_{x \to 2} \frac{x^3 - 8}{x^2 - 4}$.
\tcblower
\textbf{Solution.}
Pour $x=2$, numérateur $= 0$, dénominateur $= 0$ : forme $\frac{0}{0}$.
\begin{align*}
x^3 - 8 &= (x-2)(x^2 + 2x + 4) \\
x^2 - 4 &= (x-2)(x+2)
\end{align*}
Pour $x \neq 2$ :
\[ \frac{x^3 - 8}{x^2 - 4} = \frac{(x-2)(x^2+2x+4)}{(x-2)(x+2)} = \frac{x^2+2x+4}{x+2}. \]
Maintenant, la limite se calcule directement :
\[ \lim_{x \to 2} \frac{x^2+2x+4}{x+2} = \frac{4+4+4}{4} = \frac{12}{4} = 3. \]
\end{exoresolu}

\courssub{Forme $\dfrac{\infty}{\infty}$ : facteur de plus haut degré}

Pour les limites à $+\infty$ ou $-\infty$ de fractions rationnelles, on factorise le terme de plus haut degré au numérateur et au dénominateur.

\begin{methode}[Lever une forme $\frac{\infty}{\infty}$ à $+\infty$ ou $-\infty$]
Soit $R(x) = \frac{P(x)}{Q(x)}$ une fraction rationnelle ($\deg P = p$, $\deg Q = q$).
\begin{enumerate}
\item Factoriser $x^p$ au numérateur et $x^q$ au dénominateur.
\item Simplifier les puissances de $x$ communes.
\item Utiliser $\lim_{x \to \pm\infty} \frac{1}{x^k} = 0$ pour $k>0$.
\item Conclure selon la comparaison des degrés :
  \begin{itemize}
  \item $p < q$ : limite $0$.
  \item $p = q$ : limite = quotient des coefficients dominants.
  \item $p > q$ : limite $\pm\infty$ (signe selon parité de $p-q$ et signe de $x$).
  \end{itemize}
\end{enumerate}
\end{methode}

\begin{exoresolu}[Limite à l'infini d'une fraction rationnelle]
Calculer $\displaystyle\lim_{x \to +\infty} \frac{3x^2 - 5x + 1}{2x^2 + 7}$.
\tcblower
\textbf{Solution.}
Forme $\frac{\infty}{\infty}$. Degrés égaux ($p=q=2$). Factorisons $x^2$ :
\[ \frac{3x^2 - 5x + 1}{2x^2 + 7} = \frac{x^2(3 - \frac{5}{x} + \frac{1}{x^2})}{x^2(2 + \frac{7}{x^2})} = \frac{3 - \frac{5}{x} + \frac{1}{x^2}}{2 + \frac{7}{x^2}}. \]
Quand $x \to +\infty$, $\frac{5}{x} \to 0$, $\frac{1}{x^2} \to 0$, $\frac{7}{x^2} \to 0$.
La limite vaut $\frac{3}{2}$.
\end{exoresolu}

\courssub{Forme $0 \times \infty$ : réécriture en quotient}

On ramène le produit à une fraction pour retomber sur $\frac{0}{0}$ ou $\frac{\infty}{\infty}$.

\begin{methode}[Lever une forme $0 \times \infty$]
Si $\lim f(x) = 0$ et $\lim g(x) = \infty$, écrire $f(x)g(x) = \frac{f(x)}{1/g(x)}$ (forme $\frac{0}{0}$) ou $\frac{g(x)}{1/f(x)}$ (forme $\frac{\infty}{\infty}$). Choisir la forme la plus simple à traiter.
\end{methode}

\begin{exoresolu}[Forme $0 \times \infty$]
Calculer $\displaystyle\lim_{x \to 0^{+}} x \ln x$ (on admet $\lim_{x \to 0^{+}} \ln x = -\infty$).
\tcblower
\textbf{Solution.}
Forme $0 \times (-\infty)$. On écrit $x \ln x = \frac{\ln x}{1/x}$.
Numérateur $\to -\infty$, dénominateur $1/x \to +\infty$ : forme $\frac{-\infty}{+\infty}$.
On ne peut pas dériver (pas au programme), mais on connaît la limite classique : $\lim_{x \to 0^{+}} x \ln x = 0$.
\textit{(Note : au Probatoire, on utilise souvent le changement de variable $x = 1/t$, $t \to +\infty$, donnant $\frac{-\ln t}{t} \to 0$ car $\ln t$ croît moins vite que $t$).}
\end{exoresolu}

\courssub{Forme $\infty - \infty$ : mise au même dénominateur ou factorisation}

Souvent rencontrée avec des différences de racines carrées ou de fractions rationnelles.

\begin{methode}[Lever une forme $\infty - \infty$]
\begin{enumerate}
\item Si ce sont des fractions : mettre au même dénominateur.
\item Si ce sont des racines : multiplier par le conjugué (ex: $\sqrt{A} - \sqrt{B} = \frac{A-B}{\sqrt{A}+\sqrt{B}}$).
\item Simplifier et recalculer.
\end{enumerate}
\end{methode}

\begin{exoresolu}[Différence de racines : forme $\infty - \infty$]
Calculer $\displaystyle\lim_{x \to +\infty} \bigl(\sqrt{x^2 + x} - x\bigr)$.
\tcblower
\textbf{Solution.}
Forme $+\infty - \infty$. Multiplier par le conjugué :
\[ \sqrt{x^2+x} - x = \frac{(\sqrt{x^2+x} - x)(\sqrt{x^2+x} + x)}{\sqrt{x^2+x} + x} = \frac{x^2+x - x^2}{\sqrt{x^2+x} + x} = \frac{x}{\sqrt{x^2+x} + x}. \]
Factoriser $x$ au dénominateur ($x>0$ donc $\sqrt{x^2}=x$) :
\[ \frac{x}{x\sqrt{1+\frac{1}{x}} + x} = \frac{x}{x(\sqrt{1+\frac{1}{x}} + 1)} = \frac{1}{\sqrt{1+\frac{1}{x}} + 1}. \]
Quand $x \to +\infty$, $\frac{1}{x} \to 0$, donc $\sqrt{1+0}+1 = 2$.
La limite est $\frac{1}{2}$.
\end{exoresolu}

\begin{exercice}[Entraînement aux formes indéterminées]
Calculer les limites suivantes :
\begin{enumerate}
\item $\displaystyle\lim_{x \to 3} \frac{x^2 - 9}{x - 3}$
\item $\displaystyle\lim_{x \to +\infty} \frac{5x^3 - 2x}{2x^3 + 4x^2 - 1}$
\item $\displaystyle\lim_{x \to 1} \frac{\sqrt{x} - 1}{x - 1}$ (indication : conjugué)
\item $\displaystyle\lim_{x \to +\infty} \bigl(\sqrt{x^2 + 3x} - x\bigr)$
\end{enumerate}
\end{exercice}

\begin{corrige}[Exercice : Entraînement aux formes indéterminées]
\begin{enumerate}
\item $\frac{(x-3)(x+3)}{x-3} = x+3 \to 6$.
\item Degrés égaux (3), coeff dominants $5$ et $2$ : limite $5/2$.
\item $\frac{\sqrt{x}-1}{x-1} = \frac{\sqrt{x}-1}{(\sqrt{x}-1)(\sqrt{x}+1)} = \frac{1}{\sqrt{x}+1} \to \frac{1}{2}$.
\item Conjugué : $\frac{3x}{\sqrt{x^2+3x}+x} = \frac{3}{\sqrt{1+3/x}+1} \to \frac{3}{2}$.
\end{enumerate}
\end{corrige}

\coursec{Continuité d'une fonction}

\courssub{Définition et interprétation graphique}

L'intuition : une fonction est continue si on peut tracer sa courbe sans lever le crayon. Mathématiquement, cela signifie que la limite en un point vaut la valeur de la fonction en ce point.

\begin{definition}[Continuité en un point]
Soit $f$ une fonction définie sur un intervalle $I$ et $a \in I$. On dit que $f$ est \cle{continue en $a$} si :
\[ \lim_{x \to a} f(x) = f(a). \]
Cela implique trois conditions :
\begin{enumerate}
\item $f$ est définie en $a$ ($f(a)$ existe).
\item $f$ admet une limite finie $\ell$ en $a$.
\item Cette limite est égale à $f(a)$ : $\ell = f(a)$.
\end{enumerate}
Si l'une de ces conditions n'est pas vérifiée, $f$ est \cle{discontinue} en $a$.
\end{definition}

\begin{exemplebox}[Types de discontinuités]
\begin{enumerate}
\item $f(x) = \frac{1}{x}$ en $0$ : non définie $\Rightarrow$ discontinue (asymptote verticale).
\item $f(x) = \frac{x^2-1}{x-1}$ en $1$ : limite $=2$ mais $f(1)$ n'existe pas $\Rightarrow$ discontinue (trou « évitable »).
\item $f(x) = \begin{cases} x+1 & \text{si } x < 1 \\ 3 & \text{si } x = 1 \\ x+1 & \text{si } x > 1 \end{cases}$ en $1$ : limite $=2$, $f(1)=3$ $\Rightarrow$ discontinue (saut).
\end{enumerate}
\end{exemplebox}

\begin{propriete}[Continuité sur un intervalle]
Une fonction est \cle{continue sur un intervalle $I$} si elle est continue en tout point de $I$.
Les fonctions usuelles sont continues sur leur ensemble de définition :
\begin{itemize}
\item Fonctions polynômes : continues sur $\mathbb{R}$.
\item Fonctions rationnelles : continues sur $\mathbb{R} \setminus \{\text{zéros du dénominateur}\}$.
\item Fonction racine carrée : continue sur $[0, +\infty[$.
\item Fonctions trigonométriques ($\sin, \cos$) : continues sur $\mathbb{R}$.
\end{itemize}
\end{propriete}

\begin{propriete}[Opérations sur les fonctions continues]
Si $f$ et $g$ sont continues en $a$ (ou sur $I$), alors :
$f+g$, $f-g$, $f \times g$ sont continues en $a$ (sur $I$).
$\frac{f}{g}$ est continue en $a$ (sur $I$) \textbf{si $g(a) \neq 0$}.
La composée $g \circ f$ est continue en $a$ si $f$ est continue en $a$ et $g$ est continue en $f(a)$.
\end{propriete}

\courssub{Théorème des valeurs intermédiaires (TVI) — Application majeure}

C'est un outil puissant pour prouver l'existence de solutions d'équations $f(x) = k$ sans les calculer explicitement.

\begin{propriete}[Théorème des valeurs intermédiaires]
Soit $f$ une fonction \textbf{continue sur un intervalle fermé $[a,b]$}. Soit $k$ un nombre réel compris entre $f(a)$ et $f(b)$ (c'est-à-dire $f(a) \le k \le f(b)$ ou $f(b) \le k \le f(a)$).
Alors il existe \textbf{au moins un} $c \in [a,b]$ tel que $f(c) = k$.
\end{propriete}

\begin{attention}[Conditions d'application du TVI]
Trois conditions \textbf{indispensables} (souvent oubliées au Probatoire) :
\begin{enumerate}
\item $f$ est \textbf{continue} sur $[a,b]$ (vérifier l'absence de trou, d'asymptote, de saut).
\item L'intervalle est \textbf{fermé} $[a,b]$ (pas $]a,b[$, pas $[a,+\infty[$).
\item $k$ est \textbf{entre} $f(a)$ et $f(b)$ (vérifier l'ordre : $f(a) \le k \le f(b)$ ou l'inverse).
\end{enumerate}
Si $f$ est \textbf{strictement monotone} sur $[a,b]$, le $c$ est \textbf{unique}.
\end{attention}

\begin{exoresolu}[Application du TVI : recherche d'un zéro]
Montrer que l'équation $x^3 - 3x + 1 = 0$ admet une solution dans l'intervalle $[0,1]$.
\tcblower
\textbf{Solution.}
Posons $f(x) = x^3 - 3x + 1$.
\begin{enumerate}
\item $f$ est une fonction polynôme, donc \textbf{continue sur $\mathbb{R}$}, a fortiori sur $[0,1]$.
\item $f(0) = 1 > 0$ et $f(1) = -1 < 0$.
\item $0$ est bien entre $f(0)=1$ et $f(1)=-1$ (on a $f(1) \le 0 \le f(0)$).
\end{enumerate}
Par le TVI, il existe $c \in [0,1]$ tel que $f(c)=0$. L'équation admet bien une solution dans $[0,1]$.
\end{exoresolu}

\coursec{Applications à des situations concrètes}

\courssub{Électricité : charge d'un condensateur}

Dans un circuit RC série alimenté par une tension continue $E$, la tension $u_C(t)$ aux bornes du condensateur au cours du temps $t \ge 0$ est modélisée par :
\[ u_C(t) = E \left(1 - e^{-t/\tau}\right), \quad \tau = RC \text{ (constante de temps).} \]

\begin{exoresolu}[Analyse de la charge]
\begin{enumerate}
\item Calculer $\lim_{t \to 0^{+}} u_C(t)$ et interpréter physiquement.
\item Calculer $\lim_{t \to +\infty} u_C(t)$ et interpréter.
\item La fonction $u_C$ est-elle continue sur $[0, +\infty[$ ?
\end{enumerate}
\tcblower
\textbf{Solution.}
\begin{enumerate}
\item $u_C(0) = E(1 - e^0) = 0$. $\lim_{t \to 0^{+}} u_C(t) = 0$. Au démarrage ($t=0$), le condensateur est déchargé, tension nulle.
\item Quand $t \to +\infty$, $e^{-t/\tau} \to 0$, donc $u_C(t) \to E$. À l'infini, le condensateur est chargé à la tension source $E$. C'est le \cle{régime permanent}.
\item $u_C$ est composée de fonctions continues ($t \mapsto t$, exponentielle, opérations) sur $[0,+\infty[$. Elle est donc \textbf{continue sur $[0,+\infty[$}. La tension varie sans saut brutal.
\end{enumerate}
\end{exoresolu}

\courssub{Mécanique : coût de production et coût moyen}

Une usine produit $x$ pièces (en milliers). Le coût total (en millions de FCFA) est modélisé par $C_T(x) = 2x^2 + 50x + 1000$ pour $x > 0$. Le coût moyen par pièce (en FCFA) est $C_m(x) = \frac{C_T(x)}{x} \times 1000$ (conversion millions $\to$ unités).

\begin{exoresolu}[Comportement du coût moyen]
\begin{enumerate}
\item Exprimer $C_m(x)$ sous forme simplifiée.
\item Calculer $\lim_{x \to 0^{+}} C_m(x)$ et $\lim_{x \to +\infty} C_m(x)$. Interpréter.
\item En déduire l'existence d'un minimum du coût moyen (sans le calculer).
\end{enumerate}
\tcblower
\textbf{Solution.}
\begin{enumerate}
\item $C_m(x) = \frac{2x^2 + 50x + 1000}{x} \times 1000 = (2x + 50 + \frac{1000}{x}) \times 1000$ FCFA/pièce.
\item $\lim_{x \to 0^{+}} C_m(x) = +\infty$ (terme $1000/x$). Produire très peu coûte très cher par pièce (coûts fixes non amortis).
$\lim_{x \to +\infty} C_m(x) = +\infty$ (terme $2x$ domine). Produire trop coûte cher par pièce (surcoûts, heures sup, usure).
\item $C_m$ est continue sur $]0, +\infty[$ (somme de fonctions continues). Elle tend vers $+\infty$ aux deux bornes. Par continuité, elle admet un \textbf{minimum global} sur $]0,+\infty[$ (théorème des bornes / valeurs intermédiaires appliqué à la dérivée plus tard, ou raisonnement par l'absurde : si pas de min, elle serait $\ge M$ puis redescendrait...). Il existe un volume de production optimal minimisant le coût unitaire.
\end{enumerate}
\end{exoresolu}

\courssub{Gestion : seuil de rentabilité}

Une entreprise vend un produit au prix unitaire $p = \SI{5000}{FCFA}$. Ses coûts fixes sont $CF = \SI{1\,000\,000}{FCFA}$ et le coût variable unitaire est $c_v = \SI{3000}{FCFA}$. Le bénéfice pour $x$ unités vendues est $B(x) = (p - c_v)x - CF = 2000x - 1\,000\,000$.

\begin{exoresolu}[Seuil de rentabilité et limite]
\begin{enumerate}
\item Calculer $\lim_{x \to +\infty} B(x)$. Interpréter.
\item Déterminer le seuil de rentabilité $x_0$ (bénéfice nul).
\item Étudier la continuité de $B$ et le signe de $B(x)$ selon $x$.
\end{enumerate}
\tcblower
\textbf{Solution.}
\begin{enumerate}
\item $\lim_{x \to +\infty} B(x) = +\infty$. Si l'on vend indéfiniment, le bénéfice croît sans limite (modèle linéaire simplifié).
\item $B(x_0) = 0 \iff 2000x_0 = 1\,000\,000 \iff x_0 = 500$ unités.
\item $B$ est affine, donc continue sur $\mathbb{R}$. $B(x) < 0$ pour $x < 500$ (perte), $B(x) = 0$ pour $x=500$, $B(x) > 0$ pour $x > 500$ (bénéfice). Le passage de la perte au bénéfice est \textbf{continu} (pas de saut), ce qui correspond à la réalité économique progressive.
\end{enumerate}
\end{exoresolu}

\begin{exercice}[Synthèse : Limites, continuité, applications]
\begin{enumerate}
\item Soit $f(x) = \frac{x^2 - 4}{x - 2}$.
  \begin{enumerate}
  \item Déterminer $\lim_{x \to 2} f(x)$.
  \item Peut-on définir $f(2)$ pour que $f$ devienne continue en $2$ ? Si oui, quelle valeur donner ?
  \end{enumerate}
\item Un réservoir d'eau se vide. Le volume $V(t)$ (en $m^3$) à l'instant $t$ (en heures) est $V(t) = 100 - 5t^2$ pour $0 \le t \le \sqrt{20}$.
  \begin{enumerate}
  \item Calculer $\lim_{t \to \sqrt{20}^{-}} V(t)$. Que se passe-t-il physiquement ?
  \item $V$ est-elle continue sur $[0, \sqrt{20}]$ ?
  \end{enumerate}
\item Soit $g(x) = \frac{2x+1}{x-3}$.
  \begin{enumerate}
  \item Calculer $\lim_{x \to 3^{-}} g(x)$ et $\lim_{x \to 3^{+}} g(x)$.
  \item Calculer $\lim_{x \to +\infty} g(x)$.
  \item Tracer schématiquement la courbe de $g$ (asymptotes, intersection axes).
  \end{enumerate}
\end{enumerate}
\end{exercice}

\begin{corrige}[Exercice : Synthèse]
\begin{enumerate}
\item \begin{enumerate}
  \item $f(x) = \frac{(x-2)(x+2)}{x-2} = x+2$ pour $x \neq 2$. $\lim_{x \to 2} f(x) = 4$.
  \item Oui, en posant $f(2) = 4$. La discontinuité est « évitable » (trou).
  \end{enumerate}
\item \begin{enumerate}
  \item $\lim_{t \to \sqrt{20}^{-}} (100 - 5t^2) = 100 - 5 \times 20 = 0$. Le réservoir est vide.
  \item Oui, $V$ est polynôme, continue sur $\mathbb{R}$, donc sur $[0, \sqrt{20}]$.
  \end{enumerate}
\item \begin{enumerate}
  \item $x \to 3^{-}$ : numérateur $\to 7$, dénominateur $\to 0^{-}$ $\Rightarrow$ $-\infty$.
  $x \to 3^{+}$ : dénominateur $\to 0^{+}$ $\Rightarrow$ $+\infty$. Asymptote verticale $x=3$.
  \item $\lim_{x \to +\infty} \frac{2x+1}{x-3} = \lim \frac{2+1/x}{1-3/x} = 2$. Asymptote horizontale $y=2$.
  \item Courbe : hyperbole, centre $(3,2)$, branches dans quadrants opposés par rapport au centre. Passe par $(0, -1/3)$ et $(-0.5, 0)$.
  \end{enumerate}
\end{enumerate}
\end{corrige}

\begin{savaistu}[Continuité et monde réel : le mythe du saut]
En physique, on dit souvent « la nature ne fait pas de sauts ». La température, la pression, la position d'un mobile varient continûment. Les discontinuités dans nos modèles (asymptotes, sauts) signalent souvent une \textbf{rupture de validité du modèle} (changement de phase, rupture de matériau, saturation d'un composant électronique). Au Probatoire, on vous demandera souvent de \textbf{justifier la continuité} avant d'appliquer le TVI : « $f$ est continue sur $[a,b]$ car c'est une fonction polynôme / rationnelle dont le dénominateur ne s'annule pas sur $[a,b]$ / composée de fonctions continues ». Cette phrase vaut des points !
\end{savaistu}

\coursec{Limites de référence et interprétation graphique}

Dans la section précédente, nous avons découvert de manière \emph{intuitive} ce que signifie « une fonction tend vers une valeur » : on a observé des tableaux de valeurs, des courbes, et l'on a accepté l'idée que $f(x)$ peut s'approcher aussi près que l'on veut d'un nombre $\ell$ quand $x$ devient très grand, ou quand $x$ s'approche d'un point $a$. Nous allons maintenant donner du \cle{corps} à cette intuition : d'abord en établissant une liste de \cle{limites de référence} que l'on pourra utiliser sans les redémontrer, ensuite en montrant comment ces limites se lisent \emph{géométriquement} sur une courbe grâce à la notion d'\cle{asymptote}. C'est cette double compétence — connaître les limites de base et les interpréter graphiquement — qui servira de fondation à tous les calculs des sections suivantes.

\courssub{Limites en $+\infty$ et $-\infty$ des fonctions puissances $x \mapsto x^n$}

Commençons par les fonctions les plus simples : les \cle{fonctions puissances} $x \mapsto x^n$, où $n$ est un entier naturel non nul ($n = 1, 2, 3, \ldots$). L'intuition est immédiate : si $x$ devient un nombre positif de plus en plus grand, alors $x^2$, $x^3$, $x^4$ deviennent eux aussi de plus en plus grands, et même de plus en plus vite. Par exemple, pour $x = 100$, on a $x^2 = \num{10000}$ et $x^3 = \num{1000000}$.

En $-\infty$, il faut être plus attentif : le signe de $x^n$ dépend de la \cle{parité} de $n$. Si $n$ est \textbf{pair}, $x^n$ est positif (par exemple $(-100)^2 = \num{10000}$) ; si $n$ est \textbf{impair}, $x^n$ est négatif (par exemple $(-100)^3 = -\num{1000000}$).

\begin{propriete}[Limites de référence des fonctions puissances (admise)]
Pour tout entier naturel $n \geq 1$ :
\[ \lim_{x \to +\infty} x^n = +\infty. \]
En $-\infty$, le résultat dépend de la parité de $n$ :
\begin{itemize}
\item si $n$ est \textbf{pair} : $\displaystyle\lim_{x \to -\infty} x^n = +\infty$ ;
\item si $n$ est \textbf{impair} : $\displaystyle\lim_{x \to -\infty} x^n = -\infty$.
\end{itemize}
En particulier : $\displaystyle\lim_{x \to +\infty} x = +\infty$, $\displaystyle\lim_{x \to -\infty} x = -\infty$ et $\displaystyle\lim_{x \to -\infty} x^2 = +\infty$.
\end{propriete}

Cette propriété est \textbf{admise} : elle n'est pas à démontrer au Probatoire, mais elle doit être \emph{connue par cœur} et illustrée graphiquement.

\begin{popfigure}
\begin{center}
\begin{tikzpicture}
\begin{axis}[width=0.85\linewidth, height=7cm, axis lines=middle, xlabel={$x$}, ylabel={$y$}, xmin=-2.2, xmax=2.2, ymin=-4, ymax=6, xtick={-2,-1,1,2}, ytick={-4,-2,2,4,6}, legend style={at={(0.02,0.98)}, anchor=north west, font=\small}, samples=100]
\addplot[popcurve, popblue, domain=-2.2:2.2]{x};
\addlegendentry{$y=x$}
\addplot[popcurve, poporange, domain=-2.2:2.2]{x^2};
\addlegendentry{$y=x^2$}
\addplot[popcurve, popgreen, domain=-1.8:1.8]{x^3};
\addlegendentry{$y=x^3$}
\end{axis}
\end{tikzpicture}
\end{center}
\legende{Les courbes de $x \mapsto x$, $x \mapsto x^2$ et $x \mapsto x^3$. Toutes montent vers $+\infty$ quand $x \to +\infty$. En $-\infty$, $x^2$ (exposant pair) remonte vers $+\infty$, tandis que $x$ et $x^3$ (exposants impairs) descendent vers $-\infty$.}
\end{popfigure}

\begin{attention}[Parité de l'exposant]
L'erreur classique consiste à écrire $\lim_{x \to -\infty} x^2 = -\infty$ « parce que $x$ est négatif ». C'est faux : un carré est toujours positif ! Avant de conclure, regarde si l'exposant $n$ est \textbf{pair} ou \textbf{impair}.
\end{attention}

\courssub{Limites de la fonction inverse $x \mapsto \dfrac{1}{x}$}

La fonction inverse est définie sur $\mathbb{R}^* = ]-\infty\,; 0[ \cup ]0\,; +\infty[$. Son étude aux bornes de ce domaine fait apparaître une situation nouvelle : la limite en $0$ \emph{dépend du côté} par lequel on s'approche.

\textbf{En $+\infty$ et en $-\infty$.} Quand $x$ devient très grand, $\dfrac{1}{x}$ devient très petit : $\dfrac{1}{1000} = 0{,}001$, $\dfrac{1}{\num{1000000}} = 0{,}000001$. Le quotient s'approche de $0$ sans jamais l'atteindre. De même, quand $x$ tend vers $-\infty$, $\dfrac{1}{x}$ s'approche de $0$ par valeurs négatives.

\textbf{En $0$.} Quand $x$ s'approche de $0$ par valeurs \emph{positives} (on note $x \to 0^+$), le quotient $\dfrac{1}{x}$ explose vers le haut : $\dfrac{1}{0{,}001} = \num{1000}$, $\dfrac{1}{0{,}000001} = \num{1000000}$. En revanche, quand $x$ s'approche de $0$ par valeurs \emph{négatives} (on note $x \to 0^-$), le quotient est négatif et explose vers le bas : $\dfrac{1}{-0{,}001} = -\num{1000}$.

\begin{propriete}[Limites de référence de la fonction inverse (admise)]
\[ \lim_{x \to +\infty} \frac{1}{x} = 0 \qquad \text{et} \qquad \lim_{x \to -\infty} \frac{1}{x} = 0. \]
\[ \lim_{x \to 0^+} \frac{1}{x} = +\infty \qquad \text{et} \qquad \lim_{x \to 0^-} \frac{1}{x} = -\infty. \]
\end{propriete}

\begin{attention}[En $0$, toujours préciser le sens d'approche]
Il est \textbf{faux} d'écrire simplement « $\lim_{x \to 0} \frac{1}{x} = +\infty$ ». Le résultat dépend du côté : à droite de $0$ ($0^+$), la limite est $+\infty$ ; à gauche de $0$ ($0^-$), elle est $-\infty$. Comme les deux limites sont différentes, on dit que $\dfrac{1}{x}$ \textbf{n'a pas de limite en $0$}. Dans toute copie, précise toujours $0^+$ ou $0^-$.
\end{attention}

\begin{popfigure}
\begin{center}
\begin{tikzpicture}
\begin{axis}[width=0.85\linewidth, height=7.5cm, axis lines=middle, xlabel={$x$}, ylabel={$y$}, xmin=-5, xmax=5, ymin=-5, ymax=5, xtick={-4,-2,2,4}, ytick={-4,-2,2,4}, samples=200]
\addplot[popcurve, popblue, domain=0.2:5]{1/x};
\addplot[popcurve, popblue, domain=-5:-0.2]{1/x};
\addplot[dashed, poporange, thick] coordinates {(-5,0) (5,0)};
\addplot[dashed, poporange, thick] coordinates {(0,-5) (0,5)};
\node[poporange, font=\small] at (axis cs:3.6,0.5) {$y=0$};
\node[poporange, font=\small] at (axis cs:0.7,4.3) {$x=0$};
\end{axis}
\end{tikzpicture}
\end{center}
\legende{L'hyperbole $y = \dfrac{1}{x}$. La courbe se colle à l'axe des abscisses ($y=0$) quand $x \to \pm\infty$, et à l'axe des ordonnées ($x=0$) quand $x \to 0$.}
\end{popfigure}

\courssub{Limites de la fonction racine carrée}

La fonction $x \mapsto \sqrt{x}$ n'est définie que pour $x \geq 0$. On s'intéresse donc à son comportement en $0$ (la borne gauche de son domaine) et en $+\infty$.

\textbf{En $0$.} Quand $x$ s'approche de $0$ par valeurs positives, $\sqrt{x}$ s'approche de $0$ : $\sqrt{0{,}0001} = 0{,}01$. On a donc $\lim_{x \to 0^+} \sqrt{x} = 0 = \sqrt{0}$ : la fonction est « bien élevée » en $0$, sa limite égale son image.

\textbf{En $+\infty$.} Quand $x$ devient très grand, $\sqrt{x}$ devient très grand aussi, mais \emph{lentement} : $\sqrt{\num{10000}} = 100$, $\sqrt{\num{1000000}} = \num{1000}$.

\begin{propriete}[Limites de référence de la fonction racine carrée (admise)]
\[ \lim_{x \to 0^+} \sqrt{x} = 0 \qquad \text{et} \qquad \lim_{x \to +\infty} \sqrt{x} = +\infty. \]
\end{propriete}

\begin{aretenir}[Tableau récapitulatif des limites de référence]
\begin{center}
\begin{tabularx}{\linewidth}{|l|X|X|X|}
\hline
\rowcolor{popblueL} \textbf{Fonction} & \textbf{En $-\infty$} & \textbf{En $0$} & \textbf{En $+\infty$} \\
\hline
$x^n$ ($n$ pair) & $+\infty$ & $0$ & $+\infty$ \\
\hline
$x^n$ ($n$ impair) & $-\infty$ & $0$ & $+\infty$ \\
\hline
$\dfrac{1}{x}$ & $0$ & $+\infty$ en $0^+$ ; $-\infty$ en $0^-$ & $0$ \\
\hline
$\sqrt{x}$ & (non définie) & $0$ (en $0^+$) & $+\infty$ \\
\hline
\end{tabularx}
\end{center}
\end{aretenir}

\courssub{Limite en un point : le cas « facile » des fonctions usuelles}

Revenons sur l'observation faite pour $\sqrt{x}$ en $0$ : la limite était égale à l'image. C'est en fait le comportement \emph{normal} de toutes les fonctions usuelles en tout point où elles sont définies. Pour une fonction polynôme comme $P(x) = 2x^2 - 3x + 1$, rien de dramatique ne se produit quand $x$ s'approche de $4$ : $P(x)$ s'approche tout simplement de $P(4)$.

\begin{propriete}[Limite en un point des fonctions usuelles (admise)]
Si $f$ est une fonction polynôme, la fonction racine carrée, ou toute fonction construite à partir de celles-ci par sommes, produits et quotients, alors en tout point $a$ où $f$ est \textbf{définie} :
\[ \lim_{x \to a} f(x) = f(a). \]
\end{propriete}

\begin{exemplebox}[Calcul direct par substitution]
Calculer $\displaystyle\lim_{x \to 2} (3x^2 - 5x + 7)$.

La fonction $x \mapsto 3x^2 - 5x + 7$ est un polynôme, défini en $2$. On remplace donc simplement $x$ par $2$ :
\[ \lim_{x \to 2} (3x^2 - 5x + 7) = 3 \times 2^2 - 5 \times 2 + 7 = 12 - 10 + 7 = 9. \]
\end{exemplebox}

\begin{attention}[La substitution directe n'est possible que si $f$ est définie en $a$]
Pour $\displaystyle\lim_{x \to 1} \frac{2x+1}{x-1}$, on ne peut \textbf{pas} remplacer $x$ par $1$ : le dénominateur s'annule, $f(1)$ n'existe pas. C'est justement dans ces cas-là que les limites deviennent intéressantes (voir l'exemple résolu plus bas, puis la section 4).
\end{attention}

\courssub{Interprétation graphique : les asymptotes}

Voici maintenant le lien entre le calcul de limites et le dessin des courbes. Observons à nouveau l'hyperbole $y = \dfrac{1}{x}$ : quand $x \to +\infty$, la courbe se rapproche de la droite horizontale $y = 0$ sans jamais la toucher ; quand $x \to 0^+$, elle se rapproche de la droite verticale $x = 0$ en montant indéfiniment. Ces droites « magnétiques » s'appellent des \cle{asymptotes}.

\begin{definition}[Asymptote horizontale]
Soit $f$ une fonction et $\mathcal{C}_f$ sa courbe représentative. Si
\[ \lim_{x \to +\infty} f(x) = \ell \quad (\ell \text{ réel}) \qquad \text{ou} \qquad \lim_{x \to -\infty} f(x) = \ell, \]
alors la droite d'équation $y = \ell$ est appelée \cle{asymptote horizontale} à $\mathcal{C}_f$ en $+\infty$ (respectivement en $-\infty$).
\end{definition}

\begin{definition}[Asymptote verticale]
Soit $f$ une fonction et $a$ un réel. Si
\[ \lim_{x \to a^+} f(x) = \pm\infty \qquad \text{ou} \qquad \lim_{x \to a^-} f(x) = \pm\infty, \]
alors la droite d'équation $x = a$ est appelée \cle{asymptote verticale} à $\mathcal{C}_f$.
\end{definition}

Autrement dit, la lecture est bidirectionnelle et il faut savoir passer dans les deux sens :
\begin{itemize}
\item \textbf{Du calcul au graphique} : une limite finie $\ell$ en $\pm\infty$ donne une asymptote \emph{horizontale} $y = \ell$ ; une limite infinie en un point $a$ donne une asymptote \emph{verticale} $x = a$.
\item \textbf{Du graphique au calcul} : si l'on voit sur une courbe que celle-ci se colle à une droite horizontale $y = \ell$ aux extrémités, on en déduit $\lim_{x \to \pm\infty} f(x) = \ell$ ; si elle se colle à une droite verticale $x = a$, on en déduit $\lim_{x \to a} f(x) = \pm\infty$.
\end{itemize}

\begin{methode}[Trouver les asymptotes d'une fonction rationnelle simple]
Pour une fonction du type $f(x) = \dfrac{ax + b}{cx + d}$ avec $c \neq 0$ :
\begin{enumerate}
\item \textbf{Asymptote verticale} : chercher la valeur qui \textbf{annule le dénominateur}, c'est-à-dire résoudre $cx + d = 0$, ce qui donne $x = -\dfrac{d}{c}$. En cette valeur interdite, la limite de $f$ est infinie : la droite $x = -\dfrac{d}{c}$ est asymptote verticale.
\item \textbf{Asymptote horizontale} : pour $x$ très grand, les termes constants deviennent négligeables devant les termes en $x$, donc $f(x)$ se comporte comme $\dfrac{ax}{cx} = \dfrac{a}{c}$. La droite $y = \dfrac{a}{c}$ est asymptote horizontale.
\item \textbf{Vérifier} en calculant quelques valeurs (par exemple $f(100)$ et $f(-100)$) pour confirmer que $f(x)$ s'approche bien de $\dfrac{a}{c}$.
\end{enumerate}
\end{methode}

\begin{exoresolu}[Asymptotes de $f(x) = \dfrac{3x - 2}{x + 4}$]
Déterminer les asymptotes de la courbe de la fonction $f$ définie par $f(x) = \dfrac{3x-2}{x+4}$, et justifier.
\tcblower
\textbf{Asymptote verticale.} Le dénominateur s'annule pour $x + 4 = 0$, soit $x = -4$. La valeur $-4$ est donc interdite. Étudions la limite en $-4$ par valeurs supérieures : quand $x \to -4^+$, le numérateur tend vers $3 \times (-4) - 2 = -14$ (négatif) et le dénominateur $x + 4$ tend vers $0$ en restant \emph{positif}. Donc
\[ \lim_{x \to -4^+} f(x) = -\infty. \]
La droite d'équation $x = -4$ est asymptote verticale à la courbe.

\textbf{Asymptote horizontale.} Pour $x$ très grand, $f(x)$ se comporte comme $\dfrac{3x}{x} = 3$. Vérification numérique :
\[ f(100) = \frac{298}{104} \approx 2{,}87, \qquad f(\num{10000}) = \frac{\num{29998}}{\num{10004}} \approx 2{,}9988. \]
Les valeurs s'approchent bien de $3$ : $\displaystyle\lim_{x \to +\infty} f(x) = 3$ (et de même en $-\infty$). La droite d'équation $y = 3$ est asymptote horizontale.

\textbf{Conclusion.} La courbe de $f$ admet deux asymptotes : la droite verticale $x = -4$ et la droite horizontale $y = 3$.
\end{exoresolu}

\begin{popfigure}
\begin{center}
\begin{tikzpicture}
\begin{axis}[width=0.85\linewidth, height=8cm, axis lines=middle, xlabel={$x$}, ylabel={$y$}, xmin=-6, xmax=8, ymin=-6, ymax=9, xtick={-4,-2,2,4,6}, ytick={-4,-2,2,4,6,8}, samples=200, legend style={at={(0.98,0.02)}, anchor=south east, font=\small}]
\addplot[popcurve, popblue, domain=1.25:8]{(2*x+1)/(x-1)};
\addplot[popcurve, popblue, domain=-6:0.72]{(2*x+1)/(x-1)};
\addlegendentry{$y=\dfrac{2x+1}{x-1}$}
\addplot[dashed, poporange, thick] coordinates {(1,-6) (1,9)};
\addlegendentry{$x=1$}
\addplot[dashed, poppurple, thick] coordinates {(-6,2) (8,2)};
\addlegendentry{$y=2$}
\end{axis}
\end{tikzpicture}
\end{center}
\legende{Courbe de $f(x) = \dfrac{2x+1}{x-1}$. Le dénominateur s'annule en $x = 1$ : asymptote verticale $x = 1$. En $\pm\infty$, $f(x)$ se comporte comme $\dfrac{2x}{x} = 2$ : asymptote horizontale $y = 2$.}
\end{popfigure}

Sur cette figure, on lit directement les limites : quand $x \to 1^+$, la courbe monte vers $+\infty$ ; quand $x \to 1^-$, elle descend vers $-\infty$ ; et aux deux extrémités du graphique, elle se plaque contre la droite $y = 2$. C'est exactement le type de \emph{lecture graphique d'asymptotes} qui peut être demandée au Probatoire : on te donne une courbe, et tu dois en déduire les limites, ou inversement.

\begin{savaistu}[Des asymptotes dans l'atelier]
Les asymptotes ne sont pas qu'un jeu de mathématicien. En électricité, le courant de charge d'un condensateur décroît en s'approchant de $0$ : l'axe du temps est une asymptote horizontale de la courbe de charge. En mécanique, la vitesse d'un corps en chute avec frottements se rapproche d'une \emph{vitesse limite} sans jamais la dépasser : là encore, une asymptote horizontale. Savoir repérer une asymptote, c'est savoir prédire le comportement à long terme d'un système technique.
\end{savaistu}

\begin{exercice}
On considère la fonction $g$ définie par $g(x) = \dfrac{x + 5}{2x - 6}$.
\begin{enumerate}
\item Quel est l'ensemble de définition de $g$ ?
\item Déterminer les limites de $g$ en $3^+$ et en $3^-$. Qu'en déduit-on graphiquement ?
\item Déterminer les limites de $g$ en $+\infty$ et en $-\infty$ (on pourra calculer $g(100)$ et $g(-100)$ pour conjecturer). Qu'en déduit-on graphiquement ?
\end{enumerate}
\end{exercice}

\begin{corrige}[Exercice 1]
\begin{enumerate}
\item $g$ est définie si $2x - 6 \neq 0$, soit $x \neq 3$. Donc $D_g = \mathbb{R} \setminus \{3\}$.
\item Quand $x \to 3$, le numérateur tend vers $3 + 5 = 8 > 0$. En $3^+$, le dénominateur $2x - 6$ tend vers $0$ par valeurs positives, donc $\lim_{x \to 3^+} g(x) = +\infty$. En $3^-$, il tend vers $0$ par valeurs négatives, donc $\lim_{x \to 3^-} g(x) = -\infty$. La droite $x = 3$ est asymptote verticale.
\item $g(100) = \dfrac{105}{194} \approx 0{,}54$ et $g(-100) = \dfrac{-95}{-206} \approx 0{,}46$ : les valeurs s'approchent de $\dfrac{1}{2}$. En effet, pour $x$ très grand, $g(x)$ se comporte comme $\dfrac{x}{2x} = \dfrac{1}{2}$. Donc $\lim_{x \to \pm\infty} g(x) = \dfrac{1}{2}$ et la droite $y = \dfrac{1}{2}$ est asymptote horizontale.
\end{enumerate}
\end{corrige}

\begin{aretenir}[L'essentiel de la section]
\begin{itemize}
\item \textbf{Limites de référence} (à connaître par cœur) : $\lim_{x \to +\infty} x^n = +\infty$ ; en $-\infty$, $x^n$ tend vers $+\infty$ si $n$ est pair, vers $-\infty$ si $n$ est impair ; $\lim_{x \to \pm\infty} \dfrac{1}{x} = 0$ ; $\lim_{x \to 0^+} \dfrac{1}{x} = +\infty$, $\lim_{x \to 0^-} \dfrac{1}{x} = -\infty$ ; $\lim_{x \to +\infty} \sqrt{x} = +\infty$.
\item \textbf{En un point où $f$ est définie} : $\lim_{x \to a} f(x) = f(a)$ (substitution directe).
\item \textbf{Asymptote horizontale} $y = \ell$ : quand $\lim_{x \to \pm\infty} f(x) = \ell$.
\item \textbf{Asymptote verticale} $x = a$ : quand $\lim_{x \to a} f(x) = \pm\infty$ ; pour une fonction rationnelle, chercher les valeurs qui annulent le dénominateur.
\end{itemize}
\end{aretenir}

Ces limites de référence et ces interprétations graphiques sont nos outils de base. Mais que faire quand une fonction est une \emph{combinaison} de fonctions usuelles, comme $x^2 + \dfrac{1}{x}$ ou $\sqrt{x} - x$ ? C'est l'objet de la section suivante : les \emph{opérations sur les limites}.

\coursec{Opérations sur les limites}

Dans la section précédente, nous avons identifié les \textbf{limites de référence} (fonctions constantes, identité, puissances, inverses) et leur interprétation graphique. Nous savons maintenant « lire » la limite d'une fonction simple au voisinage d'une valeur interdite ou à l'infini. Mais la plupart des fonctions rencontrées en technique (électricité, mécanique, gestion) sont des \textbf{combinaisons} de fonctions élémentaires : sommes, produits, quotients, compositions. Cette section établit les règles qui permettent de calculer la limite d'une combinaison à partir des limites des fonctions de base.

\courssub{Limite d'une somme de fonctions}

L'intuition est simple : si $f(x)$ se rapproche de $L_1$ et $g(x)$ se rapproche de $L_2$, alors $f(x)+g(x)$ se rapproche de $L_1+L_2$. Cette règle reste valable lorsque l'une des limites est infinie, sous réserve de ne pas tomber dans le cas « $\infty - \infty$ ».

\begin{propriete}[Théorème : limite d'une somme (admis)]
Soient $f$ et $g$ deux fonctions définies sur un intervalle commun et $a \in \mathbb{R} \cup \{-\infty, +\infty\}$.
Si $\lim_{x \to a} f(x) = L_1$ et $\lim_{x \to a} g(x) = L_2$, alors :
\[
\lim_{x \to a} \bigl( f(x) + g(x) \bigr) = L_1 + L_2
\]
à condition que l'opération $L_1 + L_2$ ait un sens dans $\mathbb{R} \cup \{-\infty, +\infty\}$.
\end{propriete}

Le tableau récapitulatif ci-dessous synthétise tous les cas possibles pour la somme. Les cases \textbf{rouges} correspondent aux \textbf{\cle{formes indéterminées}} : le théorème ne permet pas de conclure, il faut transformer l'expression.

\begin{popfigure}
\legende{Tableau des limites d'une somme $f+g$ (FI = \cle{Forme Indéterminée})}
\begin{center}
\begin{tabular}{|c|ccc|}\hline
\rowcolor{popblueL}
$\lim f(x)$ \textbackslash $\lim g(x)$ & $L \in \mathbb{R}$ & $+\infty$ & $-\infty$ \\ \hline
$L' \in \mathbb{R}$ & $L+L'$ & $+\infty$ & $-\infty$ \\
$+\infty$ & $+\infty$ & $+\infty$ & \cellcolor{poppinkL}\textbf{FI : $\infty - \infty$} \\
$-\infty$ & $-\infty$ & \cellcolor{poppinkL}\textbf{FI : $\infty - \infty$} & $-\infty$ \\ \hline
\end{tabular}
\end{center}
\end{popfigure}

\begin{exemplebox}[Sommes sans indétermination]
Calculer les limites suivantes en $+\infty$ :
\begin{enumerate}
\item $\lim_{x \to +\infty} (x^2 + 3x)$
\item $\lim_{x \to +\infty} \left( \frac{1}{x} + 5 \right)$
\item $\lim_{x \to 0^+} \left( \frac{1}{x} + x^2 \right)$
\end{enumerate}
\textbf{Solution.}
\begin{enumerate}
\item $\lim x^2 = +\infty$, $\lim 3x = +\infty$. Somme : $+\infty + +\infty = +\infty$.
\item $\lim \frac{1}{x} = 0$, $\lim 5 = 5$. Somme : $0 + 5 = 5$.
\item $\lim_{x \to 0^+} \frac{1}{x} = +\infty$, $\lim_{x \to 0^+} x^2 = 0$. Somme : $+\infty + 0 = +\infty$.
\end{enumerate}
\end{exemplebox}

\courssub{Limite d'un produit de fonctions}

Si $f(x) \to L_1$ et $g(x) \to L_2$, le produit tend vers $L_1 \times L_2$. L'attention porte sur le signe (règle des signes) et sur le cas $0 \times \infty$.

\begin{propriete}[Théorème : limite d'un produit (admis)]
Soient $f$ et $g$ deux fonctions et $a \in \mathbb{R} \cup \{-\infty, +\infty\}$.
Si $\lim_{x \to a} f(x) = L_1$ et $\lim_{x \to a} g(x) = L_2$, alors :
\[
\lim_{x \to a} \bigl( f(x) \times g(x) \bigr) = L_1 \times L_2
\]
à condition que l'opération $L_1 \times L_2$ ait un sens (règle des signes pour les infinis).
\end{propriete}

\begin{popfigure}
\legende{Tableau des limites d'un produit $f \times g$}
\begin{center}
\begin{tabular}{|c|cccc|}\hline
\rowcolor{popblueL}
$\lim f$ \textbackslash $\lim g$ & $L > 0$ & $L < 0$ & $+\infty$ & $-\infty$ \\ \hline
$L' > 0$ & $+$ & $-$ & $+\infty$ & $-\infty$ \\
$L' < 0$ & $-$ & $+$ & $-\infty$ & $+\infty$ \\
$+\infty$ & $+\infty$ & $-\infty$ & $+\infty$ & $-\infty$ \\
$-\infty$ & $-\infty$ & $+\infty$ & $-\infty$ & $+\infty$ \\
$0$ & $0$ & $0$ & \cellcolor{poppinkL}\textbf{FI : $0 \times \infty$} & \cellcolor{poppinkL}\textbf{FI : $0 \times \infty$} \\ \hline
\end{tabular}
\end{center}
\end{popfigure}

\begin{attention}[Le piège du $0 \times \infty$]
On écrit souvent « $0 \times \infty = 0$ » par erreur. \textbf{Faux !} Si $f(x) \to 0$ et $g(x) \to \infty$, le produit $f(x)g(x)$ peut tendre vers $0$, vers une constante non nulle, vers $+\infty$ ou $-\infty$, ou ne pas avoir de limite.
\textbf{Exemple :} $f(x)=\frac{1}{x}$, $g(x)=x$ en $+\infty$ : $f \to 0$, $g \to +\infty$, mais $fg = 1 \to 1$.
Toujours transformer l'expression (factorisation, mise au même dénominateur) pour lever l'indétermination.
\end{attention}

\begin{exemplebox}[Produits sans indétermination]
\begin{enumerate}
\item $\lim_{x \to +\infty} (2x \times x^2) = +\infty \times +\infty = +\infty$.
\item $\lim_{x \to 0^+} (x \times \frac{1}{x^2}) = 0 \times +\infty$ \textbf{FI} (voir section 4).
\item $\lim_{x \to -\infty} (-3 \times x^3) = -3 \times (-\infty) = +\infty$.
\end{enumerate}
\end{exemplebox}

\courssub{Limite d'un quotient de fonctions}

C'est l'opération la plus délicate car elle combine les difficultés du produit (numérateur) et de l'inverse (dénominateur). Les formes indéterminées sont $\frac{0}{0}$ et $\frac{\infty}{\infty}$.

\begin{propriete}[Théorème : limite d'un quotient (admis)]
Soient $f$ et $g$ deux fonctions, $g(x) \neq 0$ au voisinage de $a$, et $a \in \mathbb{R} \cup \{-\infty, +\infty\}$.
Si $\lim_{x \to a} f(x) = L_1$ et $\lim_{x \to a} g(x) = L_2 \neq 0$, alors :
\[
\lim_{x \to a} \frac{f(x)}{g(x)} = \frac{L_1}{L_2}
\]
avec les conventions : $\frac{L}{\pm\infty} = 0$ ($L$ fini), $\frac{\pm\infty}{L} = \pm\infty$ ($L > 0$ fini), règle des signes pour les infinis.
\end{propriete}

\begin{popfigure}
\legende{Tableau des limites d'un quotient $\frac{f}{g}$}
\begin{center}
\small
\begin{tabular}{|c|ccccc|}\hline
\rowcolor{popblueL}
$\lim f$ \textbackslash $\lim g$ & $L > 0$ & $L < 0$ & $+\infty$ & $-\infty$ & $0$ \\ \hline
$L' > 0$ & $+$ & $-$ & $0$ & $0$ & $+\infty$ \\
$L' < 0$ & $-$ & $+$ & $0$ & $0$ & $-\infty$ \\
$+\infty$ & $+\infty$ & $-\infty$ & \cellcolor{poppinkL}\textbf{FI : $\frac{\infty}{\infty}$} & \cellcolor{poppinkL}\textbf{FI : $\frac{\infty}{\infty}$} & $+\infty$ \\
$-\infty$ & $-\infty$ & $+\infty$ & \cellcolor{poppinkL}\textbf{FI : $\frac{\infty}{\infty}$} & \cellcolor{poppinkL}\textbf{FI : $\frac{\infty}{\infty}$} & $-\infty$ \\
$0$ & $0$ & $0$ & $0$ & $0$ & \cellcolor{poppinkL}\textbf{FI : $\frac{0}{0}$} \\ \hline
\end{tabular}
\end{center}
\end{popfigure}

\begin{exemplebox}[Quotients sans indétermination]
\begin{enumerate}
\item $\lim_{x \to 2} \frac{x^2+1}{x-1} = \frac{5}{1} = 5$.
\item $\lim_{x \to +\infty} \frac{5}{x} = \frac{5}{+\infty} = 0$.
\item $\lim_{x \to 0^+} \frac{x^2}{x} = \frac{0}{0}$ \textbf{FI} (simplifier $x$ donne $\lim x = 0$).
\item $\lim_{x \to +\infty} \frac{x^2+3}{2x} = \frac{+\infty}{+\infty}$ \textbf{FI} (factoriser $x$).
\end{enumerate}
\end{exemplebox}

\courssub{Les formes indéterminées : synthèse}

\begin{definition}[Forme indéterminée (FI)]
On appelle \textbf{\cle{forme indéterminée}} une situation où l'application directe des théorèmes sur les opérations (somme, produit, quotient) aboutit à une expression symbolique qui \textbf{ne permet pas de déterminer la limite} (ni sa valeur finie, ni son signe infini, ni même son existence).
\end{definition}

Il existe \textbf{quatre formes indéterminées fondamentales} en Première technique :
\begin{center}
\begin{tabular}{|c|c|c|}\hline
\rowcolor{popblueL}
\textbf{Opération} & \textbf{Forme Indéterminée} & \textbf{Action à mener} \\ \hline
Somme / Différence & $\infty - \infty$ (ou $+\infty + (-\infty)$) & \cle{Factoriser}, \cle{rationaliser}, mettre au même dénominateur \\ \hline
Produit & $0 \times \infty$ & Réécrire en quotient ($\frac{0}{1/\infty}$ ou $\frac{\infty}{1/0}$) \\ \hline
Quotient & $\frac{0}{0}$ & Simplifier, factoriser, rationaliser \\ \hline
Quotient & $\frac{\infty}{\infty}$ & \cle{Factoriser par le terme de plus haut degré} \\ \hline
\end{tabular}
\end{center}

\begin{aretenir}[Règle d'or]
\textbf{Une forme indéterminée n'est PAS un résultat.} C'est un signal d'alerte : « \emph{La méthode directe échoue, il faut transformer l'expression algébriquement avant de recalculer la limite.} »
\end{aretenir}

\courssub{Limite d'une fonction composée (approche intuitive)}

Si $u(x) \to \ell$ et que la fonction $f$ admet une limite en $\ell$ (ce qui est le cas si $f$ est \textbf{continue} en $\ell$ pour $\ell$ fini, ou si $f$ a une limite en $\pm\infty$), alors $f(u(x))$ tend vers cette limite. En Première technique, on utilise cette propriété de façon intuitive pour les fonctions usuelles continues ($\sqrt{\cdot}$, $\exp$, $\ln$, $\sin$, $\cos$, polynômes).

\begin{propriete}[Limite d'une composée — Règle pratique (admise)]
Soit $u$ une fonction telle que $\lim_{x \to a} u(x) = \ell \in \mathbb{R} \cup \{-\infty, +\infty\}$.
Soit $f$ une fonction telle que $\lim_{y \to \ell} f(y)$ existe (notamment si $f$ est continue en $\ell$ pour $\ell$ fini).
Alors :
\[
\lim_{x \to a} f(u(x)) = \lim_{y \to \ell} f(y)
\]
\textit{En pratique : on « fait entrer la limite à l'intérieur de la fonction » si celle-ci admet une limite en $\ell$.}
\end{propriete}

\begin{exemplebox}[Limites de composées]
\begin{enumerate}
\item $\lim_{x \to +\infty} \sqrt{2x+3}$.
$\lim_{x \to +\infty} (2x+3) = +\infty$. La fonction racine carrée est continue sur $[0, +\infty[$ et $\lim_{X \to +\infty} \sqrt{X} = +\infty$.
Donc $\lim_{x \to +\infty} \sqrt{2x+3} = +\infty$.

\item $\lim_{x \to 0} \sin(x^2+1)$.
$\lim_{x \to 0} (x^2+1) = 1$. $\sin$ est continue sur $\mathbb{R}$, $\sin(1)$ existe.
Donc la limite est $\sin(1)$.

\item $\lim_{x \to +\infty} e^{-x}$.
$\lim_{x \to +\infty} (-x) = -\infty$. $\exp$ est continue, $\lim_{X \to -\infty} e^X = 0$.
Donc la limite est $0$.
\end{enumerate}
\end{exemplebox}

\courssub{Méthode générale de calcul de limite}

\begin{methode}[Calculer une limite en 3 étapes]
\textbf{Étape 1 : Identifier la forme.} Remplacer $x$ par la valeur visée ($a$, $+\infty$, $-\infty$) dans l'expression \emph{sans calculer}, juste pour voir les limites des parties élémentaires (constantes, $x$, $x^n$, $1/x$, etc.).

\textbf{Étape 2 : Appliquer les théorèmes (opérations).} Combiner ces limites à l'aide des tableaux (somme, produit, quotient, composée).
\begin{itemize}
\item Si le résultat est un nombre fini, $+\infty$ ou $-\infty$ \textbf{bien défini} $\to$ \textbf{C'est la limite finale.}
\item Si le résultat est une \textbf{\cle{Forme Indéterminée}} ($\infty-\infty$, $0\times\infty$, $\frac{0}{0}$, $\frac{\infty}{\infty}$) $\to$ Passer à l'étape 3.
\end{itemize}

\textbf{Étape 3 : Lever l'indétermination (transformation algébrique).}
\begin{itemize}
\item $\frac{0}{0}$ ou $\frac{\infty}{\infty}$ (fractions rationnelles) : \textbf{Factoriser par le terme de plus haut degré} (numérateur et dénominateur) ou \textbf{simplifier} un facteur commun.
\item $\infty - \infty$ (somme de fractions ou racine) : \textbf{Mettre au même dénominateur} ou \textbf{multiplier par le conjugué} (si racine).
\item $0 \times \infty$ : \textbf{Réécrire en quotient} pour retomber sur $\frac{0}{0}$ ou $\frac{\infty}{\infty}$.
\end{itemize}
Puis \textbf{revenir à l'Étape 2} avec l'expression transformée.
\end{methode}

\begin{exoresolu}[Calcul complet avec levée d'indétermination]
Calculer $\lim_{x \to +\infty} \frac{3x^2 - 5x + 2}{2x^2 + x - 1}$.
\tcblower
\textbf{Solution détaillée.}

\textbf{Étape 1 : Forme.}
Numérateur : $3x^2 \to +\infty$, $-5x \to -\infty$, $2 \to 2$ $\to$ FI $\infty - \infty$ globalement $+\infty$.
Dénominateur : $2x^2 \to +\infty$, $x \to +\infty$, $-1 \to -1$ $\to$ $+\infty$.
Quotient : $\frac{+\infty}{+\infty} = \textbf{FI de type $\frac{\infty}{\infty}$}$.

\textbf{Étape 2 : Théorèmes directs $\to$ FI. On ne peut pas conclure.}

\textbf{Étape 3 : Lever la FI.}
C'est une fraction rationnelle. Le plus haut degré est $2$ (numérateur et dénominateur).
On factorise $x^2$ au numérateur ET au dénominateur :
\[
\frac{3x^2 - 5x + 2}{2x^2 + x - 1}
= \frac{x^2 \left( 3 - \frac{5}{x} + \frac{2}{x^2} \right)}{x^2 \left( 2 + \frac{1}{x} - \frac{1}{x^2} \right)}
= \frac{3 - \frac{5}{x} + \frac{2}{x^2}}{2 + \frac{1}{x} - \frac{1}{x^2}} \quad (x \neq 0)
\]

\textbf{Étape 2 (bis) : Recalculer sur la forme transformée.}
$\lim_{x \to +\infty} \frac{5}{x} = 0$, $\lim \frac{2}{x^2} = 0$, $\lim \frac{1}{x} = 0$, $\lim \frac{1}{x^2} = 0$.
Numérateur $\to 3 - 0 + 0 = 3$.
Dénominateur $\to 2 + 0 - 0 = 2$.
Quotient $\to \frac{3}{2}$ (forme définie).

\textbf{Conclusion :}
\[
\boxed{\lim_{x \to +\infty} \frac{3x^2 - 5x + 2}{2x^2 + x - 1} = \frac{3}{2}}
\]
\end{exoresolu}

\begin{savaistu}[Application technique : Résistance équivalente]
En électricité, deux résistances $R_1(x)$ et $R_2(x)$ variables (thermistances) sont montées en parallèle. La résistance équivalente est $R_{eq}(x) = \frac{R_1(x) R_2(x)}{R_1(x) + R_2(x)}$.
Si $R_1(x) \to 0$ (court-circuit) et $R_2(x) \to 10\,\Omega$, alors $R_{eq} \to \frac{0 \times 10}{0 + 10} = 0$.
Si $R_1(x) \to +\infty$ (circuit ouvert) et $R_2(x) \to 10\,\Omega$, alors $R_{eq} \to \frac{+\infty \times 10}{+\infty + 10} = \frac{\infty}{\infty}$ (FI). En factorisant $R_1$ : $\frac{R_2}{1 + R_2/R_1} \to \frac{10}{1+0} = 10\,\Omega$. Le circuit se comporte comme la seule résistance $R_2$.
\end{savaistu}

\coursec{Calculs pratiques : lever les formes indéterminées}

Dans la section précédente, nous avons établi les règles opératoires sur les limites : somme, produit, quotient. Mais nous avons rencontré des situations bloquantes, les fameuses \cle{formes indéterminées} : $\infty - \infty$, $0 \times \infty$, $\dfrac{\infty}{\infty}$ et $\dfrac{0}{0}$. Face à elles, aucune règle directe ne s'applique. Cette section est la boîte à outils du technicien : elle donne des \cle{techniques} systématiques pour \emph{lever} chaque indétermination, c'est-à-dire transformer l'écriture de la fonction jusqu'à ce que le calcul devienne possible. Retenez l'idée directrice : \textbf{une forme indéterminée ne signifie pas « pas de limite », elle signifie « il faut travailler l'expression »}.

\courssub{Limite en $\pm\infty$ d'un polynôme : la règle du terme de plus haut degré}

\textbf{Intuition.} Prenons le polynôme $P(x) = 3x^2 - 50x + 7$. Quand $x$ devient immense (par exemple $x = 10^6$), le terme $3x^2$ vaut $3 \times 10^{12}$, alors que $-50x$ ne vaut « que » $-5 \times 10^7$ et que $7$ est négligeable. Le terme de plus haut degré \emph{écrase} tous les autres : c'est lui qui commande le comportement de $P$ à l'infini.

\begin{propriete}[Limite d'un polynôme en $\pm\infty$ — démonstration exigible]
En $+\infty$ et en $-\infty$, un polynôme a la même limite que son \cle{terme de plus haut degré}.
\end{propriete}

\begin{experience}[Démonstration par factorisation]
Soit $P(x) = a_n x^n + a_{n-1}x^{n-1} + \dots + a_1 x + a_0$ avec $a_n \neq 0$. L'idée est de \textbf{factoriser par le terme dominant} $a_n x^n$ :
\[ P(x) = a_n x^n \left( 1 + \frac{a_{n-1}}{a_n}\cdot\frac{1}{x} + \dots + \frac{a_1}{a_n}\cdot\frac{1}{x^{n-1}} + \frac{a_0}{a_n}\cdot\frac{1}{x^n} \right). \]
Or, quand $x \to \pm\infty$, chaque terme $\dfrac{1}{x}$, $\dfrac{1}{x^2}$, \dots, $\dfrac{1}{x^n}$ tend vers $0$. La parenthèse tend donc vers $1 + 0 + \dots + 0 = 1$. Par produit des limites, $P(x)$ se comporte comme $a_n x^n \times 1 = a_n x^n$. \hfill$\square$
\end{experience}

\begin{exemplebox}[Application directe]
Déterminons $\displaystyle\lim_{x \to +\infty} \left( -2x^3 + 7x^2 - x + 12 \right)$.

Le terme de plus haut degré est $-2x^3$. Donc :
\[ \lim_{x \to +\infty} \left( -2x^3 + 7x^2 - x + 12 \right) = \lim_{x \to +\infty} (-2x^3) = -\infty, \]
car $x^3 \to +\infty$ et le coefficient $-2$ est négatif.
\end{exemplebox}

\begin{attention}[Ne pas additionner les limites terme à terme]
Écrire « $-\infty + \infty - \infty + 12$ » pour $-2x^3 + 7x^2 - x + 12$ conduit à la forme indéterminée $\infty - \infty$ et ne donne \textbf{aucune} conclusion. La règle du terme de plus haut degré est précisément l'outil qui lève cette indétermination. Citez-la explicitement dans votre rédaction.
\end{attention}

\courssub{Limite en $\pm\infty$ d'une fonction rationnelle}

\textbf{Intuition.} Pour une fraction rationnelle $f(x) = \dfrac{P(x)}{Q(x)}$, numérateur et dénominateur tendent tous deux vers l'infini : c'est la forme $\dfrac{\infty}{\infty}$. Mais chacun est dominé par son propre terme de plus haut degré. Tout se joue donc entre ces deux termes dominants.

\begin{propriete}[Limite d'une fraction rationnelle en $\pm\infty$]
Soit $f(x) = \dfrac{P(x)}{Q(x)}$ une fonction rationnelle, avec $P$ et $Q$ des polynômes, $Q$ non nul. En $+\infty$ et en $-\infty$, $f$ a la même limite que le \cle{quotient de ses termes de plus haut degré} :
\[ \lim_{x \to \pm\infty} \frac{a_n x^n + \dots}{b_p x^p + \dots} = \lim_{x \to \pm\infty} \frac{a_n x^n}{b_p x^p}. \]
\end{propriete}

\begin{experience}[Démonstration guidée]
On factorise en haut et en bas par les termes dominants :
\[ \frac{a_n x^n + a_{n-1}x^{n-1} + \dots}{b_p x^p + b_{p-1}x^{p-1} + \dots} = \frac{a_n x^n}{b_p x^p} \times \frac{1 + \dfrac{a_{n-1}}{a_n x} + \dots}{1 + \dfrac{b_{p-1}}{b_p x} + \dots}. \]
Quand $x \to \pm\infty$, la grande fraction de droite tend vers $\dfrac{1}{1} = 1$ (tous les termes en $\dfrac{1}{x}$ tendent vers $0$). Il reste le quotient des termes dominants. \hfill$\square$
\end{experience}

\textbf{Les trois cas possibles.} Après simplification de $\dfrac{a_n x^n}{b_p x^p}$, trois situations se présentent :

\begin{center}
\begin{tabularx}{\linewidth}{|l|X|X|X|}
\hline
\rowcolor{popblueL} \textbf{Comparaison des degrés} & \textbf{Forme après simplification} & \textbf{Limite en $\pm\infty$} & \textbf{Exemple} \\
\hline
$\deg P > \deg Q$ & $\dfrac{a_n}{b_p}\, x^{n-p}$ avec $n-p \geq 1$ & $\pm\infty$ (signe selon coefficients et parité de $n-p$) & $\dfrac{2x^3}{x^2} = 2x \to +\infty$ \\
\hline
$\deg P = \deg Q$ & $\dfrac{a_n}{b_p}$ (constante) & $\dfrac{a_n}{b_p}$ & $\dfrac{3x^2}{x^2} = 3$ \\
\hline
$\deg P < \deg Q$ & $\dfrac{a_n}{b_p}\cdot\dfrac{1}{x^{p-n}}$ & $0$ & $\dfrac{x}{x^3} = \dfrac{1}{x^2} \to 0$ \\
\hline
\end{tabularx}
\end{center}

\begin{attention}[Signe de l'infini en $-\infty$ pour $\deg P > \deg Q$]
Lorsque la limite est $\pm\infty$ (cas $\deg P > \deg Q$), le signe en $-\infty$ dépend de la \textbf{parité de la différence des degrés} $n-p$ et du signe du quotient des coefficients $\frac{a_n}{b_p}$.
\begin{itemize}
\item Si $n-p$ est \textbf{pair} : $x^{n-p} > 0$ pour $x \to -\infty$, le signe est celui de $\frac{a_n}{b_p}$.
\item Si $n-p$ est \textbf{impair} : $x^{n-p} \to -\infty$ pour $x \to -\infty$, le signe est l'opposé de celui de $\frac{a_n}{b_p}$.
\end{itemize}
Exemple : $\dfrac{x^3}{x^2} = x \to -\infty$ en $-\infty$ (différence $1$, impaire, coeff $+$).
Exemple : $\dfrac{x^4}{x^2} = x^2 \to +\infty$ en $-\infty$ (différence $2$, paire, coeff $+$).
\end{attention}

\begin{exoresolu}[Les trois degrés en action]
\textbf{Énoncé.} Calculer les limites en $+\infty$ de :
\quad (a) $f(x) = \dfrac{3x^2 - x}{x^2 + 5}$ \qquad (b) $g(x) = \dfrac{x^3 + 1}{2x - 4}$ \qquad (c) $h(x) = \dfrac{4x + 1}{x^2 - 7}$.
\tcblower
\textbf{Solution détaillée.}

(a) Degrés égaux ($2 = 2$). Quotient des termes dominants : $\dfrac{3x^2}{x^2} = 3$.
\[ \lim_{x \to +\infty} f(x) = \lim_{x \to +\infty} \frac{3x^2}{x^2} = 3. \]

(b) Degré du numérateur supérieur ($3 > 1$). Quotient des termes dominants : $\dfrac{x^3}{2x} = \dfrac{x^2}{2}$.
\[ \lim_{x \to +\infty} g(x) = \lim_{x \to +\infty} \frac{x^2}{2} = +\infty. \]

(c) Degré du numérateur inférieur ($1 < 2$). Quotient des termes dominants : $\dfrac{4x}{x^2} = \dfrac{4}{x}$.
\[ \lim_{x \to +\infty} h(x) = \lim_{x \to +\infty} \frac{4}{x} = 0. \]
\end{exoresolu}

\begin{exemplebox}[Cas $\deg P > \deg Q$ en $-\infty$]
Calculons $\displaystyle\lim_{x \to -\infty} \frac{2x^3 - x}{x^2 + 1}$.
Degrés : $3 > 1$, différence $2$ (paire). Quotient des termes dominants : $\dfrac{2x^3}{x^2} = 2x$.
En $-\infty$, $2x \to -\infty$. Donc la limite est $-\infty$.
\end{exemplebox}

\courssub{Limite en un point $a$ : lever la forme $\dfrac{0}{0}$ par factorisation}

\textbf{Intuition.} Quand, en remplaçant $x$ par $a$ dans une fraction, on obtient $\dfrac{0}{0}$, cela signifie que $a$ est \emph{racine} du numérateur \emph{et} du dénominateur. Or un polynôme qui s'annule en $a$ se factorise par $(x - a)$. On peut donc simplifier le facteur commun $(x-a)$ en haut et en bas : c'est lui qui provoquait l'indétermination.

\begin{methode}[Lever $\dfrac{0}{0}$ en $x = a$ sur une fraction rationnelle]
\begin{enumerate}
\item Vérifier que le remplacement direct donne $\dfrac{0}{0}$ (forme indéterminée).
\item Factoriser le numérateur par $(x-a)$ : identités remarquables, racine évidente, ou division euclidienne (Ruffini).
\item Factoriser de même le dénominateur par $(x-a)$.
\item Simplifier le facteur $(x-a)$, en précisant que $x \neq a$ (on cherche une limite, pas la valeur en $a$).
\item Remplacer $x$ par $a$ dans l'expression simplifiée et conclure.
\end{enumerate}
\end{methode}

\begin{exemplebox}[Le classique du cours]
Calculons $\displaystyle\lim_{x \to 2} \frac{x^2 - 4}{x - 2}$.

Remplacement direct : $\dfrac{2^2 - 4}{2 - 2} = \dfrac{0}{0}$ : forme indéterminée.

On factorise le numérateur grâce à l'identité $a^2 - b^2 = (a-b)(a+b)$ : $x^2 - 4 = (x-2)(x+2)$. Pour $x \neq 2$ :
\[ \frac{x^2 - 4}{x - 2} = \frac{(x-2)(x+2)}{x-2} = x + 2. \]
Donc $\displaystyle\lim_{x \to 2} \frac{x^2 - 4}{x - 2} = \lim_{x \to 2} (x+2) = 4$.
\end{exemplebox}

\begin{attention}[Pourquoi a-t-on le droit de simplifier ?]
La fonction $x \mapsto \dfrac{x^2-4}{x-2}$ n'est \textbf{pas définie} en $x = 2$. Mais une limite en $2$ décrit le comportement de $f(x)$ quand $x$ \emph{s'approche} de $2$ sans jamais l'atteindre : on a donc toujours $x \neq 2$, ce qui autorise la simplification par $(x-2)$. Rédigez ainsi : « pour $x \neq 2$, on a \dots ». Ne dites jamais « $f(2) = 4$ » : la fonction n'a pas de valeur en $2$, elle a une \emph{limite} égale à $4$.
\end{attention}

\courssub{La technique de la quantité conjuguée (expressions avec racines carrées)}

\textbf{Intuition.} Quand la forme $\dfrac{0}{0}$ provient d'une racine carrée, on ne peut pas factoriser directement. L'astuce consiste à multiplier en haut et en bas par la \cle{quantité conjuguée}, afin de faire apparaître l'identité $(a-b)(a+b) = a^2 - b^2$, qui \emph{élimine les racines} au numérateur.

\begin{methode}[Multiplier par la quantité conjuguée]
Pour une expression contenant une différence avec racine carrée (ex : $\sqrt{u} - v$ ou $\sqrt{u} - \sqrt{v}$) :
\begin{enumerate}
\item Multiplier numérateur et dénominateur par la conjuguée ($\sqrt{u} + v$ ou $\sqrt{u} + \sqrt{v}$).
\item Développer le numérateur : $(\sqrt{u})^2 - v^2 = u - v^2$ (ou $u - v$) : les racines disparaissent.
\item Factoriser et simplifier le facteur qui s'annule (souvent $(x-a)$).
\item Remplacer et conclure.
\end{enumerate}
\end{methode}

\begin{exoresolu}[Racine carrée et quantité conjuguée]
\textbf{Énoncé.} Calculer $\displaystyle\lim_{x \to 3} \frac{\sqrt{x+1} - 2}{x - 3}$.
\tcblower
\textbf{Solution détaillée.} Remplacement direct : $\dfrac{\sqrt{4} - 2}{3 - 3} = \dfrac{0}{0}$ : forme indéterminée.

On multiplie par la quantité conjuguée $\sqrt{x+1} + 2$ :
\begin{align*}
\frac{\sqrt{x+1} - 2}{x - 3} &= \frac{(\sqrt{x+1} - 2)(\sqrt{x+1} + 2)}{(x - 3)(\sqrt{x+1} + 2)} \\
&= \frac{(x+1) - 4}{(x - 3)(\sqrt{x+1} + 2)} \\
&= \frac{x - 3}{(x - 3)(\sqrt{x+1} + 2)} \\
&= \frac{1}{\sqrt{x+1} + 2} \quad \text{(car } x \neq 3\text{)}.
\end{align*}
Donc :
\[ \lim_{x \to 3} \frac{\sqrt{x+1} - 2}{x - 3} = \lim_{x \to 3} \frac{1}{\sqrt{x+1} + 2} = \frac{1}{\sqrt{4} + 2} = \frac{1}{4}. \]
\end{exoresolu}

\courssub{Quelle méthode choisir ? L'arbre de décision}

Face à une forme indéterminée, le bon réflexe est de \emph{reconnaître la situation} avant de calculer. L'arbre suivant résume la stratégie à adopter.

\begin{popfigure}
\begin{center}
\begin{tikzpicture}[
  font=\small,
  decision/.style={rectangle, rounded corners, draw=popdark, fill=popblueL, text width=4.2cm, align=center, minimum height=1cm},
  method/.style={rectangle, rounded corners, draw=popgreen!60!black, fill=popgreenL, text width=4.6cm, align=center, minimum height=0.9cm},
  arr/.style={-{Stealth[length=2.5mm]}, thick, popdark}
]
\node[decision] (racine) at (0,0) {\textbf{Forme indéterminée} : quel type d'expression ?};
\node[decision, align=center] (poly) at (-6.5,-2.2){Polynôme en $\pm\infty$\\ ($\infty - \infty$)};
\node[decision, align=center] (rat) at (-2.2,-2.2){Fraction rationnelle en $\pm\infty$\\ ($\frac{\infty}{\infty}$)};
\node[decision, align=center] (rat0) at (2.2,-2.2){Fraction rationnelle en $a$\\ ($\frac{0}{0}$)};
\node[decision, align=center] (sqrt) at (6.5,-2.2){Racines carrées en $a$\\ ($\frac{0}{0}$)};
\node[method] (mpoly) at (-6.5,-4.4) {Terme de plus haut degré};
\node[method] (mrat) at (-2.2,-4.4) {Quotient des termes de plus haut degré};
\node[method] (mrat0) at (2.2,-4.4) {Factoriser par $(x-a)$, simplifier};
\node[method] (msqrt) at (6.5,-4.4) {Multiplier par la quantité conjuguée};
\draw[arr] (racine) -- (poly);
\draw[arr] (racine) -- (rat);
\draw[arr] (racine) -- (rat0);
\draw[arr] (racine) -- (sqrt);
\draw[arr] (poly) -- (mpoly);
\draw[arr] (rat) -- (mrat);
\draw[arr] (rat0) -- (mrat0);
\draw[arr] (sqrt) -- (msqrt);
\end{tikzpicture}
\end{center}
\legende{Arbre de décision : la méthode dépend du type d'expression et du point où l'on cherche la limite.}
\end{popfigure}

\begin{aretenir}[Boîte à outils : résumé des techniques de levée d'indétermination]
\begin{itemize}
\item \textbf{Polynôme en $\pm\infty$} : garder le terme de plus haut degré.
\item \textbf{Fraction rationnelle en $\pm\infty$} : comparer les degrés (quotient des termes dominants).
\item \textbf{Fraction rationnelle en $a$ (forme $0/0$)} : factoriser par $(x-a)$ (identités, Ruffini), simplifier.
\item \textbf{Expression avec racine(s) en $a$ (forme $0/0$)} : multiplier par la quantité conjuguée, simplifier.
\end{itemize}
\end{aretenir}

\courssub{Entraînement à la rédaction complète}

Au Probatoire comme au Baccalauréat technique, un calcul de limite doit être \textbf{rédigé et justifié} : annoncer la forme obtenue par remplacement, nommer la technique utilisée, détailler la transformation, conclure. Voici un modèle complet.

\begin{exoresolu}[Rédaction modèle]
\textbf{Énoncé.} Calculer $\displaystyle\lim_{x \to -1} \frac{x^2 + 3x + 2}{x^2 - 1}$.
\tcblower
\textbf{Solution rédigée.}

\emph{Étape 1 — Recherche de la forme.} En remplaçant $x$ par $-1$ : numérateur $= 1 - 3 + 2 = 0$ ; dénominateur $= 1 - 1 = 0$. On obtient la forme indéterminée $\dfrac{0}{0}$.

\emph{Étape 2 — Factorisation.} $-1$ est racine du numérateur ; l'autre racine est $-2$ (produit $2$, somme $-3$) : $x^2 + 3x + 2 = (x+1)(x+2)$. De plus $x^2 - 1 = (x-1)(x+1)$.

\emph{Étape 3 — Simplification.} Pour $x \neq -1$ :
\[ \frac{x^2 + 3x + 2}{x^2 - 1} = \frac{(x+1)(x+2)}{(x-1)(x+1)} = \frac{x+2}{x-1}. \]

\emph{Étape 4 — Conclusion.}
\[ \lim_{x \to -1} \frac{x^2 + 3x + 2}{x^2 - 1} = \lim_{x \to -1} \frac{x+2}{x-1} = \frac{-1+2}{-1-1} = -\frac{1}{2}. \]
\end{exoresolu}

\begin{exercice}[À vous de jouer]
Calculer, en rédigeant entièrement la démarche :
\begin{enumerate}
\item $\displaystyle\lim_{x \to -\infty} \left( 5x^3 - 2x^2 + x - 9 \right)$ ;
\item $\displaystyle\lim_{x \to +\infty} \frac{2x^2 - 3x + 1}{x^2 + x}$ ;
\item $\displaystyle\lim_{x \to 5} \frac{x^2 - 25}{x - 5}$ ;
\item $\displaystyle\lim_{x \to 0} \frac{\sqrt{x+4} - 2}{x}$ ;
\item $\displaystyle\lim_{x \to -\infty} \frac{x^3 + 2x}{x^2 - 1}$.
\end{enumerate}
\end{exercice}

\begin{corrige}[Exercice « À vous de jouer »]
\begin{enumerate}
\item Terme de plus haut degré : $5x^3$. $\displaystyle\lim_{x \to -\infty} 5x^3 = -\infty$ (cube négatif, coefficient positif).
\item Degrés égaux ($2=2$) : $\displaystyle\lim_{x \to +\infty} \frac{2x^2}{x^2} = 2$.
\item $\dfrac{x^2-25}{x-5} = \dfrac{(x-5)(x+5)}{x-5} = x+5$ pour $x \neq 5$, donc la limite vaut $10$.
\item Conjuguée : $\dfrac{(\sqrt{x+4}-2)(\sqrt{x+4}+2)}{x(\sqrt{x+4}+2)} = \dfrac{x}{x(\sqrt{x+4}+2)} = \dfrac{1}{\sqrt{x+4}+2}$, donc la limite vaut $\dfrac{1}{4}$.
\item Degrés : num $3$, dén $2$ ($3>2$). Quotient dominants : $\dfrac{x^3}{x^2} = x$. Différence des degrés $1$ (impair). En $-\infty$, $x \to -\infty$. Limite $= -\infty$.
\end{enumerate}
\end{corrige}

\begin{savaistu}[Vers la dérivée]
Ces techniques de levée d'indétermination sont bien plus qu'un exercice de calcul : elles sont le \textbf{premier pas vers la notion de dérivée}. En effet, le taux d'accroissement $\dfrac{f(x) - f(a)}{x - a}$ donne toujours la forme $\dfrac{0}{0}$ quand $x$ tend vers $a$ : c'est exactement la situation que vous venez d'apprendre à traiter ! La dérivée, étudiée dans la prochaine leçon et au cœur du programme de Terminale technique, permettra de modéliser vitesses, variations de coûts de production en atelier et intensités variables en électricité.
\end{savaistu}

\coursec{Limites et continuité}

\courssub{Approche intuitive de la notion de limite}

L'étude des limites permet de décrire le comportement d'une fonction \emph{au voisinage} d'une valeur (finie ou infinie) sans nécessairement calculer l'image de cette valeur. C'est l'outil indispensable pour modéliser des phénomènes qui se stabilisent, qui divergent ou qui présentent une rupture.

\begin{experience}[Mise en situation : la charge d'un condensateur]
Dans un circuit RC série alimenté par une tension continue $E$, la tension $u_C(t)$ aux bornes du condensateur pendant la charge suit la loi :
\[
u_C(t) = E \left(1 - e^{-t/\tau}\right), \quad \tau = RC.
\]
\begin{enumerate}
    \item Au démarrage ($t=0$), $u_C(0) = 0$.
    \item Lorsque le temps $t$ devient \textbf{très grand} ($t \to +\infty$), le terme $e^{-t/\tau}$ devient \textbf{très petit} (proche de 0). La tension $u_C(t)$ se rapproche de $E$ sans jamais la dépasser.
    \item On écrit : $\lim_{t \to +\infty} u_C(t) = E$. C'est une \textbf{limite finie en $+\infty$}.
\end{enumerate}
Cette situation illustre l'idée centrale : \emph{« quand $x$ tend vers quelque chose, $f(x)$ tend vers quelque chose d'autre »}.
\end{experience}

\begin{definition}[Limite finie en un point (intuitive)]
Soit $f$ une fonction définie autour de $a$ (sauf peut-être en $a$) et $\ell \in \mathbb{R}$.
On dit que $f$ admet la limite $\ell$ en $a$ si, lorsque $x$ se rapproche arbitrairement de $a$ (sans être égal à $a$), $f(x)$ se rapproche arbitrairement de $\ell$.
Notation : $\lim_{x \to a} f(x) = \ell$.
\end{definition}

\begin{definition}[Limites infinies et limites en $\pm\infty$]
\begin{itemize}
    \item \textbf{Limite infinie en $a$ :} $\lim_{x \to a} f(x) = +\infty$ (resp. $-\infty$) si $f(x)$ devient arbitrairement grand (resp. petit) quand $x \to a$. La courbe a une \textbf{asymptote verticale} $x=a$.
    \item \textbf{Limite finie en $+\infty$ :} $\lim_{x \to +\infty} f(x) = \ell$ si $f(x) \to \ell$ quand $x$ devient arbitrairement grand. La courbe a une \textbf{asymptote horizontale} $y=\ell$.
    \item \textbf{Limite infinie en $+\infty$ :} $\lim_{x \to +\infty} f(x) = +\infty$ si $f(x)$ croît sans borne.
\end{itemize}
\end{definition}

\begin{aretenir}[Vocabulaire et notations du Probatoire]
\begin{itemize}
    \item $x \to a$ : « $x$ tend vers $a$ » ($x$ approche $a$ par des valeurs $\neq a$).
    \item $x \to a^-$ (resp. $a^+$) : limite \textbf{à gauche} (resp. \textbf{à droite}).
    \item $\lim_{x \to a} f(x) = \ell$ : la limite \textbf{existe et est finie}.
    \item $\lim_{x \to a} f(x) = +\infty$ ou $-\infty$ : la limite \textbf{n'existe pas au sens fini} (on dit « limite infinie »).
    \item \textbf{Existence de la limite en $a$} $\iff$ limites latérales égales et finies : $\lim_{x \to a^-} f(x) = \lim_{x \to a^+} f(x) = \ell \in \mathbb{R}$.
\end{itemize}
\end{aretenir}

\begin{exemplebox}[Lecture graphique de limites]
\begin{popfigure}
\begin{tikzpicture}[
    axis/.style={->, >=Stealth, thick, popdark},
    grid/.style={very thin, popmuted},
    tick/.style={thick, popdark},
    curve/.style={popblue, very thick, samples=200},
    asympt/.style={poporange, dashed, thick},
    pt/.style={draw=popgreen, fill=popgreen, circle, inner sep=2pt},
    hole/.style={draw=poppink, fill=white, thick, circle, inner sep=2.5pt},
]
% Repère
\draw[grid] (-3,-3) grid (5,4);
\draw[axis] (-3,0) -- (5,0) node[right] {$x$};
\draw[axis] (0,-3) -- (0,4) node[above] {$y$};
\foreach \x in {-2,-1,1,2,3,4} \draw[tick] (\x,-0.1) -- (\x,0.1) node[below=2pt,font=\small] {\x};
\foreach \y in {-2,-1,1,2,3} \draw[tick] (-0.1,\y) -- (0.1,\y) node[left=2pt,font=\small] {\y};

% Asymptote verticale x=1
\draw[asympt] (1,-3) -- (1,4) node[above, poporange, font=\small] {$x=1$};
% Asymptote horizontale y=2
\draw[asympt] (-3,2) -- (5,2) node[right, poporange, font=\small] {$y=2$};

% Branche gauche (x<1) : hyperbole type -1/(x-1) + 2 -> tend vers +oo en 1-, vers 2 en -oo
\draw[curve, domain=-2.5:0.9] plot (\x, {2 - 1/(\x-1)});
% Branche droite (x>1) : hyperbole type 1/(x-1) + 2 -> tend vers +oo en 1+, vers 2 en +oo
\draw[curve, domain=1.1:4.5] plot (\x, {2 + 1/(\x-1)});

% Point défini en x=1 ? Non, trou.
% Point défini en x=3 ? Oui, f(3)=2.5
\fill[pt] (3, 2.5) node[above right] {$f(3)=2,5$};
% Trou en (1,2) limite horizontale
\draw[hole] (1,2) circle (3pt);
\end{tikzpicture}
\legende{Lecture de limites : $\lim_{x\to 1^-}f(x)=+\infty$, $\lim_{x\to 1^+}f(x)=+\infty$, $\lim_{x\to 1}f(x)=+\infty$. $\lim_{x\to -\infty}f(x)=2$, $\lim_{x\to +\infty}f(x)=2$. En $x=3$, limite $=2,5 = f(3)$.}
\end{popfigure}
\end{exemplebox}

\courssub{Limites de référence et interprétation graphique}

\begin{propriete}[Limites de référence (à connaître par cœur)]
Pour tout entier naturel $n \ge 1$ :
\begin{align*}
    &\lim_{x \to +\infty} x^n = +\infty, \quad \lim_{x \to -\infty} x^n = \begin{cases} +\infty & \text{si $n$ pair} \\ -\infty & \text{si $n$ impair} \end{cases} \\[4pt]
    &\lim_{x \to 0^+} \frac{1}{x} = +\infty, \quad \lim_{x \to 0^-} \frac{1}{x} = -\infty, \quad \lim_{x \to 0} \frac{1}{x} \text{ n'existe pas (pas de limite finie ni infinie unique)} \\[4pt]
    &\lim_{x \to +\infty} \frac{1}{x} = 0, \quad \lim_{x \to -\infty} \frac{1}{x} = 0 \\[4pt]
    &\lim_{x \to 0^+} \sqrt{x} = 0, \quad \lim_{x \to +\infty} \sqrt{x} = +\infty
\end{align*}
Pour une fonction polynôme $P(x) = a_n x^n + \dots + a_0$ ($a_n \neq 0$) :
\[
\lim_{x \to \pm\infty} P(x) = \lim_{x \to \pm\infty} a_n x^n.
\]
Pour une fonction rationnelle $R(x) = \frac{P(x)}{Q(x)}$ :
\[
\lim_{x \to \pm\infty} R(x) = \lim_{x \to \pm\infty} \frac{a_n x^n}{b_m x^m} \quad \text{(termes de plus haut degré)}.
\]
\end{propriete}

\begin{methode}[Déterminer une limite par lecture graphique]
\begin{enumerate}
    \item Repérer l'abscisse $a$ (ou la direction $\pm\infty$).
    \item Observer l'ordonnée vers laquelle la courbe \emph{tend} quand $x$ s'approche de $a$ (par la gauche, la droite, ou les deux).
    \item Si la courbe monte/descend sans borne : limite infinie (asymptote verticale).
    \item Si la courbe se stabilise vers une horizontale : limite finie (asymptote horizontale).
    \item Attention : la valeur $f(a)$ (point plein) peut être différente de la limite (trou) ou n'exister pas.
\end{enumerate}
\end{methode}

\courssub{Opérations sur les limites}

Les opérations algébriques se transmettent aux limites \textbf{si les limites intermédiaires existent et ne créent pas de forme indéterminée}.

\begin{propriete}[Règles de calcul des limites (admises)]
Soit $f,g$ deux fonctions et $a \in \mathbb{R} \cup \{-\infty, +\infty\}$. On suppose que $\lim_{x \to a} f(x)$ et $\lim_{x \to a} g(x)$ existent (finies ou infinies).
\begin{enumerate}
    \item \textbf{Somme :} $\lim (f+g) = \lim f + \lim g$.
    \item \textbf{Produit :} $\lim (f \times g) = \lim f \times \lim g$.
    \item \textbf{Quotient :} $\lim \frac{f}{g} = \frac{\lim f}{\lim g}$ \textbf{si $\lim g \neq 0$}.
    \item \textbf{Composition :} Si $\lim_{x \to a} f(x) = \ell$ et $g$ continue en $\ell$, alors $\lim_{x \to a} g(f(x)) = g(\ell)$.
\end{enumerate}
\end{propriete}

\begin{attention}[Formes indéterminées (FI) — Interdiction de conclure]
Les combinaisons suivantes \textbf{ne donnent pas de résultat direct} : il faut \textbf{transformer l'expression} avant de calculer la limite.
\begin{center}
\begin{tabular}{|c|c|}\hline
\textbf{Type} & \textbf{Formes indéterminées} \\ \hline
Somme / Différence & $+\infty - \infty$ \\ \hline
Produit & $0 \times \infty$ (ou $0 \times (-\infty)$) \\ \hline
Quotient & $\frac{0}{0}$, $\frac{\infty}{\infty}$, $\frac{-\infty}{\infty}$, etc. \\ \hline
Puissance & $1^{\infty}$, $0^0$, $\infty^0$ \\ \hline
\end{tabular}
\end{center}
\textbf{Règle d'or :} Ne jamais écrire « $= \text{FI}$ » comme résultat final. Écrire « Forme indéterminée, on factorise / rationalise / utilise les croissances comparées... ».
\end{attention}

\courssub{Calculs pratiques : lever les formes indéterminées}

\begin{methode}[Boîte à outils pour lever les FI (Première Technique)]
\begin{enumerate}
    \item \textbf{FI $\frac{0}{0}$ ou $\frac{\infty}{\infty}$ pour fonctions rationnelles :} Factoriser numérateur et dénominateur par le terme de plus haut degré (ou factoriser par $(x-a)$ si $x\to a$). Simplifier.
    \item \textbf{FI $\frac{0}{0}$ avec racines carrées :} \textbf{Rationaliser} (multiplier numérateur et dénominateur par le conjugué).
    \item \textbf{FI $+\infty - \infty$ (polynômes) :} Factoriser par le terme de plus haut degré.
    \item \textbf{FI $0 \times \infty$ :} Réécrire en quotient ($\frac{0}{1/\infty}$ ou $\frac{\infty}{1/0}$) pour retomber sur $\frac{0}{0}$ ou $\frac{\infty}{\infty}$.
    \item \textbf{Comparaison de croissances (à $+\infty$) :} $x^\alpha \ll x^\beta \ll e^x$ (si au programme) ; ici : tout polynôme est négligeable devant son terme de plus haut degré.
\end{enumerate}
\end{methode}

\begin{exoresolu}[FI $\frac{0}{0}$ : Factorisation]
\textbf{Calculer} $\lim_{x \to 2} \frac{x^2 - 4}{x - 2}$.
\textbf{Solution :}
Pour $x=2$, numérateur $=0$, dénominateur $=0$ $\Rightarrow$ FI $\frac{0}{0}$.
Factorisons : $x^2-4 = (x-2)(x+2)$.
\[
\frac{x^2-4}{x-2} = \frac{(x-2)(x+2)}{x-2} = x+2 \quad (x \neq 2).
\]
Donc $\lim_{x \to 2} \frac{x^2-4}{x-2} = \lim_{x \to 2} (x+2) = 4$.
\end{exoresolu}

\begin{exoresolu}[FI $\frac{0}{0}$ : Rationalisation]
\textbf{Calculer} $\lim_{x \to 4} \frac{\sqrt{x} - 2}{x - 4}$.
\textbf{Solution :}
FI $\frac{0}{0}$. Conjugué de $\sqrt{x}-2$ est $\sqrt{x}+2$.
\[
\frac{\sqrt{x}-2}{x-4} = \frac{(\sqrt{x}-2)(\sqrt{x}+2)}{(x-4)(\sqrt{x}+2)} = \frac{x-4}{(x-4)(\sqrt{x}+2)} = \frac{1}{\sqrt{x}+2} \quad (x \neq 4).
\]
Limite $= \frac{1}{\sqrt{4}+2} = \frac{1}{4}$.
\end{exoresolu}

\begin{exoresolu}[FI $\frac{\infty}{\infty}$ : Terme de plus haut degré]
\textbf{Calculer} $\lim_{x \to +\infty} \frac{3x^2 - 5x + 1}{2x^2 + 7}$.
\textbf{Solution :}
Degrés égaux (2). FI $\frac{\infty}{\infty}$.
Factorisons par $x^2$ (terme de plus haut degré) :
\[
\frac{3x^2 - 5x + 1}{2x^2 + 7} = \frac{x^2(3 - \frac{5}{x} + \frac{1}{x^2})}{x^2(2 + \frac{7}{x^2})} = \frac{3 - \frac{5}{x} + \frac{1}{x^2}}{2 + \frac{7}{x^2}}.
\]
Puisque $\lim_{x\to+\infty} \frac{1}{x} = 0$ et $\lim_{x\to+\infty} \frac{1}{x^2} = 0$ :
Limite $= \frac{3 - 0 + 0}{2 + 0} = \frac{3}{2}$.
\end{exoresolu}

\begin{exoresolu}[FI $+\infty - \infty$ : Différence de racines]
\textbf{Calculer} $\lim_{x \to +\infty} \left( \sqrt{x+1} - \sqrt{x} \right)$.
\textbf{Solution :}
FI $+\infty - \infty$. Rationalisons (multiplier par le conjugué $\sqrt{x+1}+\sqrt{x}$) :
\[
\sqrt{x+1} - \sqrt{x} = \frac{(\sqrt{x+1}-\sqrt{x})(\sqrt{x+1}+\sqrt{x})}{\sqrt{x+1}+\sqrt{x}} = \frac{x+1-x}{\sqrt{x+1}+\sqrt{x}} = \frac{1}{\sqrt{x+1}+\sqrt{x}}.
\]
Dénominateur $\to +\infty$, donc limite $= 0$.
\end{exoresolu}

\coursec{Continuité d'une fonction}

Après avoir maîtrisé le calcul des limites, notamment le traitement des formes indéterminées, nous disposons maintenant de l'outil principal pour analyser le comportement local d'une fonction. La notion de \cle{continuité} formalise l'idée intuitive d'une évolution « sans rupture » : une grandeur physique (température, tension, position) qui varie sans faire de saut brutal.

\courssub{Approche intuitive : « sans lever le crayon »}

Imaginez que vous traciez la courbe représentative d'une fonction $f$ sur une feuille de papier. Si vous pouvez dessiner cette courbe sur un intervalle donné sans jamais devoir lever votre crayon, alors la fonction est \textbf{continue} sur cet intervalle. À l'inverse, si vous devez interrompre votre trait pour « sauter » d'un point à un autre (trou, saut, asymptote verticale), la fonction présente une \cle{discontinuité}.

\begin{experience}[Découverte : continuité ou rupture ?]
Considérons trois situations concrètes au laboratoire ou sur un chantier :
\begin{enumerate}
    \item \textbf{Température d'une pièce chauffée :} Elle monte progressivement de 18 °C à 22 °C. À tout instant, elle prend \emph{toutes} les valeurs intermédiaires. Courbe sans rupture.
    \item \textbf{Nombre d'ouvriers sur un chantier :} Il passe brutalement de 10 à 11 à l'arrivée d'un nouvel ouvrier. La courbe fait un « saut » (escalier). Discontinuité.
    \item \textbf{Tension aux bornes d'un condensateur qui se charge :} Elle suit une loi exponentielle continue. Mais si le condensateur se décharge brutalement (court-circuit), la tension chute instantanément : discontinuité.
\end{enumerate}
En mathématiques, nous devons remplacer cette intuition graphique par une définition précise, utilisable sans tracer la courbe.
\end{experience}

\courssub{Définition rigoureuse : continuité en un point}

La continuité en un point $a$ exige trois conditions simultanées : la fonction existe en $a$, elle admet une limite finie en $a$, et cette limite coïncide avec la valeur de la fonction.

\begin{definition}[Continuité en un point]
Soit $f$ une fonction définie sur un intervalle $I$ et $a \in I$.
On dit que $f$ est \textbf{continue en $a$} si et seulement si :
\[
\lim_{x \to a} f(x) = f(a).
\]
Cette égalité implique trois faits :
\begin{enumerate}
    \item $f$ est \textbf{définie} en $a$ ($a$ appartient à l'ensemble de définition $D_f$).
    \item La \textbf{limite} de $f$ en $a$ \textbf{existe et est finie}.
    \item Cette limite est \textbf{égale} à l'image $f(a)$.
\end{enumerate}
Si l'une de ces conditions n'est pas vérifiée, $f$ est \textbf{discontinue} en $a$.
\end{definition}

\begin{attention}[Piège classique : « définie $\neq$ continue »]
Une fonction peut être définie en un point sans y être continue.
\textbf{Exemple :} $f(x) = \begin{cases} x^2 & \text{si } x \neq 1 \\ 5 & \text{si } x = 1 \end{cases}$.
$f$ est définie en $1$ ($f(1)=5$), $\lim_{x \to 1} f(x) = 1$, mais $1 \neq 5$.
La courbe a un « trou » en $(1,1)$ et un point isolé en $(1,5)$. Il faut lever le crayon.
\end{attention}

\courssub{Fonctions continues usuelles et contre-exemples}

La plupart des fonctions élémentaires manipulées en Première technique sont continues sur leur ensemble de définition.

\begin{propriete}[Fonctions continues de référence (admis)]
Les fonctions suivantes sont continues sur \emph{chaque intervalle} inclus dans leur ensemble de définition :
\begin{itemize}
    \item Les \textbf{fonctions polynômes} (donc aussi les fonctions affines, constantes) sur $\mathbb{R}$.
    \item La fonction \textbf{racine carrée} $x \mapsto \sqrt{x}$ sur $[0 ; +\infty[$.
    \item Les \textbf{fonctions rationnelles} $x \mapsto \frac{P(x)}{Q(x)}$ sur tout intervalle ne contenant pas de zéro du dénominateur $Q$.
    \item La fonction \textbf{valeur absolue} $x \mapsto |x|$ sur $\mathbb{R}$.
    \item La fonction \textbf{partie entière} $x \mapsto \lfloor x \rfloor$ (continue sur chaque intervalle $[n ; n+1[$ mais discontinue aux entiers relatifs $n$).
\end{itemize}
\end{propriete}

\begin{exemplebox}[Exemples de discontinuités classiques]
\begin{enumerate}
    \item \textbf{Discontinuité de la fonction inverse} $f(x) = \frac{1}{x}$ en $0$ : $f$ n'est pas définie en $0$, la limite est infinie ($\pm\infty$). Asymptote verticale.
    \item \textbf{Discontinuité de la fonction partie entière} $E(x) = \lfloor x \rfloor$ en $x = 2$ :
    $\lim_{x \to 2^-} E(x) = 1$, $\lim_{x \to 2^+} E(x) = 2$, $E(2) = 2$.
    La limite n'existe pas (limites latérales différentes) : \textbf{saut}.
    \item \textbf{Discontinuité évitable (trou)} : $f(x) = \frac{x^2-1}{x-1}$ en $1$.
    $f$ n'est pas définie en $1$. $\lim_{x \to 1} f(x) = \lim_{x \to 1} (x+1) = 2$.
    Si on \textbf{prolonge} $f$ par $f(1)=2$, la fonction devient continue. C'est une discontinuité « évitable ».
\end{enumerate}
\end{exemplebox}

\begin{popfigure}
\begin{tikzpicture}[
    axis/.style={->, >=Stealth, thick, popdark},
    grid/.style={very thin, popmuted},
    tick/.style={thick, popdark},
    cont/.style={popblue, very thick, samples=200},
    disc/.style={poporange, very thick, samples=200},
    hole/.style={draw=poppink, fill=white, thick, circle, inner sep=2pt},
    isolated/.style={draw=poppink, fill=poppink, thick, circle, inner sep=2pt},
]
% Figure 1 : Continue (x^3 - x)
\begin{scope}[xshift=0cm]
    \draw[grid] (-2.5,-2.5) grid (2.5,2.5);
    \draw[axis] (-2.5,0) -- (2.5,0) node[right] {$x$};
    \draw[axis] (0,-2.5) -- (0,2.5) node[above] {$y$};
    \foreach \x in {-2,-1,1,2} \draw[tick] (\x,-0.1) -- (\x,0.1) node[below=2pt,font=\tiny] {\x};
    \foreach \y in {-2,-1,1,2} \draw[tick] (-0.1,\y) -- (0.1,\y) node[left=2pt,font=\tiny] {\y};
    \draw[cont] plot (\x, {\x*\x*\x - \x});
    \node[popblue, font=\small\bfseries] at (0, -2.8) {(a) Continue : $x^3-x$};
\end{scope}
% Figure 2 : Saut (Partie entière)
\begin{scope}[xshift=6.5cm]
    \draw[grid] (-0.5,-0.5) grid (3.5,3.5);
    \draw[axis] (-0.5,0) -- (3.5,0) node[right] {$x$};
    \draw[axis] (0,-0.5) -- (0,3.5) node[above] {$y$};
    \foreach \x in {1,2,3} \draw[tick] (\x,-0.1) -- (\x,0.1) node[below=2pt,font=\tiny] {\x};
    \foreach \y in {1,2,3} \draw[tick] (-0.1,\y) -- (0.1,\y) node[left=2pt,font=\tiny] {\y};
    % Escalier E(x)
    \draw[disc] (0,0) -- (1,0);
    \draw[disc] (1,1) -- (2,1);
    \draw[disc] (2,2) -- (3,2);
    \draw[disc] (3,3) -- (3.5,3);
    % Points pleins (valeur à droite/entier)
    \foreach \x/\y in {0/0, 1/1, 2/2, 3/3} \fill[poporange] (\x,\y) circle (2.5pt);
    % Points vides (limite à gauche)
    \foreach \x/\y in {1/0, 2/1, 3/2} \draw[hole] (\x,\y) circle (2.5pt);
    \node[poporange, font=\small\bfseries] at (1.5, -0.8) {(b) Saut : $E(x)$};
\end{scope}
% Figure 3 : Trou (Discontinuité évitable)
\begin{scope}[xshift=13cm]
    \draw[grid] (-1.5,-0.5) grid (2.5,3.5);
    \draw[axis] (-1.5,0) -- (2.5,0) node[right] {$x$};
    \draw[axis] (0,-0.5) -- (0,3.5) node[above] {$y$};
    \foreach \x in {-1,1,2} \draw[tick] (\x,-0.1) -- (\x,0.1) node[below=2pt,font=\tiny] {\x};
    \foreach \y in {1,2,3} \draw[tick] (-0.1,\y) -- (0.1,\y) node[left=2pt,font=\tiny] {\y};
    % Courbe y = x+1 mais trou en x=1
    \draw[cont, domain=-1.5:0.9] plot (\x, {\x+1});
    \draw[cont, domain=1.1:2.5] plot (\x, {\x+1});
    % Trou en (1,2)
    \draw[hole] (1,2) circle (3pt);
    % Point isolé ailleurs (ex: f(1)=0)
    \fill[isolated] (1,0) circle (2.5pt);
    \node[poppink, font=\small\bfseries] at (0.5, -0.8) {(c) Trou : $\frac{x^2-1}{x-1}$};
\end{scope}
\end{tikzpicture}
\legende{Trois types de comportements : (a) fonction continue (polynôme), (b) fonction à sauts (partie entière), (c) fonction avec un trou (discontinuité évitable).}
\end{popfigure}

\courssub{Continuité sur un intervalle}

La définition globale découle naturellement de la définition locale.

\begin{definition}[Continuité sur un intervalle]
Soit $I$ un intervalle (ouvert, fermé, demi-ouvert, borné ou non).
Une fonction $f$ est \textbf{continue sur $I$} si et seulement si elle est continue en \textbf{tout point} de $I$.
\begin{itemize}
    \item Si $I = [a ; b]$ est fermé, on exige de plus la continuité \textbf{à droite} en $a$ ($\lim_{x \to a^+} f(x) = f(a)$) et \textbf{à gauche} en $b$ ($\lim_{x \to b^-} f(x) = f(b)$).
    \item Une fonction continue sur un intervalle fermé $[a ; b]$ est \textbf{bornée} et \textbf{atteint ses bornes} (théorème des valeurs extrêmes, admis).
\end{itemize}
\end{definition}

\begin{aretenir}[Vocabulaire technique]
\begin{itemize}
    \item \textbf{Continuité à droite en $a$ :} $\lim_{x \to a^+} f(x) = f(a)$.
    \item \textbf{Continuité à gauche en $b$ :} $\lim_{x \to b^-} f(x) = f(b)$.
    \item \textbf{Prolongement par continuité :} Si $f$ n'est pas définie en $a$ mais admet une limite finie $\ell$ en $a$, la fonction $\tilde{f}$ définie par $\tilde{f}(x)=f(x)$ pour $x\neq a$ et $\tilde{f}(a)=\ell$ est continue en $a$.
\end{itemize}
\end{aretenir}

\courssub{Théorème des valeurs intermédiaires (TVI)}

C'est le théorème majeur de cette section. Il garantit l'existence de solutions sans avoir à les calculer explicitement.

\begin{propriete}[Théorème des valeurs intermédiaires — TVI (admis, illustration graphique)]
Soit $f$ une fonction \textbf{continue sur un intervalle fermé $[a ; b]$} (avec $a < b$).
Soit $k$ un nombre réel \textbf{compris entre $f(a)$ et $f(b)$} (c'est-à-dire $f(a) \le k \le f(b)$ ou $f(b) \le k \le f(a)$).
Alors il existe \textbf{au moins un} réel $c \in [a ; b]$ tel que $f(c) = k$.
\end{propriete}

\begin{popfigure}
\begin{tikzpicture}[
    axis/.style={->, >=Stealth, thick, popdark},
    grid/.style={very thin, popmuted},
    tick/.style={thick, popdark},
    curve/.style={popblue, very thick, samples=200, domain=1:4},
    kval/.style={poporange, dashed, thick},
    pt/.style={draw=popgreen, fill=popgreen, thick, circle, inner sep=2.5pt},
]
% Repère
\draw[grid] (0.5,0.5) grid (4.5, 3.5);
\draw[axis] (0.5,0.5) -- (4.5,0.5) node[right] {$x$};
\draw[axis] (0.5,0.5) -- (0.5,3.5) node[above] {$y$};
% Graduations
\foreach \x/\xl in {1/a, 2/, 3/, 4/b} \draw[tick] (\x,0.4) -- (\x,0.6) node[below=2pt] {$\xl$};
\foreach \y/\yl in {1/f(a), 2/k, 3/f(b)} \draw[tick] (0.4,\y) -- (0.6,\y) node[left=2pt] {$\yl$};
% Courbe continue monotone croissante passant par (1,1) et (4,3) : f(x) = 1 + (2/9)*(x-1)^2
\draw[curve] plot (\x, {1 + (2/9)*(\x-1)*(\x-1)});
% Ligne y = k (k=2)
\draw[kval] (0.5,2) -- (4.5,2) node[right, poporange] {$y = k$};
% Points A, B
\fill[pt] (1,1) node[below left] {$A(a, f(a))$};
\fill[pt] (4,3) node[above right] {$B(b, f(b))$};
% Intersection (solution de f(x)=2)
% 1 + 2/9*(x-1)^2 = 2 => (x-1)^2 = 9/2 => x-1 = 3/sqrt(2) approx 2.12 => x approx 3.12
\pgfmathsetmacro{\cx}{1 + 3/sqrt(2)}
\fill[pt] (\cx, 2) node[above right] {$C(c, k)$};
% Flèches explicatives
\draw[popgreen, <->, thick] (4.7, 1) -- (4.7, 3) node[midway, right, font=\small] {$f(b)-f(a)$};
\draw[poppurple, <->, thick] (4.9, 1) -- (4.9, 2) node[midway, right, font=\small] {$k-f(a)$};
\end{tikzpicture}
\legende{Illustration du TVI : la courbe continue reliant $A$ et $B$ coupe nécessairement la droite horizontale $y=k$.}
\end{popfigure}

\begin{savaistu}[Le TVI dans la vie réelle : la température à Douala]
Le TVI n'est pas qu'un outil abstrait. Imaginez la température à Douala un après-midi d'avril. À 12h, il fait 25 °C. À 15h, il fait 32 °C. La température varie \emph{continûment} (elle ne saute pas brutalement de 25 à 32 sans passer par 28). Le TVI garantit qu'il y a eu \textbf{au moins un instant} entre 12h et 15h où la température était \textbf{exactement 28 °C} (ou 30 °C, ou n'importe quelle valeur entre 25 et 32). C'est ce principe qui permet aux thermomètres numériques d'afficher une valeur précise : ils mesurent une grandeur physique continue.
\end{savaistu}

\courssub{Application du TVI : existence de solutions d'équations}

L'usage principal du TVI en Première technique est de prouver qu'une équation $f(x) = k$ admet au moins une solution dans un intervalle donné, sans la résoudre.

\begin{methode}[Démontrer l'existence d'une solution par le TVI]
Pour montrer que l'équation $f(x) = k$ a une solution dans $[a ; b]$ :
\begin{enumerate}
    \item \textbf{Vérifier les hypothèses :} $f$ est continue sur $[a ; b]$ (citer la nature de $f$ : polynôme, rationnelle sur son ensemble de définition, racine carrée, etc.).
    \item \textbf{Calculer $f(a)$ et $f(b)$.}
    \item \textbf{Vérifier l'encadrement :} montrer que $k$ est entre $f(a)$ et $f(b)$ (souvent $f(a) < k < f(b)$ ou l'inverse).
    \item \textbf{Conclure :} « D'après le TVI, il existe au moins un $c \in [a ; b]$ tel que $f(c) = k$. L'équation admet au moins une solution dans $[a ; b]$. »
\end{enumerate}
\end{methode}

\begin{exoresolu}[Application du TVI : équation $x^3 + x - 3 = 0$]
\textbf{Énoncé :} Montrer que l'équation $x^3 + x - 3 = 0$ admet au moins une solution dans l'intervalle $[1 ; 2]$.

\textbf{Solution détaillée :}
\begin{enumerate}
    \item \textbf{Fonction associée :} Posons $f(x) = x^3 + x - 3$. C'est une \textbf{fonction polynôme}, donc \textbf{continue sur $\mathbb{R}$}, et en particulier sur l'intervalle $[1 ; 2]$.
    \item \textbf{Calcul des bornes :}
    \begin{align*}
        f(1) &= 1^3 + 1 - 3 = -1. \\
        f(2) &= 2^3 + 2 - 3 = 8 + 2 - 3 = 7.
    \end{align*}
    \item \textbf{Encadrement de 0 :} On a $f(1) = -1 < 0$ et $f(2) = 7 > 0$.
    Donc $0$ est bien compris entre $f(1)$ et $f(2)$.
    \item \textbf{Conclusion :} Les hypothèses du TVI sont satisfaites ($f$ continue sur $[1;2]$, $0 \in [f(1); f(2)]$).
    \textbf{Il existe au moins un réel $c \in [1 ; 2]$ tel que $f(c) = 0$.}
    L'équation $x^3 + x - 3 = 0$ admet donc au moins une solution dans $[1 ; 2]$.
\end{enumerate}
\textit{Remarque :} La fonction $f$ est strictement croissante sur $\mathbb{R}$ (car pour $x_2 > x_1$, $f(x_2)-f(x_1) = (x_2-x_1)(x_2^2+x_1x_2+x_1^2+1) > 0$), donc la solution est \textbf{unique}. Mais le TVI seul ne donne que l'existence.
\end{exoresolu}

\begin{exoresolu}[Recherche d'intervalle pour une équation rationnelle]
\textbf{Énoncé :} Montrer que l'équation $\frac{2x+1}{x-1} = 3$ admet une solution dans $[2 ; 4]$.

\textbf{Solution :}
\begin{enumerate}
    \item Posons $f(x) = \frac{2x+1}{x-1}$. C'est une \textbf{fonction rationnelle}. Son ensemble de définition est $\mathbb{R}^* = \mathbb{R} \setminus \{1\}$.
    L'intervalle $[2 ; 4]$ ne contient pas $1$, donc $f$ est \textbf{continue sur $[2 ; 4]$}.
    \item Calculons les images des bornes :
    \begin{align*}
        f(2) &= \frac{2(2)+1}{2-1} = \frac{5}{1} = 5. \\
        f(4) &= \frac{2(4)+1}{4-1} = \frac{9}{3} = 3.
    \end{align*}
    \item On cherche une solution pour $k=3$. On a $f(4) = 3$ et $f(2) = 5$.
    La valeur $k=3$ est atteinte en $x=4$ (borne de l'intervalle).
    \textbf{Conclusion :} $x=4$ est une solution dans $[2 ; 4]$. (Le TVI garantit l'existence, ici on la trouve explicitement car c'est une borne).
    \textit{Note :} Si on cherchait $k=4$, $f(2)=5 > 4 > 3 = f(4)$, le TVI donnerait une solution dans $]2;4[$.
\end{enumerate}
\end{exoresolu}

\begin{attention}[Conditions d'application du TVI : checklist obligatoire]
Ne jamais écrire « d'après le TVI » sans avoir explicitement vérifié :
\begin{enumerate}
    \item \textbf{Continuité sur l'intervalle FERMÉ $[a;b]$.} (Si $f$ n'est pas définie en $a$ ou $b$, ou si elle a un saut, le TVI ne s'applique pas).
    \item \textbf{Valeur $k$ entre $f(a)$ et $f(b)$.} (Si $k$ est en dehors, le TVI ne dit rien : il peut y avoir une solution ou pas, le théorème ne conclut pas).
    \item \textbf{Sens de l'inégalité :} $f(a) \le k \le f(b)$ \textbf{OU} $f(b) \le k \le f(a)$. Ne pas oublier le cas où $f$ est décroissante ($f(a) > f(b)$).
\end{enumerate}
\end{attention}

\courssub{Exercices d'application}

\begin{exercice}[Exercice 1 : Continuité de fonctions définies par morceaux]
On considère la fonction $f$ définie sur $\mathbb{R}$ par :
\[
f(x) = \begin{cases}
x^2 + 1 & \text{si } x \le 1 \\
ax + b & \text{si } x > 1
\end{cases}
\]
Déterminer les réels $a$ et $b$ pour que $f$ soit continue sur $\mathbb{R}$.
\end{exercice}

\begin{exercice}[Exercice 2 : TVI et racine carrée]
Soit $f(x) = \sqrt{x+3} - x$ définie sur $[-3 ; +\infty[$.
\begin{enumerate}
    \item Justifier que $f$ est continue sur $[0 ; 3]$.
    \item Montrer que l'équation $\sqrt{x+3} = x$ admet une unique solution dans $[0 ; 3]$.
    \item Encadrer cette solution à $10^{-1}$ près par une méthode de balayage (tableau de valeurs).
\end{enumerate}
\end{exercice}

\begin{exercice}[Exercice 3 : Situation concrète — Résistance d'un thermistance]
La résistance $R$ (en ohms) d'une sonde de température (thermistance NTC) en fonction de la température $\theta$ (en °C) est modélisée sur $[0 ; 100]$ par :
\[
R(\theta) = 10000 \cdot e^{-0.04 \theta} \quad (\text{modèle simplifié exponentiel})
\]
On admet que cette fonction est continue et strictement décroissante sur $[0 ; 100]$.
\begin{enumerate}
    \item Calculer $R(0)$ et $R(100)$.
    \item Le système d'alarme se déclenche si $R(\theta) < 2000 \Omega$. En utilisant le TVI, montrer qu'il existe une unique température critique $\theta_c$ dans $[0 ; 100]$ pour laquelle $R(\theta_c) = 2000$.
    \item Donner un encadrement de $\theta_c$ à l'unité près (à l'aide de la calculatrice ou par lecture graphique).
\end{enumerate}
\end{exercice}

\begin{corrige}[Exercice 1]
Pour que $f$ soit continue en $1$ (seul point de raccord suspect), il faut :
$\lim_{x \to 1^-} f(x) = f(1) = \lim_{x \to 1^+} f(x)$.
\begin{itemize}
    \item $f(1) = 1^2 + 1 = 2$.
    \item $\lim_{x \to 1^-} f(x) = 2$.
    \item $\lim_{x \to 1^+} f(x) = a(1) + b = a + b$.
\end{itemize}
Condition de continuité : $a + b = 2$.
Il y a une infinité de couples $(a,b)$ solutions (ex: $a=0, b=2$ ; $a=1, b=1$ ; $a=2, b=0$...).
Toute droite $y=ax+b$ passant par le point $(1; 2)$ convient.
\end{corrige}

\begin{corrige}[Exercice 2]
\begin{enumerate}
    \item $x \mapsto x+3$ est continue sur $\mathbb{R}$, $\sqrt{\cdot}$ est continue sur $[0;+\infty[$, $x \mapsto x$ est continue. La composition et la somme/différence conservent la continuité. $f$ est continue sur $[-3;+\infty[$, donc sur $[0;3]$.
    \item $f(0) = \sqrt{3} - 0 \approx 1,73 > 0$. $f(3) = \sqrt{6} - 3 \approx 2,45 - 3 = -0,55 < 0$.
    $0 \in [f(3); f(0)]$. Par le TVI, $\exists c \in [0;3], f(c)=0 \iff \sqrt{c+3}=c$.
    Unicité : On observe que sur $[0;3]$, la fonction $x \mapsto \sqrt{x+3}$ est croissante mais sa pente diminue, tandis que $x \mapsto x$ est croissante avec pente constante 1. On peut vérifier par un tableau de valeurs que $f$ ne change qu'une fois de signe (ou admettre que $f$ est strictement décroissante car la dérivée $f'(x) = \frac{1}{2\sqrt{x+3}} - 1$ est négative sur $[0;3]$, notion de Terminale). Donc la solution est unique.
    \item Balayage (valeurs approchées) :
    \begin{center}
    \begin{tabular}{|c|c|c|c|c|c|}\hline
    $x$ & 2,0 & 2,1 & 2,2 & 2,3 & 2,4 \\ \hline
    $f(x)$ & 0,236 & 0,118 & -0,004 & -0,123 & -0,239 \\ \hline
    \end{tabular}
    \end{center}
    Changement de signe entre $2,2$ et $2,3$. Solution $\in [2,2 ; 2,3]$.
    Encadrement à $10^{-1}$ près : $c \in [2,2 ; 2,3]$.
\end{enumerate}
\end{corrige}

\begin{corrige}[Exercice 3]
\begin{enumerate}
    \item $R(0) = 10000 \cdot e^0 = 10000 \Omega$. $R(100) = 10000 \cdot e^{-4} \approx 10000 \cdot 0,0183 \approx 183 \Omega$.
    \item $R$ continue sur $[0;100]$, strictement décroissante. $R(0)=10000 > 2000 > 183 \approx R(100)$.
    Par TVI, $\exists ! \theta_c \in [0;100]$ tq $R(\theta_c)=2000$.
    \item Résolution approchée (calculatrice) : $2000 = 10000 e^{-0.04\theta_c} \iff e^{-0.04\theta_c} = 0,2 \iff -0,04\theta_c = \ln(0,2) \iff \theta_c = -\frac{\ln(0,2)}{0,04} = \frac{\ln 5}{0,04} \approx \frac{1,609}{0,04} \approx 40,2$.
    Encadrement à l'unité près : $\theta_c \in [40 ; 41]$.
    (Vérification : $R(40) \approx 2019 > 2000$, $R(41) \approx 1936 < 2000$).
\end{enumerate}
\end{corrige}

\coursec{Applications à des situations concrètes}

Dans la section précédente, nous avons étudié la \cle{continuité} d'une fonction et le \cle{théorème des valeurs intermédiaires} (TVI) : une fonction continue sur un intervalle ne « saute » pas et prend toutes les valeurs intermédiaires entre deux de ses images. Nous disposons maintenant de tous les outils de la leçon : approche intuitive des limites, limites de référence, opérations, formes indéterminées, continuité. Il est temps de répondre à la question que tout élève de série technique se pose légitimement : \emph{à quoi cela sert-il dans un atelier, sur un chantier, dans une entreprise ?} Cette section montre que les limites et la continuité sont des outils de \cle{décision} : elles permettent de prévoir le comportement d'une grandeur « à long terme » ou « près d'une valeur critique », et donc de fixer des prix, de dimensionner des installations, d'éviter des ruptures.

\courssub{Modéliser un coût moyen de production}

\textbf{Situation concrète.}\par
Une menuiserie de Douala fabrique des sacs de sciure conditionnés (ou un atelier produit des sacs de ciment reconditionnés). Pour lancer la production, il faut d'abord payer des \cle{coûts fixes} : location du local, achat de la machine, installation électrique. Ces coûts existent même si l'on ne produit rien. Ensuite, chaque sac produit entraîne un \cle{coût variable} : matière première, main-d'œuvre, énergie.

Notons :
\begin{itemize}
\item $F$ le coût fixe total (en FCFA), payé une seule fois ;
\item $v$ le coût variable par sac (en FCFA par sac) ;
\item $x$ le nombre de sacs produits.
\end{itemize}

Le \cle{coût total} de production de $x$ sacs est alors :
\[ CT(x) = F + v\,x. \]

Mais ce qui intéresse le chef d'entreprise pour fixer son \emph{prix de vente}, c'est le \cle{coût moyen} d'un sac : combien me revient \emph{en moyenne} un sac produit ?

\begin{definition}[Coût moyen]
Le \cle{coût moyen} de production de $x$ unités est le quotient du coût total par le nombre d'unités produites :
\[ C(x) = \dfrac{CT(x)}{x} = \dfrac{F + v\,x}{x}, \qquad x > 0. \]
\end{definition}

\begin{attention}[Le coût moyen n'est pas le coût d'un sac supplémentaire]
Le coût moyen $C(x)$ répartit \emph{tous} les coûts (fixes compris) sur les $x$ sacs. Ce n'est pas le coût du $(x+1)$-ième sac, qui vaut simplement $v$. Confondre les deux est une erreur classique en gestion.
\end{attention}

\courssub{Comportement quand la production devient grande : asymptote horizontale}

Étudions ce qui se passe quand l'atelier produit \emph{de plus en plus}, c'est-à-dire quand $x$ tend vers $+\infty$. Transformons l'écriture du coût moyen :
\[ C(x) = \dfrac{F + v\,x}{x} = \dfrac{F}{x} + v. \]

Quand $x \to +\infty$, le quotient $\dfrac{F}{x}$ tend vers $0$ (limite de référence : l'inverse d'une quantité infiniment grande est infiniment petit). Par somme de limites :
\[ \lim_{x \to +\infty} C(x) = 0 + v = v. \]

\begin{propriete}[Coût moyen plancher]
Si le coût total est $CT(x) = F + vx$ (avec $F > 0$ et $v > 0$), alors
\[ \lim_{x \to +\infty} C(x) = v. \]
La droite d'équation $y = v$ est \cle{asymptote horizontale} à la courbe de $C$ en $+\infty$. De plus, comme $C(x) = v + \dfrac{F}{x} > v$, le coût moyen est \emph{toujours strictement supérieur} à $v$ : il s'en approche sans jamais l'atteindre.
\end{propriete}

\begin{aretenir}[Interprétation économique d'une asymptote horizontale]
En économie, une asymptote horizontale $y = \ell$ s'interprète comme un \cle{plancher} (ou un plafond) :
\begin{itemize}
\item le coût moyen ne descendra \textbf{jamais} en dessous de $\ell$, même en produisant énormément ;
\item $\ell$ est le \cle{prix limite} : vendre en dessous de $\ell$ FCFA l'unité garantit une perte, quelle que soit la quantité produite.
\end{itemize}
Les coûts fixes se « diluent » dans la masse produite : $\dfrac{F}{x}$ devient négligeable, et seul le coût variable $v$ subsiste.
\end{aretenir}

\begin{exemplebox}[Coût moyen de production de sacs de ciment]
Un atelier produit des sacs de ciment reconditionnés. Coûts fixes : $F = \num{50000}$ FCFA ; coût variable : $v = \num{2000}$ FCFA par sac. Le coût moyen de $x$ sacs est :
\[ C(x) = \dfrac{\num{50000} + \num{2000}\,x}{x} = \num{2000} + \dfrac{\num{50000}}{x} \quad \text{(en FCFA par sac)}. \]
Quelques valeurs :
\begin{center}
\begin{tabularx}{\linewidth}{|l|X|X|X|X|}
\hline
\rowcolor{popblueL} $x$ (sacs) & $10$ & $100$ & $1000$ & $10000$ \\
\hline $C(x)$ (FCFA) & $\num{7000}$ & $\num{2500}$ & $\num{2050}$ & $\num{2005}$ \\
\hline
\end{tabularx}
\end{center}
On observe que $C(x)$ décroît et se rapproche de $\num{2000}$ FCFA sans jamais passer en dessous. Conclusion commerciale : \textbf{le prix de vente d'un sac doit être strictement supérieur à \num{2000} FCFA}, sinon l'atelier perd de l'argent même en produisant énormément.
\end{exemplebox}

\begin{popfigure}
\begin{center}
\begin{tikzpicture}
\begin{axis}[
width=0.95\linewidth, height=7cm,
xlabel={$x$ (nombre de sacs)}, ylabel={$C(x)$ (FCFA)},
xmin=0, xmax=600, ymin=1500, ymax=7500,
domain=10:600, samples=200,
grid=both, grid style={popline!40},
axis lines=left,
]
\addplot[popcurve] {2000 + 50000/x};
\addplot[poporange, dashed, thick, domain=0:600] {2000};
\addplot[popgreen, thick, mark=*, only marks] coordinates {(100,2500) (50,3000)};
\node[poporange, anchor=south west] at (axis cs:330,2050) {\small asymptote $y = 2000$ : plancher du coût moyen};
\node[popblue, anchor=west] at (axis cs:60,3400) {\small $C(50) = 3000$ FCFA};
\node[popblue, anchor=west] at (axis cs:110,2600) {\small $C(100) = 2500$ FCFA};
\node[popdark, anchor=south east] at (axis cs:590,2150) {\small $C(x) \to 2000$ quand $x \to +\infty$};
\end{axis}
\end{tikzpicture}
\end{center}
\legende{Courbe du coût moyen $C(x) = \dfrac{\num{50000} + \num{2000}x}{x}$ et son asymptote horizontale $y = \num{2000}$ : le coût moyen décroît vers \num{2000} FCFA sans jamais l'atteindre.}
\end{popfigure}

\courssub{Asymptote verticale : explosion d'une grandeur près d'une valeur critique}

Certaines grandeurs, au contraire, \emph{explosent} lorsqu'on s'approche d'une valeur critique. C'est le cas chaque fois qu'une formule contient un \emph{dénominateur qui s'annule}.

\begin{exemplebox}[Débit de vidange d'un réservoir]
Un réservoir de $V = 600$ litres doit être vidangé par une pompe. Si la pompe enlève déjà $q$ litres par heure par un autre circuit, le débit disponible est $(q - q_0)$... Prenons un modèle simple : le temps de vidange en fonction du débit $d$ (en L/h) est $T(d) = \dfrac{600}{d}$. Si le débit devient très faible ($d \to 0^+$), alors
\[ \lim_{d \to 0^+} T(d) = +\infty : \]
la vidange devient \emph{interminable}. La droite $d = 0$ est asymptote verticale. Autre exemple : le coût moyen $C(x) = 2000 + \dfrac{\num{50000}}{x}$ vérifie $\displaystyle\lim_{x \to 0^+} C(x) = +\infty$ : produire très peu coûte extrêmement cher par unité, car les coûts fixes sont répartis sur presque rien.
\end{exemplebox}

\begin{aretenir}[Interprétation d'une asymptote verticale]
Une asymptote verticale $x = a$ signale une \cle{valeur critique} du paramètre $x$ : près de $a$, la grandeur étudiée devient arbitrairement grande (ou négative). En pratique : un coût qui explose, un temps qui devient infini, une contrainte mécanique qui dépasse la résistance du matériau. \textbf{Identifier les valeurs critiques, c'est identifier les dangers.}
\end{aretenir}

\courssub{Continuité et vie courante : des grandeurs continues par nature}

Pourquoi la continuité est-elle si naturelle ? Parce que la plupart des grandeurs physiques \emph{ne sautent pas} :
\begin{itemize}
\item la \cle{température} d'un atelier au cours de la journée passe progressivement de 24 °C à 35 °C : elle prend toutes les valeurs intermédiaires ;
\item le \cle{niveau d'eau} d'un réservoir qui se remplit monte de façon continue : pour passer de 40 cm à 90 cm, il est forcément passé par 60 cm ;
\item la \cle{distance parcourue} par un véhicule entre Yaoundé et Bafoussam augmente continûment avec le temps : un compteur ne saute jamais de 120 km à 130 km instantanément.
\end{itemize}

\begin{exemplebox}[Application du TVI : à quelle heure le niveau atteint-il 60 cm ?]
Le niveau d'eau d'un réservoir est une fonction continue du temps $t$. À 8 h, le niveau est de 40 cm ; à 12 h, il est de 90 cm. Comme $40 < 60 < 90$ et que la fonction niveau est continue sur $[8\,; 12]$, le théorème des valeurs intermédiaires garantit qu'il existe (au moins) un instant $t_0$ entre 8 h et 12 h où le niveau vaut exactement 60 cm. On peut donc affirmer au gardien : « l'eau a atteint 60 cm dans la matinée », sans avoir observé l'instant précis.
\end{exemplebox}

\begin{attention}[Continuité physique, discontinuité tarifaire]
Toutes les grandeurs de la vie courante ne sont pas continues ! Les \emph{tarifs} sont souvent discontinus : le prix d'un billet de transport est le même pour 50 km et pour 99 km, puis saute d'un cran à 100 km (fonction « en escaliers »). Les mathématiques de cette leçon (TVI en particulier) ne s'appliquent qu'aux grandeurs \textbf{continues}. Toujours vérifier cette hypothèse avant de conclure.
\end{attention}

\courssub{Rédiger une démarche complète : modéliser, calculer, interpréter, conclure}

Au Probatoire et au Baccalauréat technique, un problème concret est noté autant sur la \emph{rédaction} que sur le calcul. Voici la démarche attendue.

\begin{methode}[Résoudre un problème concret avec les limites en 4 étapes]
\begin{enumerate}
\item \textbf{Modéliser} : identifier la grandeur étudiée, la variable, et écrire la fonction $f(x)$ qui les relie, \emph{en précisant les unités}.
\item \textbf{Calculer} la limite pertinente (en $+\infty$ pour un comportement à long terme, en une valeur $a$ pour un comportement critique), en appliquant les règles opératoires et en levant les formes indéterminées si nécessaire.
\item \textbf{Interpréter} le résultat \emph{dans le contexte} : traduire $\lim f = \ell$ par une phrase concrète (asymptote = plancher, plafond, valeur critique).
\item \textbf{Conclure} par une recommandation ou une décision argumentée (prix minimal, capacité maximale, danger à éviter).
\end{enumerate}
\end{methode}

\begin{attention}[Une limite sans unité est une réponse incomplète]
Dans l'interprétation, précise \textbf{toujours} les unités et le sens concret : écrire « la limite est 2000 » ne vaut rien ; il faut écrire « le coût moyen tend vers \num{2000} FCFA par sac : il ne descendra jamais sous ce seuil ». De même : 60 km/h, 600 litres, 35 °C. Un résultat numérique nu, sans unité ni phrase d'interprétation, perd des points à l'examen.
\end{attention}

\begin{exoresolu}[Tarification dans un atelier de soudure]
\textbf{Énoncé.} Un atelier de soudure à Bafoussam engage des coûts fixes de \num{80000} FCFA par mois. La fabrication de chaque grille métallique coûte \num{5000} FCFA en matière et main-d'œuvre. On note $x$ le nombre de grilles produites dans le mois.
\begin{enumerate}
\item Exprimer le coût moyen $C(x)$ d'une grille.
\item Calculer $\displaystyle\lim_{x \to +\infty} C(x)$ et interpréter.
\item L'atelier veut vendre chaque grille \num{5400} FCFA. À partir de combien de grilles le coût moyen devient-il inférieur à ce prix ?
\end{enumerate}
\tcblower
\textbf{Solution détaillée.}
\begin{enumerate}
\item \textbf{Modélisation.} Coût total : $CT(x) = \num{80000} + \num{5000}x$ (FCFA). Coût moyen, pour $x > 0$ :
\[ C(x) = \dfrac{\num{80000} + \num{5000}x}{x} = \num{5000} + \dfrac{\num{80000}}{x} \quad \text{(FCFA par grille)}. \]
\item \textbf{Calcul de la limite.} $\displaystyle\lim_{x \to +\infty} \dfrac{\num{80000}}{x} = 0$, donc par somme :
\[ \lim_{x \to +\infty} C(x) = \num{5000}. \]
\textbf{Interprétation.} La droite $y = \num{5000}$ est asymptote horizontale : le coût moyen d'une grille ne descendra \emph{jamais} sous \num{5000} FCFA, même en produisant énormément. C'est le \cle{coût moyen plancher}.
\item \textbf{Conclusion chiffrée.} On résout $C(x) < \num{5400}$ :
\begin{align*}
\num{5000} + \dfrac{\num{80000}}{x} &< \num{5400} \\
\dfrac{\num{80000}}{x} &< 400 \\
x &> \dfrac{\num{80000}}{400} = 200.
\end{align*}
\textbf{Conclusion.} Il faut produire et vendre \emph{plus de 200 grilles par mois} (au moins 201) pour que le coût moyen passe sous \num{5400} FCFA et que la vente à ce prix devienne rentable. En dessous de 200 grilles, l'atelier vend à perte.
\end{enumerate}
\end{exoresolu}

\begin{exercice}[Niveau d'eau et continuité]
Un réservoir se vide régulièrement. À 6 h, il contient 800 litres ; à 18 h, il n'en contient plus que 200 litres. On admet que la quantité d'eau est une fonction continue du temps.
\begin{enumerate}
\item Justifier qu'il existe un instant où le réservoir contient exactement 500 litres.
\item Le débit moyen de vidange est supposé constant. Le coût de pompage par heure est donné par $C(t) = \dfrac{\num{12000}}{t}$ FCFA pour une vidange étalée sur $t$ heures. Que vaut $\displaystyle\lim_{t \to +\infty} C(t)$ ? Interpréter.
\end{enumerate}
\end{exercice}

\begin{corrige}[Niveau d'eau et continuité]
\begin{enumerate}
\item La fonction $Q(t)$ (quantité d'eau) est continue sur $[6\,; 18]$, avec $Q(6) = 800$ et $Q(18) = 200$. Comme $200 < 500 < 800$, le théorème des valeurs intermédiaires garantit l'existence d'un instant $t_0 \in [6\,; 18]$ tel que $Q(t_0) = 500$ litres.
\item $\displaystyle\lim_{t \to +\infty} \dfrac{\num{12000}}{t} = 0$. Interprétation : plus on étale la vidange dans le temps, plus le coût horaire de pompage diminue, en tendant vers 0 FCFA par heure (le coût total se répartit sur de plus en plus d'heures).
\end{enumerate}
\end{corrige}

\courssub{Synthèse de la leçon}

\begin{savaistu}[Les limites sauvent des vies]
Les ingénieurs utilisent les limites chaque jour pour \cle{dimensionner} les ponts, les pylônes et les réseaux électriques : ils étudient le comportement « à la limite » — charge maximale, courant maximal, température extrême — pour s'assurer que la structure reste dans la zone de sécurité. Comprendre qu'une contrainte tend vers une valeur critique (asymptote verticale !) permet de fixer les marges de sécurité et d'éviter les ruptures. Au Cameroun, le dimensionnement des ouvrages d'art comme les ponts sur la Sanaga ou le Wouri repose sur ces calculs de comportements limites.
\end{savaistu}

\begin{popfigure}
\begin{center}
\begin{tikzpicture}[
mindmap,
grow cyclic,
every node/.style={concept, text=white, font=\small},
concept color=popblue,
level 1/.append style={level distance=3.6cm, sibling angle=90},
level 2/.append style={level distance=2.6cm, sibling angle=50},
]
\node[concept color=popdark, font=\small\bfseries] {Limites et continuité}
child[concept color=popblue] { node {Notion de limite}
  child { node {en un point $a$} }
  child { node {en $+\infty$ / $-\infty$} }
}
child[concept color=popgreen] { node {Opérations}
  child { node {somme, produit, quotient} }
  child { node {formes indéterminées} }
}
child[concept color=poporange] { node {Asymptotes}
  child { node {horizontale : plancher} }
  child { node {verticale : valeur critique} }
}
child[concept color=poppurple] { node {Continuité}
  child { node {pas de saut} }
  child { node {TVI : valeurs intermédiaires} }
};
\end{tikzpicture}
\end{center}
\legende{Carte mentale de synthèse : de la notion de limite aux applications concrètes, en passant par les opérations, les formes indéterminées, les asymptotes et la continuité.}
\end{popfigure}

\begin{aretenir}[L'essentiel de la section]
\begin{itemize}
\item Le \cle{coût moyen} $C(x) = \dfrac{F + vx}{x} = v + \dfrac{F}{x}$ tend vers $v$ quand $x \to +\infty$ : l'asymptote horizontale donne le \textbf{coût moyen plancher} et le \textbf{prix limite} de vente.
\item Une \cle{asymptote verticale} signale une \textbf{valeur critique} près de laquelle une grandeur explose (coût, temps, contrainte).
\item Température, niveau d'eau, distance parcourue sont des grandeurs \textbf{continues} : le TVI permet d'affirmer qu'une valeur intermédiaire est atteinte.
\item Toute démarche complète suit 4 étapes : \textbf{modéliser} $\to$ \textbf{calculer} la limite $\to$ \textbf{interpréter} avec les unités $\to$ \textbf{conclure} par une décision.
\end{itemize}
\end{aretenir}

\newpage

\coursec{Activité d'intégration}

\coursec{Leçon 21 : Limites et continuité}

\courssub{Objectifs de l'activité}

À l'issue de cette activité d'intégration, l'élève sera capable de :
\begin{itemize}
\item calculer des valeurs numériques d'une fonction et \cle{conjecturer} une limite à partir d'un tableau de valeurs ;
\item déterminer la \cle{limite} d'une fonction rationnelle en $+\infty$ et en $0^+$ par les méthodes du cours ;
\item interpréter graphiquement les limites en termes d'\cle{asymptotes} horizontale et verticale ;
\item justifier la \cle{continuité} d'une fonction sur un intervalle ;
\item résoudre une inéquation simple et rédiger une \cle{conclusion argumentée} destinée à un décideur (savoir-faire : rédiger et justifier une démarche, une réponse ou une conclusion).
\end{itemize}

\courssub{Situation}

La communauté villageoise de \textbf{Nkongsamba}, dans le département du Moungo (région du Littoral), fait face à des difficultés d'approvisionnement en eau potable, surtout pendant la saison sèche. La coopérative des planteurs de la localité décide de financer le \textbf{forage d'un puits} équipé d'une pompe.

Le bureau d'études fournit les données suivantes :
\begin{itemize}
\item le \textbf{coût fixe d'installation} (forage, pompe, réservoir, mise en service) s'élève à \num{1200000} FCFA, payé une seule fois ;
\item le \textbf{coût de fonctionnement} (énergie, entretien, réparations) est de \num{3000} FCFA par mètre cube d'eau pompé.
\end{itemize}

On note $x$ le volume total d'eau pompé, exprimé en mètres cubes ($x > 0$). Le coût total engagé est donc, en FCFA :
\[ \num{1200000} + \num{3000}\,x. \]

Le \textbf{coût moyen} du mètre cube d'eau, noté $C(x)$, est obtenu en divisant le coût total par le volume pompé :
\[ C(x) = \dfrac{\num{1200000} + \num{3000}\,x}{x} \qquad (x > 0). \]

Le président de la coopérative se pose deux questions :
\begin{itemize}
\item « En pompant beaucoup d'eau, jusqu'où peut descendre le coût moyen du mètre cube ? »
\item « À partir de quel volume pompé le coût moyen devient-il inférieur à \num{3500} FCFA ? »
\end{itemize}

Tu es membre du comité technique de la coopérative. Ta mission : analyser mathématiquement la situation et rédiger une note claire pour le président.

\courssub{Consignes}

\begin{enumerate}
\item \textbf{Approche numérique.} Calculer $C(x)$ pour $x = 100$ ; $500$ ; $\num{1000}$ ; $\num{5000}$ ; $\num{10000}$. Présenter les résultats dans un tableau. Que semble devenir $C(x)$ lorsque $x$ devient de plus en plus grand ? Formuler une conjecture sur la limite de $C$ en $+\infty$.
\item \textbf{Limite en $+\infty$.} En écrivant $C(x) = \dfrac{\num{1200000}}{x} + \num{3000}$, calculer $\displaystyle\lim_{x \to +\infty} C(x)$. En déduire l'équation de l'asymptote horizontale à la courbe de $C$ et interpréter ce résultat pour la coopérative.
\item \textbf{Limite en $0^+$.} Calculer $\displaystyle\lim_{x \to 0^+} C(x)$. Quelle asymptote cela révèle-t-il ? Expliquer concrètement pourquoi le coût moyen « explose » quand on pompe très peu d'eau.
\item \textbf{Continuité.} Justifier que la fonction $C$ est continue sur $]0\,;\,+\infty[$. Pourquoi le point $x = 0$ est-il exclu ?
\item \textbf{Décision.} Résoudre l'inéquation $C(x) < \num{3500}$ sur $]0\,;\,+\infty[$. Rédiger ensuite une \textbf{note d'une dizaine de lignes} au président de la coopérative, expliquant :
\begin{itemize}
\item pourquoi le coût moyen ne pourra jamais descendre sous \num{3000} FCFA par mètre cube ;
\item le volume minimal de pompage (arrondi au mètre cube supérieur) à atteindre pour que le coût moyen soit inférieur à \num{3500} FCFA.
\end{itemize}
\end{enumerate}

\courssub{Ressources à mobiliser}

\begin{aretenir}[Boîte à outils de la leçon]
\begin{itemize}
\item \textbf{Limites de référence :} $\displaystyle\lim_{x \to +\infty} \dfrac{1}{x} = 0$ et $\displaystyle\lim_{x \to 0^+} \dfrac{1}{x} = +\infty$.
\item \textbf{Opérations :} la limite d'une somme, d'un produit ou d'un quotient se calcule terme à terme (sauf formes indéterminées).
\item \textbf{Fonction rationnelle en $+\infty$ :} la limite est celle du quotient des termes de plus haut degré.
\item \textbf{Asymptotes :} si $\displaystyle\lim_{x \to +\infty} f(x) = b$, la droite $y = b$ est \cle{asymptote horizontale} ; si $\displaystyle\lim_{x \to a^+} f(x) = +\infty$, la droite $x = a$ est \cle{asymptote verticale}.
\item \textbf{Continuité :} une fonction rationnelle est continue sur tout intervalle de son ensemble de définition.
\end{itemize}
\end{aretenir}

\courssub{Démarche et corrigé commenté}

\begin{exoresolu}[Le forage de Nkongsamba : corrigé complet]
\textbf{Question 1.} Calculer $C(x)$ pour $x \in \{100\,;\,500\,;\,\num{1000}\,;\,\num{5000}\,;\,\num{10000}\}$ et conjecturer la limite en $+\infty$.

\textbf{Question 2.} Calculer $\displaystyle\lim_{x \to +\infty} C(x)$ et identifier l'asymptote horizontale.

\textbf{Question 3.} Calculer $\displaystyle\lim_{x \to 0^+} C(x)$ et interpréter.

\textbf{Question 4.} Justifier la continuité de $C$ sur $]0\,;\,+\infty[$.

\textbf{Question 5.} Résoudre $C(x) < \num{3500}$ et rédiger la note au président.
\tcblower
\textbf{Question 1 : approche numérique.}

On applique la formule $C(x) = \dfrac{\num{1200000} + \num{3000}\,x}{x}$ :
\begin{align*}
C(100) &= \dfrac{\num{1200000} + \num{300000}}{100} = \num{15000} \text{ FCFA} \\
C(500) &= \dfrac{\num{1200000} + \num{1500000}}{500} = \num{5400} \text{ FCFA} \\
C(\num{1000}) &= \dfrac{\num{1200000} + \num{3000000}}{\num{1000}} = \num{4200} \text{ FCFA} \\
C(\num{5000}) &= \dfrac{\num{1200000} + \num{15000000}}{\num{5000}} = \num{3240} \text{ FCFA} \\
C(\num{10000}) &= \dfrac{\num{1200000} + \num{30000000}}{\num{10000}} = \num{3120} \text{ FCFA}
\end{align*}

\begin{center}
\begin{tabularx}{\linewidth}{|l|X|X|X|X|X|}
\hline
\rowcolor{popblueL} $x$ (m$^3$) & $100$ & $500$ & $\num{1000}$ & $\num{5000}$ & $\num{10000}$ \\
\hline $C(x)$ (FCFA) & $\num{15000}$ & $\num{5400}$ & $\num{4200}$ & $\num{3240}$ & $\num{3120}$ \\
\hline
\end{tabularx}
\end{center}

\textbf{Conjecture :} lorsque $x$ devient très grand, $C(x)$ diminue et semble se rapprocher de $\num{3000}$ FCFA, sans jamais l'atteindre.

\textbf{Question 2 : limite en $+\infty$.}

On réécrit la fonction en séparant les termes :
\[ C(x) = \dfrac{\num{1200000}}{x} + \num{3000}. \]
Or $\displaystyle\lim_{x \to +\infty} \dfrac{\num{1200000}}{x} = 0$ (limite de référence : une constante divisée par un nombre qui tend vers $+\infty$ tend vers $0$). Par somme de limites :
\[ \lim_{x \to +\infty} C(x) = 0 + \num{3000} = \num{3000}. \]
\textit{(Vérification par la méthode des termes de plus haut degré : $C(x) = \dfrac{\num{3000}\,x + \num{1200000}}{x}$, et $\dfrac{\num{3000}\,x}{x} = \num{3000}$.)}

\textbf{Interprétation graphique :} la droite d'équation $y = \num{3000}$ est \cle{asymptote horizontale} à la courbe de $C$ en $+\infty$.

\textbf{Interprétation concrète :} plus la coopérative pompe d'eau, plus le coût fixe de \num{1200000} FCFA est « dilué » sur un grand nombre de mètres cubes ; le coût moyen se rapproche alors du seul coût de fonctionnement, soit \num{3000} FCFA/m$^3$, mais ne descend jamais en dessous.

\textbf{Question 3 : limite en $0^+$.}

On utilise la même écriture : $C(x) = \dfrac{\num{1200000}}{x} + \num{3000}$.

Quand $x \to 0^+$, on a $\dfrac{\num{1200000}}{x} \to +\infty$ (limite de référence : $\displaystyle\lim_{x \to 0^+} \dfrac{1}{x} = +\infty$). Donc :
\[ \lim_{x \to 0^+} C(x) = +\infty. \]

\textbf{Interprétation graphique :} la droite d'équation $x = 0$ (l'axe des ordonnées) est \cle{asymptote verticale} à la courbe.

\textbf{Interprétation concrète :} si l'on pompe très peu d'eau, l'investissement de \num{1200000} FCFA se répartit sur un tout petit nombre de mètres cubes : chaque mètre cube coûte alors extrêmement cher (par exemple \num{15000} FCFA/m$^3$ pour seulement 100 m$^3$ pompés).

\textbf{Question 4 : continuité.}

La fonction $C$ est une \cle{fonction rationnelle} (quotient de deux fonctions polynômes : $x \mapsto \num{1200000} + \num{3000}\,x$ et $x \mapsto x$). Elle est donc définie et continue partout où son dénominateur ne s'annule pas, c'est-à-dire sur $\mathbb{R} \setminus \{0\}$, et en particulier sur $]0\,;\,+\infty[$.

Le point $x = 0$ est exclu car la division par zéro est impossible — et d'ailleurs, concrètement, un « coût moyen pour zéro mètre cube pompé » n'a aucun sens. La fonction $C$ est donc \textbf{continue sur $]0\,;\,+\infty[$}.

\textbf{Question 5 : résolution de l'inéquation et note au président.}

On résout sur $]0\,;\,+\infty[$ :
\begin{align*}
C(x) < \num{3500} &\iff \dfrac{\num{1200000}}{x} + \num{3000} < \num{3500} \\
&\iff \dfrac{\num{1200000}}{x} < 500 \\
&\iff \num{1200000} < 500\,x \quad (\text{car } x > 0,\ \text{le sens ne change pas}) \\
&\iff x > \dfrac{\num{1200000}}{500} = \num{2400}.
\end{align*}
L'ensemble des solutions est $S = ]\num{2400}\,;\,+\infty[$. Le volume minimal entier est donc $\num{2401}$ m$^3$.

\textbf{Note au président de la coopérative (modèle de rédaction) :}

\begin{quote}
Monsieur le Président,

L'étude mathématique du coût moyen $C(x) = \dfrac{\num{1200000} + \num{3000}\,x}{x}$ conduit aux conclusions suivantes. Le coût moyen du mètre cube diminue au fur et à mesure que le volume pompé augmente, et il se rapproche de \num{3000} FCFA sans jamais l'atteindre : en effet, $\displaystyle\lim_{x \to +\infty} C(x) = \num{3000}$, et pour tout $x > 0$, $C(x) = \num{3000} + \dfrac{\num{1200000}}{x} > \num{3000}$. Il est donc impossible de descendre sous \num{3000} FCFA/m$^3$, car la part du coût fixe, $\dfrac{\num{1200000}}{x}$, reste toujours strictement positive. En revanche, dès que le volume total pompé dépasse \num{2400} m$^3$, le coût moyen devient inférieur à \num{3500} FCFA/m$^3$. Je recommande donc de planifier un pompage d'au moins \num{2401} m$^3$ sur la durée de vie du forage, volume largement atteignable au vu des besoins de la communauté. Le projet est rentable à moyen terme.
\end{quote}
\end{exoresolu}

\begin{attention}[Pièges fréquents dans cette activité]
\begin{itemize}
\item Multiplier une inéquation par $x$ sans vérifier son signe : ici $x > 0$, donc le sens de l'inégalité est conservé. Le préciser dans la rédaction est exigé.
\item Conclure que le coût moyen « atteint » \num{3000} FCFA : non, il s'en \emph{approche}. Une asymptote n'est (en général) jamais atteinte.
\item Oublier que $C(0)$ n'existe pas : on ne peut pas « prolonger » $C$ en $0$, car la limite y est infinie.
\end{itemize}
\end{attention}

\courssub{Bilan de l'activité}

Cette activité a permis de mettre en évidence les points essentiels de la leçon :
\begin{itemize}
\item la \textbf{conjecture numérique} (tableau de valeurs) donne une intuition de la limite, que le calcul rigoureux vient ensuite confirmer ;
\item la \textbf{limite en $+\infty$} d'une fonction rationnelle se lit sur les termes de plus haut degré et se traduit graphiquement par une \textbf{asymptote horizontale}, ici $y = \num{3000}$ ;
\item la \textbf{limite infinie en $0^+$} se traduit par une \textbf{asymptote verticale} et s'interprète concrètement : répartir un coût fixe sur un très petit volume donne un coût unitaire énorme ;
\item la \textbf{continuité} sur $]0\,;\,+\infty[$ garantit que le coût moyen évolue « sans saut » : toute valeur intermédiaire entre \num{3000} et les grandes valeurs est atteinte pour un certain volume ;
\item enfin, les mathématiques fournissent un \textbf{outil d'aide à la décision} : l'inéquation $C(x) < \num{3500}$ donne un seuil précis ($\num{2400}$ m$^3$) que le gestionnaire peut utiliser pour planifier l'exploitation du forage.
\end{itemize}

\begin{savaistu}[L'eau potable, un enjeu national]
Selon les programmes nationaux d'hydraulique villageoise, des milliers de forages sont réalisés chaque année au Cameroun, notamment dans les régions du Littoral, du Centre et de l'Extrême-Nord. L'analyse du coût moyen que tu viens d'effectuer est exactement celle que mènent les bureaux d'études et les ONG pour décider de la viabilité d'un point d'eau : un forage n'est « rentable » que si le volume pompé sur sa durée de vie est suffisamment grand pour amortir l'investissement initial. Les limites de fonctions sont ainsi un outil concret de gestion des projets de développement.
\end{savaistu}

\newpage

\coursec{Méthodes et exercices résolus}

\coursec{Limites et continuité}

\courssub{Approche intuitive de la notion de limite}

\begin{definition}[Limite finie en un point]
Soit $f$ une fonction définie sur un intervalle pointé $I\setminus\{a\}$ (autour de $a$). On dit que $f$ admet la limite $\ell\in\mathbb{R}$ en $a$ si, lorsque $x$ se rapproche de $a$ (sans être égal à $a$), $f(x)$ se rapproche arbitrairement de $\ell$. On note :
\[\lim_{x\to a}f(x)=\ell.\]
\end{definition}

\begin{definition}[Limite infinie en un point]
On dit que $f$ admet pour limite $+\infty$ (resp. $-\infty$) en $a$ si, lorsque $x\to a$, $f(x)$ devient arbitrairement grand (resp. petit). On note :
\[\lim_{x\to a}f(x)=+\infty \quad\text{ou}\quad \lim_{x\to a}f(x)=-\infty.\]
\end{definition}

\begin{definition}[Limite en $+\infty$ (ou $-\infty$)]
On dit que $f$ admet la limite $\ell$ (finie ou infinie) en $+\infty$ si, lorsque $x$ devient arbitrairement grand, $f(x)$ se rapproche de $\ell$. On note :
\[\lim_{x\to+\infty}f(x)=\ell.\]
\end{definition}

\begin{propriete}[Asymptotes]
\begin{itemize}
\item \textbf{Asymptote verticale :} si $\displaystyle\lim_{x\to a^+}f(x)=\pm\infty$ ou $\displaystyle\lim_{x\to a^-}f(x)=\pm\infty$, la droite d'équation $x=a$ est une asymptote verticale à la courbe $\mathcal{C}_f$.
\item \textbf{Asymptote horizontale :} si $\displaystyle\lim_{x\to+\infty}f(x)=\ell\in\mathbb{R}$ (ou pour $x\to-\infty$), la droite d'équation $y=\ell$ est une asymptote horizontale à $\mathcal{C}_f$.
\end{itemize}
\end{propriete}

\begin{exemplebox}[Interprétation graphique : $f(x)=\dfrac{1}{x}$]
\begin{itemize}
\item $\displaystyle\lim_{x\to0^+}f(x)=+\infty$ : la courbe monte vers le haut près de l'axe $x=0$ (à droite).
\item $\displaystyle\lim_{x\to0^-}f(x)=-\infty$ : la courbe descend vers le bas près de l'axe $x=0$ (à gauche).
\item $\displaystyle\lim_{x\to+\infty}f(x)=0$ : la courbe se colle à l'axe des abscisses ($y=0$) vers la droite.
\item $\displaystyle\lim_{x\to-\infty}f(x)=0$ : la courbe se colle à l'axe des abscisses vers la gauche.
\end{itemize}
La droite $x=0$ (axe des ordonnées) est asymptote verticale ; la droite $y=0$ (axe des abscisses) est asymptote horizontale.
\end{exemplebox}

\courssub{Limites de référence et interprétation graphique}

\begin{aretenir}[Limites de référence à connaître par cœur]
\begin{center}
\begin{tabular}{|c|c|c|}\hline
Fonction $f(x)$ & $\displaystyle\lim_{x\to+\infty}f(x)$ & $\displaystyle\lim_{x\to0^+}f(x)$ \\ \hline
$x^n\ (n\in\mathbb{N}^*)$ & $+\infty$ & $0$ \\ \hline
$\dfrac{1}{x}$ & $0$ & $+\infty$ \\ \hline
$\dfrac{1}{x^n}\ (n\in\mathbb{N}^*)$ & $0$ & $+\infty$ \\ \hline
$\sqrt{x}$ & $+\infty$ & $0$ \\ \hline
\end{tabular}
\end{center}
\textbf{Autres repères utiles :} $\displaystyle\lim_{x\to a}k=k$ (constante), $\displaystyle\lim_{x\to a}x=a$.
\end{aretenir}

\begin{methode}[Calculer une limite simple en un point ou à l'infini (lecture directe)]
\begin{enumerate}
\item \textbf{Remplacer} directement $x$ par la valeur $a$ (ou raisonner sur le comportement quand $x$ devient très grand).
\item Si le résultat est un nombre fini, c'est la limite : on écrit $\displaystyle\lim_{x\to a}f(x)=\ell$.
\item Si le dénominateur s'annule mais pas le numérateur, \textbf{étudier le signe} du dénominateur (à gauche et à droite de $a$) : la limite est $+\infty$ ou $-\infty$.
\item Pour les fonctions de référence, appliquer les résultats connus : $\displaystyle\lim_{x\to+\infty}x^n=+\infty$ ($n\in\mathbb{N}^*$), $\displaystyle\lim_{x\to+\infty}\frac{1}{x}=0$, $\displaystyle\lim_{x\to0^+}\frac{1}{x^2}=+\infty$.
\item Conclure par une phrase : \og la droite d'équation $x=a$ est une asymptote verticale \fg{} ou \og la droite d'équation $y=\ell$ est une asymptote horizontale \fg{} si cela est demandé.
\end{enumerate}
\end{methode}

\begin{exoresolu}[Limite en un point et interprétation graphique]
Soit la fonction $f$ définie par $f(x)=\dfrac{2x+1}{x-3}$.
\begin{enumerate}
\item Déterminer l'ensemble de définition de $f$.
\item Calculer $\displaystyle\lim_{x\to3^+}f(x)$ et $\displaystyle\lim_{x\to3^-}f(x)$.
\item En déduire une asymptote à la courbe de $f$.
\end{enumerate}
\tcblower
\textbf{1. Ensemble de définition.} $f$ est définie si et seulement si $x-3\neq0$, c'est-à-dire $x\neq3$. Donc $D_f=\mathbb{R}\setminus\{3\}$.

\textbf{2. Limites en $3$.} Quand $x$ tend vers $3$, le numérateur tend vers $2\times3+1=7$ (nombre strictement positif) et le dénominateur tend vers $0$. On étudie le signe du dénominateur :
\begin{itemize}
\item Si $x\to3^+$ (c'est-à-dire $x>3$), alors $x-3>0$, donc le quotient est positif et très grand : $\displaystyle\lim_{x\to3^+}f(x)=+\infty$.
\item Si $x\to3^-$ (c'est-à-dire $x<3$), alors $x-3<0$, donc le quotient est négatif et très grand en valeur absolue : $\displaystyle\lim_{x\to3^-}f(x)=-\infty$.
\end{itemize}

\textbf{3. Conclusion.} Comme les limites en $3$ sont infinies, la droite d'équation $x=3$ est une \cle{asymptote verticale} à la courbe représentative de $f$.
\end{exoresolu}

\courssub{Opérations sur les limites}

\begin{propriete}[Théorèmes sur les opérations (limites finies ou infinies)]
Soient $f$ et $g$ deux fonctions et $a\in\mathbb{R}\cup\{+\infty,-\infty\}$.
\begin{itemize}
\item \textbf{Somme :} Si $\lim f = \ell_1$ et $\lim g = \ell_2$ (pas de forme $\infty-\infty$), alors $\lim(f+g) = \ell_1+\ell_2$.
\item \textbf{Produit :} Si $\lim f = \ell_1$ et $\lim g = \ell_2$ (pas de forme $0\times\infty$), alors $\lim(fg) = \ell_1\ell_2$.
\item \textbf{Quotient :} Si $\lim f = \ell_1$ et $\lim g = \ell_2\neq0$ (pas de forme $\frac{\infty}{\infty}$ ni $\frac{0}{0}$), alors $\lim\frac{f}{g} = \frac{\ell_1}{\ell_2}$.
\end{itemize}
\textbf{Règles usuelles avec l'infini :} $\infty+\infty=\infty$ ; $\infty\times\infty=\infty$ ; $\infty\times(-\infty)=-\infty$ ; $\frac{k}{\infty}=0$ ($k\in\mathbb{R}$) ; $\frac{\infty}{k}=\infty$ ($k>0$).
\end{propriete}

\begin{methode}[Utiliser les théorèmes sur les opérations]
\begin{enumerate}
\item \textbf{Identifier} la limite de chaque terme ou facteur séparément.
\item \textbf{Appliquer} le théorème correspondant : somme, produit, quotient.
\item \textbf{Vérifier} que l'on n'est pas en présence d'une forme indéterminée : $+\infty-\infty$, $0\times\infty$, $\dfrac{\infty}{\infty}$, $\dfrac{0}{0}$.
\item Si aucune indétermination : conclure directement en rédigeant \og par somme de limites\fg{} ou \og par produit de limites\fg{}.
\end{enumerate}
\end{methode}

\begin{exoresolu}[Limites par opérations]
Calculer les limites suivantes en justifiant :
\begin{enumerate}
\item $\displaystyle\lim_{x\to+\infty}\left(x^2+3x-5\right)$
\item $\displaystyle\lim_{x\to+\infty}\left(2+\frac{1}{x}\right)\left(x+1\right)$
\end{enumerate}
\tcblower
\textbf{1.} On a $\displaystyle\lim_{x\to+\infty}x^2=+\infty$, $\displaystyle\lim_{x\to+\infty}3x=+\infty$ et $\displaystyle\lim_{x\to+\infty}(-5)=-5$.

Par somme de limites ($+\infty$ plus un nombre fini donne $+\infty$) :
\[\lim_{x\to+\infty}\left(x^2+3x-5\right)=+\infty.\]

\textbf{2.} D'une part $\displaystyle\lim_{x\to+\infty}\frac{1}{x}=0$, donc par somme : $\displaystyle\lim_{x\to+\infty}\left(2+\frac{1}{x}\right)=2$.

D'autre part $\displaystyle\lim_{x\to+\infty}(x+1)=+\infty$.

Par produit de limites (un nombre strictement positif multiplié par $+\infty$ donne $+\infty$) :
\[\lim_{x\to+\infty}\left(2+\frac{1}{x}\right)(x+1)=+\infty.\]
\end{exoresolu}

\courssub{Calculs pratiques : lever les formes indéterminées}

\begin{aretenir}[Formes indéterminées principales]
\begin{center}
\begin{tabular}{|c|c|}\hline
Forme & Signification \\ \hline
$\dfrac{\infty}{\infty}$ & Quotient de deux grandeurs infinies \\ \hline
$\dfrac{0}{0}$ & Quotient de deux grandeurs tendant vers $0$ \\ \hline
$+\infty-\infty$ & Différence de deux grandeurs infinies \\ \hline
$0\times\infty$ & Produit d'une grandeur tendant vers $0$ par une infinie \\ \hline
\end{tabular}
\end{center}
\emph{Attention :} Ces formes ne donnent \textbf{pas} de résultat direct ; il faut transformer l'expression.
\end{aretenir}

\begin{methode}[Lever les formes indéterminées classiques]
\begin{enumerate}
\item \textbf{Constater} l'indétermination en écrivant explicitement \og on obtient la forme $\dfrac{\infty}{\infty}$ \fg{} (ou $0/0$, $+\infty-\infty$, $0\times\infty$).
\item \textbf{Polynôme ou fraction rationnelle en $+\infty$ ou $-\infty$} : mettre en facteur le terme de \cle{plus haut degré} au numérateur et au dénominateur, puis simplifier.
\item \textbf{Forme $0/0$ en un point $a$} : factoriser par $(x-a)$ au numérateur et au dénominateur (souvent avec une identité remarquable), puis simplifier \textbf{en précisant $x\neq a$}.
\item \textbf{Conclure} en calculant la limite de l'expression simplifiée.
\end{enumerate}
\end{methode}

\begin{exoresolu}[Forme indéterminée $\infty/\infty$]
Calculer $\displaystyle\lim_{x\to+\infty}\frac{3x^2-5x+1}{2x^2+4}$.
\tcblower
Quand $x\to+\infty$, le numérateur et le dénominateur tendent tous les deux vers $+\infty$ : on obtient la forme indéterminée $\dfrac{\infty}{\infty}$.

On met en facteur le terme de plus haut degré, $x^2$, au numérateur et au dénominateur :
\[\frac{3x^2-5x+1}{2x^2+4}=\frac{x^2\left(3-\dfrac{5}{x}+\dfrac{1}{x^2}\right)}{x^2\left(2+\dfrac{4}{x^2}\right)}=\frac{3-\dfrac{5}{x}+\dfrac{1}{x^2}}{2+\dfrac{4}{x^2}} \quad (x\neq0).\]

Or $\displaystyle\lim_{x\to+\infty}\frac{5}{x}=0$, $\displaystyle\lim_{x\to+\infty}\frac{1}{x^2}=0$ et $\displaystyle\lim_{x\to+\infty}\frac{4}{x^2}=0$.

Par somme puis quotient de limites :
\[\lim_{x\to+\infty}\frac{3x^2-5x+1}{2x^2+4}=\frac{3}{2}.\]

\textbf{Conclusion :} $\displaystyle\lim_{x\to+\infty}\frac{3x^2-5x+1}{2x^2+4}=\dfrac{3}{2}$. La droite d'équation $y=\dfrac{3}{2}$ est asymptote horizontale à la courbe.
\end{exoresolu}

\begin{exoresolu}[Forme indéterminée $0/0$]
Calculer $\displaystyle\lim_{x\to2}\frac{x^2-4}{x-2}$.
\tcblower
En remplaçant $x$ par $2$ : le numérateur vaut $2^2-4=0$ et le dénominateur vaut $2-2=0$. On obtient la forme indéterminée $\dfrac{0}{0}$.

On factorise le numérateur grâce à l'identité remarquable $a^2-b^2=(a-b)(a+b)$ :
\[x^2-4=(x-2)(x+2).\]

Donc, \textbf{pour $x\neq2$} :
\[\frac{x^2-4}{x-2}=\frac{(x-2)(x+2)}{x-2}=x+2.\]

Alors :
\[\lim_{x\to2}\frac{x^2-4}{x-2}=\lim_{x\to2}(x+2)=2+2=4.\]

\textbf{Conclusion :} $\displaystyle\lim_{x\to2}\dfrac{x^2-4}{x-2}=4$.
\end{exoresolu}

\courssub{Continuité d'une fonction}

\begin{definition}[Continuité en un point]
Une fonction $f$ est \cle{continue} en $a$ si et seulement si les trois conditions suivantes sont réunies :
\begin{enumerate}
\item $f$ est définie en $a$ ($a\in D_f$) ;
\item $\displaystyle\lim_{x\to a}f(x)$ existe (finie) ;
\item $\displaystyle\lim_{x\to a}f(x)=f(a)$.
\end{enumerate}
Si l'une de ces conditions n'est pas vérifiée, $f$ est \cle{discontinue} en $a$.
\end{definition}

\begin{propriete}[Fonctions usuelles continues]
\begin{itemize}
\item Les fonctions \cle{polynômes} sont continues sur $\mathbb{R}$.
\item Les fonctions \cle{rationnelles} sont continues sur chaque intervalle de leur ensemble de définition.
\item La fonction racine carrée $x\mapsto\sqrt{x}$ est continue sur $[0,+\infty[$.
\end{itemize}
\end{propriete}

\begin{methode}[Justifier la continuité en un point ou sur un intervalle]
\begin{enumerate}
\item \textbf{Vérifier} que la fonction est définie au point $a$ : $f(a)$ existe.
\item \textbf{Calculer} $\displaystyle\lim_{x\to a}f(x)$ (éventuellement limites à gauche et à droite).
\item \textbf{Comparer} : $f$ est continue en $a$ si et seulement si $\displaystyle\lim_{x\to a}f(x)=f(a)$.
\item Pour une fonction définie par morceaux, vérifier la continuité au point de raccordement en comparant les deux expressions.
\item Sur un intervalle : $f$ est continue sur $I$ si elle est continue en tout point de $I$.
\end{enumerate}
\end{methode}

\begin{exoresolu}[Continuité d'une fonction définie par morceaux]
Soit la fonction $g$ définie sur $\mathbb{R}$ par :
\[g(x)=\begin{cases} 2x+1 & \text{si } x\leq 1 \\ x^2+2 & \text{si } x>1 \end{cases}\]
La fonction $g$ est-elle continue en $1$ ?
\tcblower
\textbf{Étape 1 :} $g$ est définie en $1$ (première expression) : $g(1)=2\times1+1=3$.

\textbf{Étape 2 :} calcul des limites à gauche et à droite de $1$.
\begin{itemize}
\item À gauche : $\displaystyle\lim_{x\to1^-}g(x)=\lim_{x\to1^-}(2x+1)=3$.
\item À droite : $\displaystyle\lim_{x\to1^+}g(x)=\lim_{x\to1^+}(x^2+2)=1^2+2=3$.
\end{itemize}

\textbf{Étape 3 :} les deux limites sont égales et valent $g(1)=3$ :
\[\lim_{x\to1^-}g(x)=\lim_{x\to1^+}g(x)=g(1)=3.\]

\textbf{Conclusion :} $\displaystyle\lim_{x\to1}g(x)=g(1)$, donc la fonction $g$ est \cle{continue} en $1$.
\end{exoresolu}

\courssub{Applications à des situations concrètes}

\begin{methode}[Résoudre un problème concret avec les limites]
\begin{enumerate}
\item \textbf{Traduire} la situation : identifier la grandeur, la variable et la fonction qui la modélise.
\item \textbf{Préciser} le domaine de validité du modèle (valeurs possibles de la variable dans le contexte).
\item \textbf{Calculer} la limite demandée en appliquant les méthodes précédentes.
\item \textbf{Interpréter} le résultat dans le langage de la situation (coût, température, stock, tension...) avec une phrase complète et les unités.
\end{enumerate}
\end{methode}

\begin{exoresolu}[Coût moyen dans un atelier de Douala]
Un atelier de fabrication de meubles à Douala produit $x$ tables par semaine. Le coût moyen de fabrication d'une table, en milliers de francs CFA, est donné par :
\[C(x)=\frac{50x+200}{x}\quad\text{pour } x\geq 1.\]
\begin{enumerate}
\item Calculer le coût moyen pour $x=10$ tables.
\item Calculer $\displaystyle\lim_{x\to+\infty}C(x)$ et interpréter le résultat pour l'atelier.
\end{enumerate}
\tcblower
\textbf{1.} Pour $x=10$ :
\[C(10)=\frac{50\times10+200}{10}=\frac{700}{10}=70.\]
Le coût moyen de fabrication est de $70$ milliers de francs CFA, soit $\SI{70000}{F}$ par table, lorsque l'atelier produit $10$ tables par semaine.

\textbf{2.} Quand $x\to+\infty$, le numérateur et le dénominateur tendent vers $+\infty$ : forme indéterminée $\dfrac{\infty}{\infty}$. On factorise par $x$ au numérateur :
\[C(x)=\frac{x\left(50+\dfrac{200}{x}\right)}{x}=50+\frac{200}{x} \quad (x>0).\]

Or $\displaystyle\lim_{x\to+\infty}\frac{200}{x}=0$, donc par somme de limites :
\[\lim_{x\to+\infty}C(x)=50.\]

\textbf{Interprétation :} lorsque la production devient très grande, le coût moyen de fabrication d'une table se rapproche de $50$ milliers de francs CFA ($\SI{50000}{F}$), sans jamais descendre en dessous. Les frais fixes de $200$ milliers de francs se répartissent sur un grand nombre de tables et deviennent négligeables par table.
\end{exoresolu}

\begin{attention}[Les 5 erreurs qui coûtent des points au Probatoire]
\begin{enumerate}
\item \textbf{Écrire $\dfrac{1}{0}=\infty$ sans préciser le signe.} Quand le dénominateur tend vers $0$, il faut étudier son \cle{signe} à gauche et à droite : $\displaystyle\lim_{x\to0^+}\frac{1}{x}=+\infty$ mais $\displaystyle\lim_{x\to0^-}\frac{1}{x}=-\infty$.
\item \textbf{Conclure directement face à une forme indéterminée.} Écrire \og $\infty-\infty=0$ \fg{} ou \og $\dfrac{\infty}{\infty}=1$ \fg{} est faux et sanctionné. Il faut annoncer la forme indéterminée puis factoriser (terme de plus haut degré, ou facteur $(x-a)$).
\item \textbf{Simplifier par $(x-a)$ sans préciser $x\neq a$.} La simplification $\dfrac{(x-2)(x+2)}{x-2}=x+2$ n'est valable que pour $x\neq2$ : il faut l'écrire.
\item \textbf{Confondre limite et valeur.} Dire qu'une fonction est continue en $a$ exige de vérifier les \textbf{trois} conditions : $f(a)$ existe, la limite en $a$ existe, et les deux sont égales. Une seule vérification ne suffit pas.
\item \textbf{Oublier la conclusion et l'interprétation.} Au Probatoire, toute limite infinie en un point doit être suivie de la conclusion \og asymptote verticale d'équation $x=a$ \fg{} si la question le demande, et tout problème concret exige une phrase d'interprétation avec les unités.
\end{enumerate}
\end{attention}

\newpage

\coursec{Exercices}

\courssub{Je vérifie mes connaissances}

\begin{exercice}[QCM : Limites à l'infini]
Parmi les affirmations suivantes, une seule est exacte. Laquelle ? Justifier brièvement.
\begin{enumerate}
\item $\displaystyle\lim_{x\to +\infty} (2x^3 - 5x + 1) = -\infty$
\item $\displaystyle\lim_{x\to +\infty} \frac{3x^2 + x}{x^2 - 4} = 3$
\item $\displaystyle\lim_{x\to +\infty} \frac{x+2}{2x-3} = +\infty$
\item $\displaystyle\lim_{x\to -\infty} (x^2 - 3x) = -\infty$
\end{enumerate}
\end{exercice}

\begin{exercice}[Vrai ou Faux : Opérations sur les limites]
Dire si les affirmations suivantes sont vraies ou fausses. Justifier chaque réponse par une règle de calcul ou un contre-exemple.
\begin{enumerate}
\item Si $\lim_{x\to a} f(x) = +\infty$ et $\lim_{x\to a} g(x) = 0$, alors $\lim_{x\to a} f(x)g(x) = 0$.
\item $\displaystyle\lim_{x\to +\infty} \frac{5x^2 - 2}{2x^3 + x} = 0$.
\item Si $f$ est continue sur $\mathbb{R}$ et $\lim_{x\to +\infty} f(x) = 5$, alors $f$ est bornée sur $\mathbb{R}$.
\item La fonction $x \mapsto \frac{1}{x}$ est continue sur $\mathbb{R}$.
\end{enumerate}
\end{exercice}

\begin{exercice}[Définitions et vocabulaire]
Compléter les phrases suivantes en utilisant le vocabulaire précis du cours.
\begin{enumerate}
\item On dit qu'une fonction $f$ admet pour limite $+\infty$ en $+\infty$ si \emph{...}
\item Une fonction $f$ est continue en un point $a$ de son ensemble de définition si \emph{...}
\item La droite d'équation $y = ax + b$ est une asymptote oblique à la courbe $\mathcal{C}_f$ en $+\infty$ si \emph{...}
\item On appelle \emph{forme indéterminée} une expression du type \emph{...} qui ne permet pas de conclure directement sur la valeur de la limite.
\end{enumerate}
\end{exercice}

\begin{exercice}[Calculs directs de limites]
Calculer les limites suivantes en justifiant chaque étape par les théorèmes de opérations ou les limites de référence.
\begin{enumerate}
\item $\displaystyle\lim_{x\to +\infty} (3x^4 - 2x^2 + 7)$
\item $\displaystyle\lim_{x\to -\infty} \frac{4x^3 - x}{2x^3 + 5}$
\item $\displaystyle\lim_{x\to 2} \frac{x^2 - 4}{x - 2}$
\item $\displaystyle\lim_{x\to 0} \frac{\sin(3x)}{x}$ \quad (On admet $\lim_{u\to 0} \frac{\sin u}{u} = 1$)
\end{enumerate}
\end{exercice}

\courssub{Je m'entraîne}

\begin{exercice}[Limites de fonctions rationnelles à l'infini]
Soit la fonction $f(x) = \frac{2x^3 - 5x + 1}{x^2 + 3}$.
\begin{enumerate}
\item Déterminer $\lim_{x\to +\infty} f(x)$ et $\lim_{x\to -\infty} f(x)$.
\item En déduire l'existence et l'équation d'une éventuelle asymptote oblique à la courbe $\mathcal{C}_f$ en $+\infty$ et en $-\infty$.
\item Construire le tableau de comportement de $f$ aux bornes de son ensemble de définition.
\end{enumerate}
\end{exercice}

\begin{exercice}[Levée d'indétermination et prolongement par continuité]
Soit la fonction $g$ définie sur $\mathbb{R} \setminus \{2\}$ par $g(x) = \frac{x^2 - 5x + 6}{x - 2}$.
\begin{enumerate}
\item Factoriser le numérateur et simplifier l'expression de $g(x)$ pour $x \neq 2$.
\item Calculer $\lim_{x\to 2} g(x)$.
\item La fonction $g$ peut-elle être prolongée par continuité en $2$ ? Si oui, donner la fonction prolongée $\tilde{g}$.
\item Tracer schématiquement la courbe de $g$ au voisinage de $x=2$ en indiquant le « trou ».
\end{enumerate}
\end{exercice}

\begin{exercice}[Asymptotes et tableau de comportement]
On considère la fonction $h$ définie sur $\mathbb{R} \setminus \{-2\}$ par $h(x) = \frac{4x - 1}{x + 2}$.
\begin{enumerate}
\item Calculer $\lim_{x\to -2^-} h(x)$ et $\lim_{x\to -2^+} h(x)$. En déduire l'asymptote verticale.
\item Calculer $\lim_{x\to \pm\infty} h(x)$. En déduire l'asymptote horizontale.
\item Construire le tableau de comportement complet de $h$ (variations non demandées, seulement limites aux bornes et aux valeurs interdites).
\item Placer les asymptotes et tracer schématiquement la courbe $\mathcal{C}_h$.
\end{enumerate}
\end{exercice}

\begin{exercice}[Continuité et définition]
On donne ci-dessous la représentation graphique d'une fonction $k$ définie par morceaux.
\begin{center}
\begin{tikzpicture}[scale=0.8, >=Stealth]
\draw[popline] (-3,0) -- (4,0) node[right] {$x$};
\draw[popline] (0,-2) -- (0,4) node[above] {$y$};
\foreach \x in {-2,-1,1,2,3} \draw (\x,0.1) -- (\x,-0.1) node[below] {\small $\x$};
\foreach \y in {-1,1,2,3} \draw (0.1,\y) -- (-0.1,\y) node[left] {\small $\y$};
% Branche 1 : x < 1, y = x+1
\draw[popcurve, thick, domain=-2.5:1] plot (\x, {\x+1});
\fill[popcream] (1,2) circle (2.5pt); % trou
\draw[popdark, thick] (1,2) circle (2.5pt); % cercle vide
% Branche 2 : x = 1, y = 0 (point isolé)
\fill[popblue] (1,0) circle (2.5pt);
% Branche 3 : x > 1, y = 3
\draw[popcurve, thick, domain=1:3.5] plot (\x, {3});
\fill[popcream] (1,3) circle (2.5pt);
\draw[popdark, thick] (1,3) circle (2.5pt);
% Point plein à (1,3) pour la branche de droite ? Non, x>1 strict.
% On met un point plein à (1,3) si x>=1, mais ici x>1.
% Pour montrer le saut, on laisse le trou à (1,3) et point isolé à (1,0).
\node[popdark] at (1, -0.5) {$1$};
\end{tikzpicture}
\end{center}
\begin{enumerate}
\item En utilisant la définition de la continuité, expliquer pourquoi $k$ n'est pas continue en $x=1$.
\item Calculer $\lim_{x\to 1^-} k(x)$ et $\lim_{x\to 1^+} k(x)$.
\item Quel type de discontinuité observe-t-on en $1$ ?
\end{enumerate}
\end{exercice}

\courssub{Je me prépare au Probatoire}

\begin{exercice}[Théorème des valeurs intermédiaires : Encadrement de racine]
On considère l'équation $x^3 - 2x - 5 = 0$ et la fonction $f(x) = x^3 - 2x - 5$ définie sur $\mathbb{R}$.
\begin{enumerate}
\item Montrer que $f$ est continue et strictement croissante sur l'intervalle $[2 ; 3]$.
\item En déduire, à l'aide du théorème des valeurs intermédiaires, que l'équation admet une unique solution $\alpha$ dans $[2 ; 3]$.
\item Encadrer $\alpha$ à $0,5$ près (donner un intervalle $[a ; b]$ de longueur $0,5$ contenant $\alpha$).
\item Affiner l'encadrement à $0,1$ près en utilisant une méthode de dichotomie (deux itérations suffisent).
\end{enumerate}
\end{exercice}

\begin{exercice}[Modélisation : Bénéfice d'un atelier de menuiserie]
Le bénéfice mensuel (en milliers de FCFA) réalisé par un artisan menuisier en vendant $x$ meubles est modélisé par la fonction $B$ définie sur $[0 ; 80]$ par :
\[ B(x) = -x^2 + 80x - 300 \]
\begin{enumerate}
\item Étudier la continuité de $B$ sur son ensemble de définition.
\item Calculer $\lim_{x\to 0^+} B(x)$ et $\lim_{x\to 80^-} B(x)$. Interpréter ces résultats dans le contexte de l'atelier.
\item Déterminer le nombre de meubles $x$ à vendre pour maximiser le bénéfice (on admet que $B$ est dérivable et on utilise le fait qu'un trinôme $-x^2+80x-300$ admet un maximum au sommet de sa parabole).
\item Calculer le bénéfice maximal.
\item L'artisan a un coût fixe de fonctionnement de $500$ milliers de FCFA par mois. À partir de combien de meubles vendus l'activité devient-elle rentable (bénéfice net positif) ? Résoudre l'inéquation $B(x) > 500$ et arrondir la réponse à l'entier supérieur.
\end{enumerate}
\end{exercice}

\begin{exercice}[Électricité : Charge d'un condensateur]
Dans un circuit RC série alimenté par une tension continue $E$, la tension $u_C(t)$ aux bornes du condensateur au cours de la charge s'exprime par :
\[ u_C(t) = E \left( 1 - e^{-t/\tau} \right) \quad \text{pour } t \ge 0 \]
où $\tau = RC$ est la constante de temps du circuit ($R$ résistance, $C$ capacité).
\begin{enumerate}
\item Étudier la continuité de la fonction $u_C$ sur $[0 ; +\infty[$.
\item Calculer $\lim_{t\to 0^+} u_C(t)$ et $\lim_{t\to +\infty} u_C(t)$. Donner l'interprétation physique de ces deux limites (état initial et régime permanent).
\item On souhaite que la tension aux bornes du condensateur atteigne $95\%$ de la tension d'alimentation $E$. Exprimer cette condition sous forme d'inéquation et déterminer le temps $t_{95}$ nécessaire en fonction de $\tau$ (on donnera une valeur approchée de $\ln 20$).
\item Tracer schématiquement la courbe de $u_C$ en indiquant l'asymptote horizontale, la tangente à l'origine et le point correspondant à $t = \tau$.
\end{enumerate}
\end{exercice}

\newpage

\coursec{Corrigés des exercices}

\coursec{Limites et continuité}

\courssub{Approche intuitive de la notion de limite}

\begin{definition}[Limite finie en un point]
Soit $f$ une fonction définie sur un intervalle (sauf peut-être en $a$) et $\ell\in\mathbb{R}$. On dit que $f$ a pour limite $\ell$ en $a$ (noté $\lim_{x\to a}f(x)=\ell$) si, lorsque $x$ se rapproche de $a$ (sans y être égal), les valeurs $f(x)$ se rapprochent arbitrairement de $\ell$.
\end{definition}

\begin{definition}[Limite infinie en un point]
On dit que $\lim_{x\to a}f(x)=+\infty$ si, quand $x\to a$, les valeurs $f(x)$ deviennent arbitrairement grandes (dépassent tout seuil fixé). De même pour $-\infty$.
\end{definition}

\begin{definition}[Limite en $+\infty$ (ou $-\infty$)]
$\lim_{x\to +\infty}f(x)=\ell$ signifie que lorsque $x$ devient arbitrairement grand, $f(x)$ se rapproche de $\ell$.
$\lim_{x\to +\infty}f(x)=+\infty$ signifie que $f(x)$ devient arbitrairement grand quand $x$ grandit.
\end{definition}

\begin{propriete}[Asymptotes]
\begin{itemize}
\item \textbf{Verticale :} La droite $x=a$ est asymptote verticale à $\mathcal{C}_f$ si $\lim_{x\to a^-}f(x)=\pm\infty$ ou $\lim_{x\to a^+}f(x)=\pm\infty$.
\item \textbf{Horizontale :} La droite $y=\ell$ est asymptote horizontale si $\lim_{x\to +\infty}f(x)=\ell$ ou $\lim_{x\to -\infty}f(x)=\ell$ ($\ell\in\mathbb{R}$).
\item \textbf{Oblique :} La droite $y=ax+b$ ($a\neq 0$) est asymptote oblique en $+\infty$ si $\lim_{x\to +\infty}\big[f(x)-(ax+b)\big]=0$. De même pour $-\infty$.
\end{itemize}
\end{propriete}

\begin{exoresolu}[Exercice 3 — Définitions et vocabulaire]
\begin{enumerate}
\item \textbf{Limite $+\infty$ en $+\infty$ :} $\lim_{x\to +\infty}f(x)=+\infty$ signifie : pour tout réel $A$ (aussi grand soit-il), il existe un réel $B$ tel que pour tout $x>B$, $f(x)>A$. Autrement dit : $f(x)$ finit par dépasser n'importe quel seuil.
\item \textbf{Continuité en un point :} $f$ est continue en $a$ si $f$ est définie en $a$ et $\lim_{x\to a}f(x)=f(a)$ (la limite en $a$ existe, est finie et égale à l'image de $a$).
\item \textbf{Asymptote oblique :} La droite $y=ax+b$ est asymptote oblique à $\mathcal{C}_f$ en $+\infty$ si $\lim_{x\to +\infty}\big[f(x)-(ax+b)\big]=0$, c'est-à-dire que la distance verticale entre la courbe et la droite tend vers $0$.
\item \textbf{Formes indéterminées :} Les principales sont $\ll\dfrac{0}{0}\gg$, $\ll\dfrac{\infty}{\infty}\gg$, $\ll\infty-\infty\gg$ et $\ll 0\times\infty\gg$. Elles ne donnent pas le résultat directement ; il faut transformer l'expression.
\end{enumerate}
\end{exoresolu}

\courssub{Limites de référence et interprétation graphique}

\begin{propriete}[Limites de référence (admises en Première technique)]
\begin{itemize}
\item \textbf{Polynôme :} $\lim_{x\to \pm\infty}P(x)=\lim_{x\to \pm\infty}a_nx^n$ (terme de plus haut degré).
\item \textbf{Fraction rationnelle :} $\lim_{x\to \pm\infty}\frac{P(x)}{Q(x)}=\lim_{x\to \pm\infty}\frac{a_nx^n}{b_mx^m}$ (quotient des termes de plus haut degré).
\item \textbf{Limites remarquables :} $\lim_{x\to 0}\frac{\sin x}{x}=1$ ; $\lim_{x\to 0}\frac{e^x-1}{x}=1$ ; $\lim_{x\to +\infty}e^{-x}=0$ ; $\lim_{x\to -\infty}e^{x}=0$.
\end{itemize}
\end{propriete}

\begin{exoresolu}[Exercice 1 — QCM : Limites à l'infini]
La bonne réponse est la \textbf{b)}.
\begin{enumerate}
\item \textbf{Fausse.} Un polynôme a en $\pm\infty$ la limite de son terme de plus haut degré : $\lim_{x\to +\infty} 2x^3 = +\infty$, donc $\lim_{x\to +\infty}(2x^3-5x+1)=+\infty$.
\item \textbf{Exacte.} Une fraction rationnelle a en $\pm\infty$ la limite du quotient des termes de plus haut degré : $\lim_{x\to +\infty}\dfrac{3x^2}{x^2}=3$.
\item \textbf{Fausse.} $\lim_{x\to +\infty}\dfrac{x+2}{2x-3}=\lim_{x\to +\infty}\dfrac{x}{2x}=\dfrac{1}{2}$.
\item \textbf{Fausse.} $\lim_{x\to -\infty}(x^2-3x)=\lim_{x\to -\infty}x^2=+\infty$.
\end{enumerate}
\end{exoresolu}

\begin{exoresolu}[Exercice 4 — Calculs directs de limites]
\begin{enumerate}
\item Polynôme : on garde le terme de plus haut degré.
\[\lim_{x\to +\infty}(3x^4-2x^2+7)=\lim_{x\to +\infty}3x^4=+\infty.\]
\item Fraction rationnelle : quotient des termes de plus haut degré.
\[\lim_{x\to -\infty}\frac{4x^3-x}{2x^3+5}=\lim_{x\to -\infty}\frac{4x^3}{2x^3}=\frac{4}{2}=2.\]
\item Forme $\ll\frac{0}{0}\gg$ : on factorise. Pour $x\neq 2$ : $\dfrac{x^2-4}{x-2}=\dfrac{(x-2)(x+2)}{x-2}=x+2$. Donc
\[\lim_{x\to 2}\frac{x^2-4}{x-2}=\lim_{x\to 2}(x+2)=4.\]
\item On pose $u=3x$ ; quand $x\to 0$, $u\to 0$ :
\[\frac{\sin(3x)}{x}=3\times\frac{\sin(3x)}{3x}\xrightarrow[x\to 0]{}3\times 1=3.\]
\end{enumerate}
\end{exoresolu}

\begin{exoresolu}[Exercice 5 — Limites de fonctions rationnelles à l'infini et asymptote oblique]
Soit $f(x)=\dfrac{2x^3-5x+1}{x^2+3}$.
\begin{enumerate}
\item Quotient des termes de plus haut degré :
\[\lim_{x\to +\infty}f(x)=\lim_{x\to +\infty}\frac{2x^3}{x^2}=\lim_{x\to +\infty}2x=+\infty ;\qquad \lim_{x\to -\infty}f(x)=\lim_{x\to -\infty}2x=-\infty.\]
\item Cherchons une asymptote oblique. Le degré du numérateur est supérieur de 1 à celui du dénominateur : on effectue la division euclidienne ou on cherche $a,b$ tels que $f(x)-(ax+b)\to 0$.
Pour $x\neq 0$, on teste $y=2x$ :
\[f(x)-2x=\frac{2x^3-5x+1-2x(x^2+3)}{x^2+3}=\frac{-11x+1}{x^2+3}.\]
Or $\lim_{x\to\pm\infty}\dfrac{-11x+1}{x^2+3}=\lim_{x\to\pm\infty}\dfrac{-11x}{x^2}=0$. Donc la droite d'équation $y=2x$ est \cle{asymptote oblique} à $\mathcal{C}_f$ en $+\infty$ et en $-\infty$.
\item $f$ est définie sur $\mathbb{R}$ (car $x^2+3>0$ pour tout $x$). Tableau de comportement :
\begin{center}
\begin{tabular}{|c|ccc|}\hline
$x$ & $-\infty$ & & $+\infty$ \\ \hline
$f(x)$ & $-\infty$ & & $+\infty$ \\ \hline
\end{tabular}
\end{center}
La courbe rejoint son asymptote oblique $y=2x$ aux deux extrémités.
\end{enumerate}
\end{exoresolu}

\courssub{Opérations sur les limites}

\begin{propriete}[Opérations sur les limites (cas où cela a un sens)]
Si $\lim f = \ell_1$ et $\lim g = \ell_2$ (finies ou infinies) :
\begin{itemize}
\item $\lim (f+g) = \ell_1+\ell_2$ (sauf $\ll\infty-\infty\gg$).
\item $\lim (f\times g) = \ell_1\times\ell_2$ (sauf $\ll 0\times\infty\gg$).
\item $\lim (f/g) = \ell_1/\ell_2$ (sauf $\ll 0/0\gg$ et $\ll\infty/\infty\gg$).
\item \textbf{Théorème de composition :} Si $\lim_{x\to a}g(x)=L$ et $f$ continue en $L$, alors $\lim_{x\to a}f(g(x))=f(L)$.
\end{itemize}
\end{propriete}

\begin{aretenir}[Formes indéterminées principales]
$\ll\dfrac{0}{0}\gg$, $\ll\dfrac{\infty}{\infty}\gg$, $\ll\infty-\infty\gg$, $\ll 0\times\infty\gg$, $\ll 1^\infty\gg$, $\ll 0^0\gg$, $\ll \infty^0\gg$.
\emph{Dans chaque cas, il faut transformer l'expression (factorisation, conjugaison, changement de variable, limites remarquables) pour lever l'indétermination.}
\end{aretenir}

\begin{exoresolu}[Exercice 2 — Vrai ou Faux : Opérations sur les limites]
\begin{enumerate}
\item \textbf{Fausse.} C'est une \cle{forme indéterminée} du type $\ll\infty\times 0\gg$. Contre-exemple : $f(x)=x^2$ et $g(x)=\dfrac{1}{x^2}$ en $+\infty$ : $f(x)\to+\infty$, $g(x)\to 0$, mais $f(x)g(x)=1\to 1\neq 0$.
\item \textbf{Vraie.} Limite du quotient des termes de plus haut degré : $\lim_{x\to +\infty}\dfrac{5x^2}{2x^3}=\lim_{x\to +\infty}\dfrac{5}{2x}=0$.
\item \textbf{Vraie \emph{sous l'hypothèse que $f$ est continue sur $\mathbb{R}$}.} Puisque $\lim_{x\to +\infty}f(x)=5$, il existe $A>0$ tel que $f$ soit bornée (par exemple entre $4$ et $6$) sur $[A;+\infty[$. Sur le segment $[0;A]$, une fonction continue est bornée (théorème des bornes) ; de même sur $]-\infty;0]$ si l'on suppose la continuité sur $\mathbb{R}$. \textbf{Attention :} sans continuité, c'est faux (ex : $f(x)=1/x$ pour $x\neq 0$, $f(0)=0$, limite $0$ en $+\infty$ mais non bornée près de $0$).
\item \textbf{Fausse.} La fonction inverse $x\mapsto 1/x$ n'est pas définie en $0$, donc ne peut pas être continue sur $\mathbb{R}$ tout entier. Elle est continue sur $\mathbb{R}^*=]-\infty;0[\cup]0;+\infty[$.
\end{enumerate}
\end{exoresolu}

\courssub{Calculs pratiques : lever les formes indéterminées}

\begin{methode}[Principales techniques pour lever les indéterminations]
\begin{enumerate}
\item \textbf{Factorisation / Simplification :} Pour $\ll\frac{0}{0}\gg$ (racines communes) ou $\ll\frac{\infty}{\infty}\gg$ (factoriser par la plus haute puissance).
\item \textbf{Conjugaison :} Pour $\ll\infty-\infty\gg$ avec racines carrées (multiplier par l'expression conjuguée).
\item \textbf{Changement de variable :} Pour ramener à une limite remarquable (ex : $u=3x$ pour $\sin(3x)/x$).
\item \textbf{Division par la plus haute puissance :} Pour les fractions rationnelles à l'infini.
\item \textbf{Utilisation des limites remarquables :} $\frac{\sin u}{u}\to 1$, $\frac{e^u-1}{u}\to 1$, $\frac{\ln(1+u)}{u}\to 1$ ($u\to 0$).
\end{enumerate}
\end{methode}

\begin{exoresolu}[Exercice 6 — Levée d'indétermination et prolongement par continuité]
Soit $g(x)=\dfrac{x^2-5x+6}{x-2}$.
\begin{enumerate}
\item Le numérateur $x^2-5x+6$ admet $x=2$ pour racine ($4-10+6=0$) ; l'autre racine est $3$ (produit $=6$). Donc $x^2-5x+6=(x-2)(x-3)$ et pour $x\neq 2$ :
\[g(x)=\frac{(x-2)(x-3)}{x-2}=x-3.\]
\item $\lim_{x\to 2}g(x)=\lim_{x\to 2}(x-3)=-1$.
\item Oui : la limite en $2$ existe et est finie. Le prolongement par continuité est la fonction $\tilde{g}$ définie sur $\mathbb{R}$ par :
\[\tilde{g}(x)=x-3\quad\text{pour tout }x\in\mathbb{R},\qquad \tilde{g}(2)=-1.\]
\item La courbe de $g$ est la droite d'équation $y=x-3$ \emph{privée du point} $(2;-1)$ : on représente un petit cercle vide (le « trou ») en ce point.
\begin{popfigure}
\begin{center}
\begin{tikzpicture}[scale=0.8, >=Stealth]
\draw[popline] (-1,0) -- (5,0) node[right] {$x$};
\draw[popline] (0,-4) -- (0,2) node[above] {$y$};
\foreach \x in {1,2,3,4} \draw (\x,0.08) -- (\x,-0.08) node[below] {\small $\x$};
\foreach \y in {-3,-2,-1,1} \draw (0.08,\y) -- (-0.08,\y) node[left] {\small $\y$};
\draw[popcurve, thick, domain=-0.5:4.5] plot (\x, {\x-3});
\fill[popcream] (2,-1) circle (2.5pt);
\draw[popdark, thick] (2,-1) circle (2.5pt);
\node[popdark, right] at (2.1,-1.4) {\small trou en $(2;-1)$};
\end{tikzpicture}
\end{center}
\legende{Courbe de $g$ : la droite $y=x-3$ avec un trou en $x=2$.}
\end{popfigure}
\end{enumerate}
\end{exoresolu}

\begin{exoresolu}[Exercice 7 — Asymptotes verticales, horizontales et tableau de comportement]
Soit $h(x)=\dfrac{4x-1}{x+2}$.
\begin{enumerate}
\item En $x=-2$, le dénominateur s'annule et le numérateur vaut $4\times(-2)-1=-9\neq 0$.
\begin{itemize}
\item Si $x\to -2^-$ ($x<-2$) : $x+2<0$, donc $h(x)=\dfrac{-9}{\text{petit négatif}}\to +\infty$.
\item Si $x\to -2^+$ ($x>-2$) : $x+2>0$, donc $h(x)\to -\infty$.
\end{itemize}
La droite d'équation $x=-2$ est \cle{asymptote verticale}.
\item $\lim_{x\to\pm\infty}h(x)=\lim_{x\to\pm\infty}\dfrac{4x}{x}=4$. La droite d'équation $y=4$ est \cle{asymptote horizontale} en $+\infty$ et en $-\infty$.
\item Tableau de comportement (limites aux bornes du domaine et aux asymptotes) :
\begin{center}
\begin{tabular}{|c|ccccc|}\hline
$x$ & $-\infty$ & & $-2$ & & $+\infty$ \\ \hline
$h(x)$ & $4$ & & $\pm\infty$ & & $4$ \\ \hline
\end{tabular}
\end{center}
(à gauche de $-2$, $h\to+\infty$ ; à droite, $h\to-\infty$).
\item On trace les asymptotes $x=-2$ et $y=4$ en pointillés ; la courbe est une hyperbole dont les deux branches se rapprochent de ces asymptotes.
\begin{popfigure}
\begin{center}
\begin{tikzpicture}[scale=0.55, >=Stealth]
\draw[popline] (-7,0) -- (4,0) node[right] {$x$};
\draw[popline] (0,-5) -- (0,9) node[above] {$y$};
\draw[popmuted, dashed] (-2,-5) -- (-2,9);
\draw[popmuted, dashed] (-7,4) -- (4,4);
\node[popmuted, below right] at (-2,9) {\small $x=-2$};
\node[popmuted, above left] at (4,4) {\small $y=4$};
\draw[popcurve, thick, domain=-7:-2.9, samples=60] plot (\x, {(4*\x-1)/(\x+2)});
\draw[popcurve, thick, domain=-1.15:3.8, samples=60] plot (\x, {(4*\x-1)/(\x+2)});
\end{tikzpicture}
\end{center}
\legende{Courbe de $h$ et ses deux asymptotes.}
\end{popfigure}
\end{enumerate}
\end{exoresolu}

\courssub{Continuité d'une fonction}

\begin{definition}[Continuité en un point / sur un intervalle]
$f$ est \textbf{continue en $a$} si : $f$ est définie en $a$, $\lim_{x\to a}f(x)$ existe (finie) et $\lim_{x\to a}f(x)=f(a)$.
$f$ est \textbf{continue sur un intervalle $I$} si elle est continue en tout point de $I$ (aux bornes, on considère les limites unilatérales).
\end{definition}

\begin{propriete}[Propriétés des fonctions continues]
\begin{itemize}
\item Somme, produit, quotient (si dénominateur non nul), composée de fonctions continues sont continues.
\item Polynômes, fonctions rationnelles (sur leur ensemble de définition), $\sin$, $\cos$, $\exp$, $\ln$ (sur $]0;+\infty[$) sont continues.
\item \textbf{Théorème des valeurs intermédiaires (TVI) :} Si $f$ est continue sur $[a;b]$, alors pour tout $k$ entre $f(a)$ et $f(b)$, il existe $c\in[a;b]$ tel que $f(c)=k$.
\item \textbf{Corollaire (Bijection) :} Si de plus $f$ est strictement monotone sur $[a;b]$, l'équation $f(x)=k$ a une \emph{unique} solution dans $[a;b]$.
\end{itemize}
\end{propriete}

\begin{attention}[Types de discontinuités (Probatoire)]
\begin{itemize}
\item \textbf{Première espèce (saut) :} Limites finies à gauche et à droite, mais différentes (ou égales mais différentes de $f(a)$).
\item \textbf{Deuxième espèce :} Au moins une limite unilatérale est infinie ou n'existe pas.
\item \textbf{Discontinuité évitable :} Limite finie $\ell$ en $a$, mais $f$ non définie en $a$ ou $f(a)\neq\ell$. On peut prolonger par continuité.
\end{itemize}
\end{attention}

\begin{exoresolu}[Exercice 8 — Continuité et définition par morceaux]
Soit $k(x)=\begin{cases} x+1 & \text{si } x<1 \\ 0 & \text{si } x=1 \\ 3 & \text{si } x>1 \end{cases}$.
\begin{enumerate}
\item $k$ est définie en $1$ avec $k(1)=0$ (point plein). Mais la limite à gauche vaut $2$ : $\lim_{x\to 1^-}k(x)=2 \neq k(1)$. La condition $\lim_{x\to 1}k(x)=k(1)$ n'est pas satisfaite, donc $k$ n'est pas continue en $1$.
\item $\lim_{x\to 1^-}k(x)=1+1=2$ (branche $y=x+1$) ; $\lim_{x\to 1^+}k(x)=3$ (branche constante $y=3$).
\item Les limites à gauche et à droite existent, sont finies, mais sont \emph{différentes} ($2\neq 3$) : c'est une \cle{discontinuité de première espèce}, appelée aussi \emph{discontinuité par saut} (saut de $2$ à $3$). De plus la valeur $k(1)=0$ ne coïncide avec aucune des deux limites.
\end{enumerate}
\end{exoresolu}

\begin{exoresolu}[Exercice 9 — TVI : Encadrement de racine par dichotomie]
Soit $f(x)=x^3-2x-5$ sur $[2;3]$.
\textbf{Barème indicatif (sur 5 pts) :} continuité et monotonie (1,5 pt) ; application du TVI avec hypothèses citées (1,5 pt) ; encadrement à $0,5$ (1 pt) ; dichotomie à $0,1$ (1 pt).
\begin{enumerate}
\item $f$ est une fonction \cle{polynôme}, donc continue sur $\mathbb{R}$, en particulier sur $[2;3]$. $f$ est dérivable et $f'(x)=3x^2-2$. Pour $x\in[2;3]$, $3x^2-2\ge 3\times 4-2=10>0$, donc $f$ est \emph{strictement croissante} sur $[2;3]$.
\item $f(2)=8-4-5=-1<0$ et $f(3)=27-6-5=16>0$. $f$ est continue et strictement monotone sur $[2;3]$, et $0\in[f(2);f(3)]$ : d'après le \cle{théorème des valeurs intermédiaires} (corollaire de la bijection), l'équation $f(x)=0$ admet une \emph{unique} solution $\alpha\in[2;3]$.
\item $f(2)=-1<0$ ; $f(2,5)=15,625-5-5=5,625>0$. Donc $\alpha\in[2\,;\,2,5]$, intervalle de longueur $0,5$.
\item Dichotomie :
\begin{itemize}
\item Milieu de $[2;2,5]$ : $m_1=2,25$. $f(2,25)=11,390\,625-4,5-5=1,890\,625>0$, donc $\alpha\in[2;2,25]$.
\item Milieu de $[2;2,25]$ : $m_2=2,125$. $f(2,125)=9,591\ldots-4,25-5=0,341\ldots>0$, donc $\alpha\in[2;2,125]$.
\end{itemize}
Conclusion : $\alpha\in[2\,;\,2,125]$, soit un encadrement à $0,1$ près : $2,0<\alpha<2,1$ (valeur approchée $\alpha\approx 2,09$).
\end{enumerate}
\textbf{Points de vigilance :} citer explicitement les trois hypothèses du TVI (continuité, stricte monotonie, changement de signe) ; vérifier les signes de $f$ aux bornes avant de conclure ; à chaque itération de dichotomie, préciser dans quel sous-intervalle le signe change.
\end{exoresolu}

\courssub{Applications à des situations concrètes}

\begin{exoresolu}[Exercice 10 — Modélisation : Bénéfice d'un atelier de menuiserie]
Un atelier produit et vend $x$ meubles par mois ($0\le x\le 80$). Le bénéfice (en milliers de FCFA) est modélisé par $B(x)=-x^2+80x-300$.
\textbf{Barème indicatif (sur 5 pts) :} continuité (0,5 pt) ; limites et interprétation (1,5 pt) ; maximum (1 pt) ; bénéfice maximal (0,5 pt) ; inéquation de rentabilité (1,5 pt).
\begin{enumerate}
\item $B$ est une fonction polynôme, donc continue sur $\mathbb{R}$, en particulier sur $[0;80]$.
\item $\lim_{x\to 0^+}B(x)=B(0)=-300$ : sans vente, l'atelier perd $300$ milliers de FCFA (coûts fixes). $\lim_{x\to 80^-}B(x)=B(80)=-6400+6400-300=-300$ : à $80$ meubles, les charges de production absorbent tout le gain, la perte est à nouveau de $300$ milliers de FCFA.
\item Le trinôme $-x^2+80x-300$ (coefficient de $x^2$ négatif) atteint son maximum au sommet : $x_0=\dfrac{-80}{2\times(-1)}=40$. Il faut vendre $40$ meubles.
\item $B(40)=-1600+3200-300=1300$ : le bénéfice maximal est de $1\,300$ milliers de FCFA, soit $1\,300\,000$ FCFA.
\item $B(x)>500\iff -x^2+80x-800>0\iff x^2-80x+800<0$. Discriminant : $\Delta=6400-3200=3200$, $\sqrt{\Delta}=40\sqrt{2}\approx 56,57$. Racines : $x_1=\dfrac{80-56,57}{2}\approx 11,72$ et $x_2=\dfrac{80+56,57}{2}\approx 68,28$. Le trinôme est négatif entre ses racines : $11,72<x<68,28$. L'activité devient rentable (bénéfice $>500$ milliers FCFA) à partir de $\boxed{12}$ meubles vendus (entier supérieur à $11,72$).
\end{enumerate}
\textbf{Points de vigilance :} ne pas oublier l'interprétation concrète des limites (en FCFA) ; pour l'inéquation, bien changer le sens en multipliant par $-1$ ; arrondir à l'entier \emph{supérieur} car la rentabilité exige $x>11,72$.
\end{exoresolu}

\begin{exoresolu}[Exercice 11 — Électricité : Charge d'un condensateur]
Un condensateur de capacité $C$ se charge à travers une résistance $R$ sous une tension $E$. La tension à ses bornes est $u_C(t)=E\left(1-e^{-t/\tau}\right)$ avec $\tau=RC$ (constante de temps), pour $t\ge 0$.
\textbf{Barème indicatif (sur 5 pts) :} continuité (0,5 pt) ; limites et interprétation physique (1,5 pt) ; calcul de $t_{95}$ (2 pts) ; schéma (1 pt).
\begin{enumerate}
\item La fonction $t\mapsto e^{-t/\tau}$ est continue sur $\mathbb{R}$ ; $u_C$ est obtenue par produit et somme de fonctions continues, donc $u_C$ est continue sur $[0;+\infty[$.
\item $\lim_{t\to 0^+}u_C(t)=E(1-e^0)=E(1-1)=0$ : à l'instant initial, le condensateur est déchargé (tension nulle). $\lim_{t\to +\infty}e^{-t/\tau}=0$, donc $\lim_{t\to +\infty}u_C(t)=E$ : en régime permanent, le condensateur est chargé et se comporte comme un interrupteur ouvert ; la tension à ses bornes égale la tension d'alimentation. La droite $u=E$ est asymptote horizontale.
\item Condition : $u_C(t)\ge 0,95\,E$ :
\begin{align*}
E\left(1-e^{-t/\tau}\right)\ge 0,95\,E &\iff e^{-t/\tau}\le 0,05\\
&\iff -\frac{t}{\tau}\le \ln(0,05)=-\ln 20\\
&\iff t\ge \tau\ln 20.
\end{align*}
Avec $\ln 20\approx 3,0$ (plus précisément $2,996$) : $t_{95}\approx 3\tau$. Le condensateur atteint $95\,\%$ de sa charge au bout d'environ $3$ constantes de temps.
\item Schéma : courbe croissante partant de l'origine, tangente à l'origine coupant l'asymptote $u=E$ en $t=\tau$, point $(\tau\,;\,0,63E)$, asymptote horizontale $u=E$.
\begin{popfigure}
\begin{center}
\begin{tikzpicture}[scale=1, >=Stealth]
\draw[popline] (0,0) -- (7,0) node[right] {$t$};
\draw[popline] (0,0) -- (0,3.2) node[above] {$u_C$};
\draw[popmuted, dashed] (0,2.5) -- (7,2.5) node[right, popdark] {\small $u=E$};
\draw[popcurve, thick, domain=0:6.8, samples=80] plot (\x, {2.5*(1-exp(-\x/2))});
\draw[poporange, thick, domain=0:2.6] plot (\x, {1.25*\x});
\draw[popmuted, dashed] (2,0) node[below] {\small $\tau$} -- (2,{2.5*(1-exp(-1))}) -- (0,{2.5*(1-exp(-1))}) node[left] {\small $0,63E$};
\fill[popblue] (2,{2.5*(1-exp(-1))}) circle (2pt);
\node[poporange, above right] at (1.6,1.7) {\small tangente en $0$};
\end{tikzpicture}
\end{center}
\legende{Charge du condensateur : $u_C(\tau)\approx 0,63E$ ; asymptote $u=E$.}
\end{popfigure}
\end{enumerate}
\textbf{Points de vigilance :} pour l'inéquation avec l'exponentielle, le sens change lorsqu'on multiplie par $-1$ ; citer $\lim_{u\to +\infty}e^{-u}=0$ ; bien distinguer état transitoire et régime permanent dans l'interprétation.
\end{exoresolu}

\newpage

\coursec{Fiche bilan}

\begin{popfigure}
\begin{tikzpicture}[
  mindmap,
  grow cyclic,
  every node/.style={concept, execute at begin node=\hskip0pt},
  concept color=popblue,
  level 1/.append style={level distance=4.5cm, sibling angle=60},
  level 2/.append style={level distance=3cm, sibling angle=45},
  branch/.style={concept color=#1, edge from parent/.style={draw=#1, thick}}
]
\node[concept, text width=3.2cm, font=\bfseries\large] {Limites \& Continuité}
  child[branch=poporange] { node[concept, text width=2.8cm] {Approche intuitive\\voisinage, tendance} }
  child[branch=popgreen] { node[concept, text width=2.8cm] {Limites de référence\\$x^n$, $1/x$, $\sqrt{x}$} }
  child[branch=poppurple] { node[concept, text width=2.8cm] {Opérations sur limites\\somme, produit, quotient} }
  child[branch=poppink] { node[concept, text width=2.8cm] {Formes indéterminées\\factorisation, conjugaison} }
  child[branch=popgold] { node[concept, text width=2.8cm] {Continuité\\def, TVI, asymptotes} }
  child[branch=popdark] { node[concept, text width=2.8cm] {Applications techniques\\seuil, modélisation} };
\end{tikzpicture}
\legende{Carte mentale : Limites et continuité — Première Technique}
\end{popfigure}

\begin{bilan}
\courssub{Définitions clés}
\begin{itemize}
  \item \cle{Limite finie en $+\infty$} : $\lim_{x\to+\infty}f(x)=\ell \iff \forall \varepsilon>0,\ \exists A>0,\ \forall x>A,\ |f(x)-\ell|<\varepsilon$.
  \item \cle{Limite infinie en $+\infty$} : $\lim_{x\to+\infty}f(x)=+\infty \iff \forall M>0,\ \exists A>0,\ \forall x>A,\ f(x)>M$.
  \item \cle{Limite en $x_0$ (finie ou infinie)} : même idée avec $|x-x_0|<\delta$ (voisinage de $x_0$).
  \item \cle{Continuité en $x_0$} : $f$ définie en $x_0$, $\lim_{x\to x_0}f(x)$ existe finie et $\lim_{x\to x_0}f(x)=f(x_0)$.
  \item \cle{Continuité sur un intervalle} : continue en tout point de l'intervalle (aux bornes : continuité à droite/gauche).
\end{itemize}

\courssub{Théorèmes \& Formules à connaître par cœur}
\begin{center}
\begin{tabularx}{\linewidth}{|>{\bfseries}l|X|}\hline
\textbf{Sujet} & \textbf{Énoncé / Formule} \\\hline
Limites de référence & $\lim_{x\to+\infty}x^n=+\infty\ (n>0)$ ; $\lim_{x\to+\infty}\frac{1}{x}=0$ ; $\lim_{x\to+\infty}\sqrt{x}=+\infty$ \\\hline
Opérations (limites finies) & $\lim(f+g)=\ell_1+\ell_2$ ; $\lim(fg)=\ell_1\ell_2$ ; $\lim(1/g)=1/\ell_2\ (\ell_2\neq0)$ \\\hline
Formes indéterminées (FI) & $0/0$, $\infty/\infty$, $0\times\infty$, $\infty-\infty$, $1^\infty$ \textbf{: calculer autrement} \\\hline
Théorème des Valeurs Intermédiaires (TVI) & $f$ continue sur $[a,b]$, $k\in[f(a),f(b)] \Rightarrow \exists c\in[a,b],\ f(c)=k$ (si $f$ strictement mono : $c$ unique) \\\hline
Asymptotes & Verticale $x=x_0$ si limite infinie en $x_0$ ; Horizontale $y=\ell$ si limite finie $\ell$ en $\pm\infty$ ; Oblique $y=ax+b$ si $\lim_{x\to\pm\infty}\frac{f(x)}{x}=a$ et $\lim_{x\to\pm\infty}(f(x)-ax)=b$ \\\hline
\end{tabularx}
\end{center}

\courssub{Méthodes express}
\begin{enumerate}
  \item \textbf{Limite d'un polynôme / fonction rationnelle en $\pm\infty$} : factoriser par le terme de plus haut degré $\rightarrow$ limite = limite du terme dominant.
  \item \textbf{FI $0/0$ (forme rationnelle)} : factoriser numérateur et dénominateur, simplifier le facteur commun $(x-x_0)$ ou $x^n$.
  \item \textbf{FI $0/0$ (racines)} : multiplier par l'expression \cle{conjuguée} (numérateur ou dénominateur).
  \item \textbf{FI $\infty-\infty$ (racines/polynômes)} : factoriser par le terme dominant ou utiliser conjugaison.
  \item \textbf{Étudier la continuité en $x_0$} : (1) $x_0\in D_f$ ? (2) Calculer $\lim_{x\to x_0}f(x)$ (g/d si nécessaire). (3) Comparer à $f(x_0)$.
  \item \textbf{Résoudre $f(x)=k$ (TVI)} : Vérifier continuité + monotonie sur $[a,b]$, encadrer $k$ par $f(a),f(b)$, conclure existence (et unicité).
\end{enumerate}
\end{bilan}

\begin{aretenir}[Checklist avant le Probatoire]
\begin{itemize}
  \item $\square$ Définir avec ses mots : limite finie, limite infinie, voisinage, continuité en un point.
  \item $\square$ Énoncer et appliquer les 5 limites de référence (puissances, inverse, racine, exponentielle si au programme, logarithme si au programme).
  \item $\square$ Réciter le tableau des opérations sur les limites et \textbf{identifier les 5 formes indéterminées}.
  \item $\square$ Lever une FI $0/0$ par factorisation (remarquables, factorisation par $x-x_0$) ou conjugaison.
  \item $\square$ Lever une FI $\infty-\infty$ ou $0\times\infty$ par factorisation du terme dominant ou changement de variable ($t=1/x$).
  \item $\square$ Calculer une limite en $+\infty$ d'une fonction rationnelle : factoriser par $x^n$ (plus haut degré).
  \item $\square$ Vérifier la continuité en un point : 3 étapes (appartenance $D_f$, existence limite, égalité).
  \item $\square$ Appliquer le TVI : hypothèses (continuité, intervalle), thèse (valeurs intermédiaires), unicité si monotonie stricte.
  \item $\square$ Déterminer les asymptotes (verticale, horizontale, oblique) et les tracer.
  \item $\square$ Modéliser une situation technique (concentration, tension, coût, population) par une fonction et interpréter la limite/continuité.
  \item $\square$ Rédiger un calcul de limite \textbf{étape par étape} avec justifications (théorème utilisé, forme non indéterminée obtenue).
  \item $\square$ Différencier limite à gauche ($\lim_{x\to x_0^-}$) et à droite ($\lim_{x\to x_0^+}$) pour les points de rupture de $D_f$.
\end{itemize}
\end{aretenir}

\findecours

\end{document}
